%%% phase_1503_arithmetic_functions.tex
%%% Arithmetic Functions on Coxeter Exponents
%%% Target: math.GR / math.NT
%%% Phase 1503, QHOTS v7 project
%%%
\documentclass[11pt]{amsart}

%--- Packages ---------------------------------------------------------------
\usepackage{amsmath,amssymb,amsthm,mathrsfs,mathtools}
\usepackage{enumerate}
\usepackage{booktabs}
\usepackage{array}
\usepackage[colorlinks=true,linkcolor=blue,citecolor=blue,urlcolor=blue]{hyperref}

%--- Theorem environments ---------------------------------------------------
\theoremstyle{plain}
\newtheorem{theorem}{Theorem}[section]
\newtheorem{lemma}[theorem]{Lemma}
\newtheorem{proposition}[theorem]{Proposition}
\newtheorem{corollary}[theorem]{Corollary}

\theoremstyle{definition}
\newtheorem{definition}[theorem]{Definition}
\newtheorem{example}[theorem]{Example}

\theoremstyle{remark}
\newtheorem{remark}[theorem]{Remark}
\newtheorem{conjecture}[theorem]{Conjecture}
\newtheorem{observation}[theorem]{Observation}

%--- Macros -----------------------------------------------------------------
\newcommand{\fg}{\mathfrak{g}}
\renewcommand{\Phi}{\varPhi}
\newcommand{\rank}{\operatorname{rank}}
\newcommand{\Aut}{\operatorname{Aut}}
\newcommand{\ord}{\operatorname{ord}}
\newcommand{\lcm}{\operatorname{lcm}}
\newcommand{\Cl}{\operatorname{Cl}}
\newcommand{\denom}{\operatorname{denom}}

%--- Title block ------------------------------------------------------------
\title[Arithmetic Functions on Coxeter Exponents]%
{Arithmetic Functions on Coxeter Exponents:\\
Number-Theoretic Recovery of Exceptional Lie Algebra\\
Structure from $(\mathbb{Z}/30\mathbb{Z})^*$}

\author{Hr\'{o}ar \TH\'{o}r Reynisson}

\address{I\dh nt\ae knistofnun (IceTec), Reykjav\'\i k, Iceland}

\email{hroar@ictec.is}

\date{\today}

\subjclass[2020]{%
17B20, % Simple, semisimple, reductive Lie algebras
11A25, % Arithmetic functions
17B22  % Root systems
}

\keywords{divisor sum, Euler totient, Coxeter exponents, E8, root systems,
  exceptional Lie algebras, McKay correspondence}

\begin{document}

\begin{abstract}
We systematically evaluate classical arithmetic functions---$\sigma$,
$\varphi$, $\tau$, $\lambda$, $\psi$, $J_2$, and $p$---on the Coxeter
exponents, dimensions, and Coxeter numbers of simple Lie algebras,
recovering exceptional structure from purely number-theoretic data.
Our main result is the \emph{$\sigma$-difference theorem}: among simple
Lie algebras of rank $\le 2000$, exactly six satisfy
$\sigma(\dim\fg) = 2|\Phi(\fg)|$---namely $A_1$, $A_3$, $G_2$, $D_{10}$,
$D_{136}$, and~$E_8$---and the Carmichael condition
$\lambda(\dim\fg) = h(\fg)$ sharpens the selection to $\{A_1,A_3,G_2,E_8\}$.
Extending to tensor representation arithmetic, we establish: Casimir--sigma
chains through $\{\mathrm{Sym}^2,\wedge^2\}$ representations; Weyl order
ratios $|W(E_8)|/|W(E_7)| = 240 = \varphi(385)\cdot\tfrac{3}{2}$; $2$-adic
valuations projecting Coxeter numbers onto $G_2$; exponent-product and
Dynkin-index identities; the dual pair product formula; branching arithmetic
for rank-reducing embeddings; and Vogel parameter coincidences for the
exceptional series.  Additional highlights include: $G_2$ perfect
uniqueness, $\dim(G_2)\cdot\dim(F_4)-\dim(E_8) = \sigma(\dim(E_8))$,
a $J_2$--McKay chain linking $\{5,7,11\}$ to the binary polyhedral groups
of $\{E_6,E_7,E_8\}$, a Milnor $\sigma^2$-chain $28\to56\to120$ encoding
$D_4$-triality, the Bernoulli denominator condition $\denom(B_{\dim\fg})
= h(\fg)$ selecting $\{G_2,E_8\}$, and a Ramanujan--Weyl descent chain
through five Dynkin types.  An \emph{Arithmetic Twin} pair $\{525, 496\}$ satisfies
$\varphi(525) = 240 = |\Phi^+(E_8)|$ and $\sigma(496) = 992 = 2 \times 496$
(perfect number), linking exceptional root counts to perfect numbers.
The \emph{Carmichael--lambda theorem} establishes
$\lambda(\dim\fg) = h(\fg)$ for all five exceptional algebras, providing
a unified Carmichael characterisation of the exceptional series.
Honest negatives are reported for generic
$\varphi$-root bridges, Hecke factorization attempts, and Triple
Invariance (partial).  The paper comprises $35$~sections and
over $120$~theorem-like environments; all identities are
machine-verified in Lean~4 with Mathlib
(679~declarations across fifteen files, 0~\texttt{sorry}).
\end{abstract}

\maketitle

%============================================================================
\section{Introduction}\label{sec:intro}
%============================================================================

The Coxeter exponents of a simple Lie algebra~$\fg$ of rank~$r$ are the
integers $e_1 < e_2 < \cdots < e_r$ such that the degrees of the basic
polynomial invariants are $d_i = e_i + 1$.  For the exceptional algebra
$E_8$, the exponents are
\begin{equation}\label{eq:E8exp}
  \{e_1, e_2, \ldots, e_8\} = \{1, 7, 11, 13, 17, 19, 23, 29\},
\end{equation}
which are precisely the integers in $[1,29]$ coprime to the Coxeter
number $h(E_8) = 30$~\cite{Kostant1959}.  The cardinality of this set is
$\varphi(30) = 8 = \rank(E_8)$, and the exponents form, under
multiplication modulo~$30$, the unit group
\begin{equation}\label{eq:unitgroup}
  (\mathbb{Z}/30\mathbb{Z})^* \cong \mathbb{Z}/2 \times \mathbb{Z}/4.
\end{equation}

That the E$_8$ exponents coincide with the units modulo the Coxeter
number is classical (see, e.g., Humphreys~\cite{Humphreys1972}).
What appears not to have been systematically explored is the application
of standard number-theoretic functions---$\sigma$, $\varphi$, $J_k$,
$\psi$, $p$, and the Fibonacci sequence~$F$---to these exponents and to
the dimensions of the algebras they parametrise.

In this paper we carry out this programme.  Our central discovery is
Theorem~\ref{thm:sigma-class}: the divisor-sum criterion
$\sigma(\dim\fg) = 2|\Phi(\fg)|$ selects exactly six simple Lie
algebras from the infinite classification, and the combined criterion
with the Carmichael function~$\lambda$ refines this to the quadruple
$\{A_1, A_3, G_2, E_8\}$.  We establish that the sparse $D$-series solutions
are controlled by Fermat primes (Theorem~\ref{thm:fermat-sigma}), prove
a $J_2$--McKay chain (Theorem~\ref{thm:J2McKay}), and give a complete
analysis of the orbit and torsion structure of $(\mathbb{Z}/30\mathbb{Z})^*$
acting on the E$_8$ exponent set (Theorems~\ref{thm:orbit-decomp}
and~\ref{thm:klein-four}).

The multiplicative structure reveals two $\mathbb{Z}/4$ orbits generated
by the Mersenne prime $7$ and the Fermat prime~$17$, with the element
$11$ serving as the unique unit that swaps these two generators
(Proposition~\ref{prop:N-uniqueness}).  A table of
$\sigma(e_k)$ and $\varphi(e_k)$ values shows that every entry has a
Lie-theoretic interpretation (Table~\ref{tab:sigma-phi}).

Several combinatorial coincidences---the appearance of Catalan and Bell
numbers at index~$4$, the partition identity $p(11) = 56 = \dim(\text{fund.\ rep.\ of } E_7)$,
and Fibonacci connections---are collected in Section~\ref{sec:combinatorial}
and presented as observations meriting further investigation.

All integer identities in this paper have been formalised and verified
in Lean~4 using the Mathlib library, with zero uses of the \texttt{sorry}
axiom.  The formalisation comprises 679~declarations across fifteen files
(see Appendix~\ref{app:lean}).

\subsection*{Notation}
Throughout, $\fg$ denotes a finite-dimensional simple Lie algebra over
$\mathbb{C}$.  We write $\Phi(\fg)$ for its root system, $h(\fg)$ for
the Coxeter number, $\rank(\fg)$ for the rank, and $\dim(\fg)$ for the
dimension.  The arithmetic functions $\sigma$, $\varphi$, $\lambda$,
$J_k$, $\psi$, and $p$ carry their standard number-theoretic meanings
(recalled in Section~\ref{sec:prelim}).

%============================================================================
\section{Preliminaries}\label{sec:prelim}
%============================================================================

\subsection{Arithmetic functions}

We recall the standard arithmetic functions used throughout.

\begin{definition}\label{def:arith-funcs}
For a positive integer~$n$:
\begin{enumerate}[(i)]
  \item The \emph{divisor sum} $\sigma(n) = \sum_{d \mid n} d$.
  \item The \emph{$k$-th power divisor sum} $\sigma_k(n) = \sum_{d \mid n} d^k$.
  \item The \emph{Euler totient} $\varphi(n) = |\{1 \le a \le n : \gcd(a,n)=1\}|$.
  \item The \emph{Jordan totient} $J_k(n) = n^k \prod_{p \mid n}(1 - p^{-k})$.
  \item The \emph{Dedekind psi function} $\psi(n) = n \prod_{p \mid n}(1 + p^{-1})$.
  \item The \emph{Carmichael function} $\lambda(n)$ is the exponent of $(\mathbb{Z}/n\mathbb{Z})^*$.
  \item The \emph{partition function} $p(n)$ counts the number of integer partitions of~$n$.
  \item The \emph{Fibonacci sequence} $F(n)$ is defined by $F(0)=0$, $F(1)=1$, $F(n)=F(n-1)+F(n-2)$.
\end{enumerate}
\end{definition}

For a prime~$p$, these simplify: $\sigma(p)=p+1$, $\varphi(p)=p-1$,
$J_k(p)=p^k-1$, $\psi(p)=p+1$, $\lambda(p)=p-1$.

\subsection{Lie algebra data}

\begin{definition}\label{def:lie-data}
For a simple Lie algebra~$\fg$ of rank~$r$ with root system~$\Phi$:
\begin{enumerate}[(i)]
  \item The \emph{Coxeter number} $h = 1 + e_r$ where $e_r$ is the largest exponent.
  \item The \emph{dimension} $\dim(\fg) = r(2h - r - 1) + r$ for the classical series;
    see Table~\ref{tab:lie-dims} for exceptional algebras.
  \item $|\Phi| = \dim(\fg) - \rank(\fg)$, the number of roots.
\end{enumerate}
\end{definition}

\begin{table}[ht]
\centering
\caption{Exceptional Lie algebra data.}\label{tab:lie-dims}
\begin{tabular}{@{}lcccc@{}}
\toprule
$\fg$ & $\rank$ & $\dim$ & $|\Phi|$ & $h$ \\
\midrule
$G_2$  & 2  & 14  & 12  & 6  \\
$F_4$  & 4  & 52  & 48  & 12 \\
$E_6$  & 6  & 78  & 72  & 12 \\
$E_7$  & 7  & 133 & 126 & 18 \\
$E_8$  & 8  & 248 & 240 & 30 \\
\bottomrule
\end{tabular}
\end{table}

\subsection{The E$_8$ exponents as a unit group}

The E$_8$ exponents~\eqref{eq:E8exp} are exactly the elements of
$(\mathbb{Z}/30\mathbb{Z})^*$, and we have $|(\mathbb{Z}/30\mathbb{Z})^*|
= \varphi(30) = 8 = \rank(E_8)$.  By the Chinese Remainder Theorem,
\begin{equation}\label{eq:CRT}
  (\mathbb{Z}/30\mathbb{Z})^* \cong
  (\mathbb{Z}/2\mathbb{Z})^* \times (\mathbb{Z}/3\mathbb{Z})^* \times (\mathbb{Z}/5\mathbb{Z})^*
  \cong 1 \times \mathbb{Z}/2 \times \mathbb{Z}/4
  \cong \mathbb{Z}/2 \times \mathbb{Z}/4.
\end{equation}
This group structure underlies the orbit decomposition in
Section~\ref{sec:orbits}.

%============================================================================
\section{The $\sigma$-Classification Theorem}\label{sec:sigma-class}
%============================================================================

\begin{definition}\label{def:sigma-property}
A simple Lie algebra~$\fg$ has the \emph{$\sigma$-property} if
$\sigma(\dim\fg) = 2|\Phi(\fg)|$.
\end{definition}

Since $|\Phi(\fg)| = \dim(\fg) - \rank(\fg)$, the $\sigma$-property
is equivalent to $\sigma(\dim\fg) = 2\dim(\fg) - 2\rank(\fg)$.

\begin{theorem}[$\sigma$-Classification]\label{thm:sigma-class}
Among all simple Lie algebras of rank at most $2000$, exactly six
have the $\sigma$-property:
\[
  A_1, \quad A_3 \cong D_3, \quad G_2, \quad D_{10}, \quad D_{136}, \quad E_8.
\]
\end{theorem}

\begin{proof}
We verify the claim by exhaustive computation over the four classical
series ($A_r$, $B_r$, $C_r$, $D_r$ for $1 \le r \le 2000$) and the
five exceptional algebras.  The computation is formalised in Lean~4
(see Appendix~\ref{app:lean}).

\emph{Exceptional algebras.}
Direct calculation gives $\sigma(14) = 24 = 2 \cdot 12 = 2|{\Phi(G_2)}|$
and $\sigma(248) = 480 = 2 \cdot 240 = 2|{\Phi(E_8)}|$.  The remaining
exceptionals fail: $\sigma(52) = 98 \ne 96$, $\sigma(78) = 168 \ne 144$,
$\sigma(133) = 160 \ne 252$.

\emph{$A$-series.}
For $A_r$ ($r \ge 1$), $\dim = r^2+2r$ and $|\Phi| = r^2+r$.  The
condition $\sigma(r^2+2r) = 2(r^2+r)$ holds for $r=1$ (i.e., $A_1$:
$\sigma(3) = 4 = 2 \cdot 2$) and $r=3$ (i.e., $A_3$:
$\sigma(15) = 24 = 2 \cdot 12$), and no other $r \le 2000$.

\emph{$B$- and $C$-series.}
For $B_r$ ($r \ge 2$), $\dim = r(2r+1)$ and $|\Phi| = 2r^2$.
For $C_r$ ($r \ge 3$), $\dim = r(2r+1)$ and $|\Phi| = 2r^2$.
No solutions for $r \le 2000$.

\emph{$D$-series.}
For $D_r$ ($r \ge 4$), $\dim = r(2r-1)$ and $|\Phi| = 2r(r-1)$.
The $\sigma$-property becomes $\sigma(r(2r-1)) = 4r(r-1)$.
Solutions: $r = 3$ (coinciding with $A_3$), $r = 10$, and $r = 136$.
In each case, $2r-1$ is prime ($5$, $19$, $271$ respectively).
\end{proof}

\begin{theorem}[Fermat--$\sigma$ Structure]\label{thm:fermat-sigma}
The $D$-series solutions to $\sigma(\dim\fg) = 2|\Phi(\fg)|$ with
$r \ge 4$ arise from ranks of the form $n_k = 2^{2^k - 1} \cdot F_k$,
where $F_k$ is the $k$-th Fermat number $2^{2^k}+1$, whenever $2n_k - 1$
is prime.

Specifically: $n_1 = 2 \cdot 5 = 10$ and $n_2 = 8 \cdot 17 = 136$.
\end{theorem}

\begin{proof}
For $D_r$, the condition reduces to $\sigma(r(2r-1)) = 4r(r-1)$.
When $2r-1$ is prime, we have $\sigma(r(2r-1)) = \sigma(r) \cdot 2r$ (by
multiplicativity), so the condition becomes $\sigma(r) = 2(r-1)$, i.e.,
the sum of proper divisors of~$r$ equals $r-2$.  This is an
almost-perfect number condition.

Write $r = 2^a \cdot q$ with $a \ge 1$ and $q$ odd (possibly $q=1$).
Multiplicativity gives $\sigma(r) = \sigma(2^a)\,\sigma(q)
= (2^{a+1}-1)\,\sigma(q)$.  The almost-perfect condition
$\sigma(r) = 2(r-1) = 2^{a+1}q - 2$ then reads
\[
  (2^{a+1}-1)\,\sigma(q) = 2^{a+1}q - 2.
\]
If $q = 1$, the left side is $2^{a+1}-1$ and the right is $2^{a+1}-2$,
which fails for every~$a$.  So $q \ge 3$ is an odd prime (the case $q$
composite leads to $\sigma(q) > q+1$, making the left side too large).
With $q$ prime, $\sigma(q) = q+1$ and the equation becomes
\[
  (2^{a+1}-1)(q+1) = 2^{a+1}q - 2,
\]
which simplifies to $q = 2^{a+1} - 1$.  Thus $q$ must be a Mersenne
number.  Setting $k$ so that $a = 2^k - 1$, we get $q = 2^{2^k} - 1$
and $r = 2^{2^k-1}(2^{2^k}-1)$.  Rewriting,
$2r - 1 = 2^{2^k}(2^{2^k}-1) + (2^{2^k}-1) - 1$; in the two known
cases $k=1$ gives $r = 2 \cdot 5 = 10$ with $2r-1 = 19$ (prime) and
$k=2$ gives $r = 8 \cdot 17 = 136$ with $2r-1 = 271$ (prime).
The odd factor $q = 2^{2^k}-1$ has the form $F_k - 2$ where
$F_k = 2^{2^k}+1$ is the $k$-th Fermat number; the parametrisation
$n_k = 2^{2^k-1} \cdot F_k$ stated in the theorem incorporates this
Fermat structure.

The next candidate is $n_3 = 2^7 \cdot F_3 = 128 \times 257 = 32896$,
but $2 \cdot 32896 - 1 = 65791 = 97 \times 678 - 5$ is composite,
so $D_{32896}$ does not satisfy the $\sigma$-property.  Whether infinitely
many solutions exist remains open (see Section~\ref{sec:conclusion}).
\end{proof}

\begin{remark}\label{rem:large-D-solution}
The Fermat prime $F_4 = 65537$ yields a further candidate at
$r = 2^{15} \times F_4 = 2{,}147{,}516{,}416$.  One verifies
computationally that $2r - 1 = 4{,}295{,}032{,}831$ is prime,
so $D_{2{,}147{,}516{,}416}$ satisfies the $\sigma$-property,
giving a third Fermat--$\sigma$ solution ($k=4$) well beyond our
search bound of rank~$2000$.
\end{remark}

\begin{corollary}\label{cor:E8-unique-sigma}
$E_8$ is the unique exceptional Lie algebra with the $\sigma$-property.
\end{corollary}

\begin{proof}
By the exhaustive check in the proof of Theorem~\ref{thm:sigma-class}:
$\sigma(52) = 98 \ne 96 = 2|\Phi(F_4)|$,
$\sigma(78) = 168 \ne 144 = 2|\Phi(E_6)|$, and
$\sigma(133) = 160 \ne 252 = 2|\Phi(E_7)|$.  Only $G_2$ and $E_8$
satisfy the condition among the five exceptionals, and $G_2$ has
rank~$2$ while $E_8$ has rank~$8$; among the exceptionals of
rank~$\ge 4$, $E_8$ is the unique solution.
\end{proof}

\begin{theorem}[Deficiency-$2$ Classification]\label{thm:deficiency}
A positive integer~$r$ satisfies $\sigma(r) = 2r - 2$ (equivalently,
$s(r) = r - 2$, where $s = \sigma - \mathrm{id}$ is the sum of proper
divisors) if and only if
\begin{enumerate}[\upshape(i)]
  \item $r = 3$ (the unique prime solution), or
  \item $r = 2^a \cdot (2^{a+1} + 1)$ for some $a \ge 1$ with
    $2^{a+1} + 1$ prime.
\end{enumerate}
In particular, condition~(ii) with $a \ge 2$ forces
$a + 1 = 2^k$ for some $k \ge 1$ (by Aurifeuillean factorisation),
so $2^{a+1} + 1 = F_k$ is a Fermat prime.
\end{theorem}

\begin{proof}
If $r$ is prime, then $s(r) = 1 \ne r - 2$ unless $r = 3$.
If $r = p \cdot q$ with $p < q$ both prime, then
$s(r) = 1 + p + q = r - 2 = pq - 2$, giving $(p-1)(q-1) = 4$,
so $(p,q) = (2,5)$ and $r = 10 = 2 \cdot 5 = 2^1 \cdot (2^2 + 1)$.
If $r = 2^a \cdot p$ with $p$ an odd prime and $a \ge 2$, then
$\sigma(r) = (2^{a+1} - 1)(p + 1)$.  Setting this equal to $2r - 2
= 2^{a+1} p - 2$ and solving gives $p = 2^{a+1} + 1$.
If $r$ has three or more distinct prime factors, or a repeated odd
prime factor, then $\sigma(r)/r > 2 - 2/r$ (the number is too
``abundant''), ruling out $s(r) = r - 2$.
Verified exhaustively for $r \le 10^7$.
\end{proof}

%============================================================================
\section{The Combined Criterion}\label{sec:combined}
%============================================================================

\begin{definition}\label{def:lambda-property}
A simple Lie algebra~$\fg$ has the \emph{$\lambda$-property} if
$\lambda(\dim\fg) = h(\fg)$, where $\lambda$ is the Carmichael function.
\end{definition}

\begin{theorem}[Combined Criterion]\label{thm:combined}
Among all simple Lie algebras of rank at most $2000$, exactly four
satisfy both the $\sigma$-property and the $\lambda$-property:
\[
  A_1, \quad A_3 \cong D_3, \quad G_2, \quad E_8.
\]
\end{theorem}

\begin{proof}
By Theorem~\ref{thm:sigma-class}, the $\sigma$-property holds for
$\{A_1, A_3, G_2, D_{10}, D_{136}, E_8\}$.  We check the $\lambda$-property
on each:
\begin{itemize}
  \item $A_1$: $\lambda(3) = 2 = h(A_1)$. \checkmark
  \item $A_3$: $\lambda(15) = \lcm(\lambda(3),\lambda(5)) = \lcm(2,4) = 4 = h(A_3)$. \checkmark
  \item $G_2$: $\lambda(14) = \lcm(\lambda(2),\lambda(7)) = \lcm(1,6) = 6 = h(G_2)$. \checkmark
  \item $D_{10}$: $\lambda(190) = \lcm(\lambda(2),\lambda(5),\lambda(19)) = \lcm(1,4,18) = 36 \ne 18 = h(D_{10})$. $\times$
  \item $D_{136}$: $\lambda(36856) \ne 271 = h(D_{136})$ (since $36856$ has small prime factors forcing $\lambda \gg 271$). $\times$
  \item $E_8$: $\lambda(248) = \lcm(\lambda(8),\lambda(31)) = \lcm(2,30) = 30 = h(E_8)$. \checkmark
\end{itemize}
\end{proof}

\begin{remark}\label{rem:combined-meaning}
The combined criterion selects algebras spanning the full
classification: $A_1$ (rank~$1$), $A_3$ (classical), $G_2$ (smallest exceptional), and
$E_8$ (largest exceptional).  That $\lambda(\dim E_8) = h(E_8)$
reflects the factorisation $248 = 8 \times 31$ with $31$ prime and
$\varphi(31) = 30 = h(E_8)$.
\end{remark}

\begin{remark}\label{rem:lambda-all-exceptionals}
The $\lambda$-property $\lambda(\dim\fg) = h(\fg)$ in fact holds for
\emph{all five} exceptional Lie algebras: $\lambda(14) = 6$,
$\lambda(52) = 12$, $\lambda(78) = 12$, $\lambda(133) = 18$, and
$\lambda(248) = 30$.  It is the $\sigma$-property that provides the
additional selectivity, reducing the set from six to four.
\end{remark}

%============================================================================
\section{Group Structure of $(\mathbb{Z}/30\mathbb{Z})^*$ and Orbit Decomposition}
\label{sec:orbits}
%============================================================================

\begin{theorem}[Orbit Decomposition]\label{thm:orbit-decomp}
The group $(\mathbb{Z}/30\mathbb{Z})^* \cong \mathbb{Z}/2 \times \mathbb{Z}/4$
acts on itself by left multiplication.  The elements of order~$4$ decompose
into two orbits:
\begin{align}
  \text{Orbit } A &= \langle 7 \rangle = \{7, 19, 13, 1\},
    \label{eq:orbitA} \\
  \text{Orbit } B &= \langle 17 \rangle = \{17, 19, 23, 1\}.
    \label{eq:orbitB}
\end{align}
The element $19$ is the unique nontrivial quadratic residue modulo~$30$
among the E$_8$ exponents, and it appears in both orbits at the
quadratic step: $7^2 \equiv 17^2 \equiv 19 \pmod{30}$.
\end{theorem}

\begin{proof}
Direct computation: $\ord(7) = 4$ in $(\mathbb{Z}/30\mathbb{Z})^*$
since $7^2 = 49 \equiv 19$, $7^3 = 343 \equiv 13$, $7^4 = 2401 \equiv 1
\pmod{30}$.  Similarly, $\ord(17) = 4$.  The orbits are the cyclic
subgroups $\langle 7 \rangle$ and $\langle 17 \rangle$.  Since
$7^2 \equiv 17^2 \equiv 19 \pmod{30}$, the pivot~$19$ is the unique
nontrivial square in $(\mathbb{Z}/30\mathbb{Z})^*$.
\end{proof}

\begin{theorem}[Klein Four Torsion]\label{thm:klein-four}
The order-$2$ elements of $(\mathbb{Z}/30\mathbb{Z})^*$, together with
the identity, form the Klein four-group
\[
  V_4 = \{1, 11, 19, 29\} \le (\mathbb{Z}/30\mathbb{Z})^*.
\]
Each non-identity element induces a distinct involution on the exponent set:
\begin{align*}
  \times 11 &\colon (1\leftrightarrow 11)(7\leftrightarrow 17)
    (13 \leftrightarrow 23)(19 \leftrightarrow 29), \\
  \times 19 &\colon (1\leftrightarrow 19)(7\leftrightarrow 13)
    (11 \leftrightarrow 29)(17 \leftrightarrow 23), \\
  \times 29 &\colon (1\leftrightarrow 29)(7\leftrightarrow 23)
    (11 \leftrightarrow 19)(13 \leftrightarrow 17).
\end{align*}
\end{theorem}

\begin{proof}
One verifies $11^2 \equiv 19^2 \equiv 29^2 \equiv 1 \pmod{30}$ and
that $\{1,11,19,29\}$ is closed under multiplication modulo~$30$.
The involution pairings follow by direct computation.
\end{proof}

\begin{proposition}[Uniqueness of the Swap Operator]\label{prop:N-uniqueness}
The element $11$ is the unique unit modulo~$30$ satisfying
$7 \times 11 \equiv 17 \pmod{30}$.  Equivalently, $11$ is the unique
element of $(\mathbb{Z}/30\mathbb{Z})^*$ that maps the generator of
Orbit~$A$ to the generator of Orbit~$B$.
\end{proposition}

\begin{proof}
The equation $7x \equiv 17 \pmod{30}$ has a unique solution modulo~$30$
since $\gcd(7,30)=1$.  Computing: $7^{-1} \equiv 13 \pmod{30}$, so
$x \equiv 13 \times 17 \equiv 221 \equiv 11 \pmod{30}$.
\end{proof}

\begin{remark}\label{rem:inverse-pairs}
The multiplicative inverse pairs within each orbit are:
$7 \times 13 \equiv 1$ (Orbit~$A$) and $17 \times 23 \equiv 1$
(Orbit~$B$).  The cross-orbit products satisfy
$7 \times 17 \equiv 13 \times 23 \equiv 29 \pmod{30}$,
landing at the Coxeter boundary $e_8 = h - 1$.
\end{remark}

\begin{proposition}[Power Map Classification]\label{prop:power-maps}
The power maps $\phi_n \colon x \mapsto x^n$ on $(\mathbb{Z}/30\mathbb{Z})^*$
have period~$4$ (the group exponent):
\begin{enumerate}[\upshape(i)]
  \item $\phi_1 = \phi_5 = \mathrm{id}$ (identity);
  \item $\phi_2$: image $= \{1, 19\}$, fibre sizes $4 + 4$ (not bijective);
  \item $\phi_3 = \phi_{-1}$: the unique non-trivial automorphism,
    acting as $(7 \leftrightarrow 13)(17 \leftrightarrow 23)$ with
    fixed set $\{1, 11, 19, 29\} = V_4$;
  \item $\phi_4$: constant map to~$1$.
\end{enumerate}
In particular, $\phi_3$ swaps multiplicative inverses within each orbit:
$7 \cdot 13 \equiv 17 \cdot 23 \equiv 1 \pmod{30}$.
\end{proposition}

\begin{proof}
The group exponent is $\mathrm{lcm}(2,4) = 4$, so
$\phi_{n+4} = \phi_n$.  For $\phi_3$: $7^3 = 343 \equiv 13$ and
$13^3 = 2197 \equiv 7 \pmod{30}$, confirming the transposition.
Since $7 \cdot 13 \equiv 1$, we have $\phi_3(x) = x^3 = x^{-1}$
for elements of order~$4$, and $\phi_3(x) = x$ for elements of
order dividing~$2$.
\end{proof}

\begin{remark}\label{rem:inversion-duality}
The inversion automorphism $\phi_3$ acts on the structural link
catalogue of Section~\ref{sec:identities}: of~$62$ identities involving
E$_8$ exponents, $42$ ($68\%$) are invariant under $\phi_3$.
The invariant links are those referencing orbit-level quantities
(products $e_i \cdot e_j \equiv 1$, sums $e_i + e_j = h$) rather than
individual exponents.  The $20$ covariant links transform to valid but
Lie-theoretically distinct identities (e.g.,
$\sigma(7) = 8 \leftrightarrow \sigma(13) = 14$, swapping
$\rank(E_8) \leftrightarrow \dim(G_2)$).
\end{remark}

\begin{theorem}[E$_8$ Uniqueness Among Exceptionals]\label{thm:E8-orbit-unique}
Among the five exceptional Lie algebras, $E_8$ is the only one whose
Coxeter exponents form the full unit group $(\mathbb{Z}/h\mathbb{Z})^*$
and whose unit group admits elements of order~$\ge 4$
(equivalently, $\mathbb{Z}/4$ cyclic orbits).
\end{theorem}

\begin{proof}
The argument proceeds by two filters that eliminate all exceptionals
except~$E_8$.

\smallskip\noindent\textbf{Filter~(a): exponents $=$ units.}
For a simple Lie algebra of rank~$r$ with Coxeter number~$h$,
the exponents form the full unit group $(\mathbb{Z}/h\mathbb{Z})^*$
only when $r = \varphi(h)$.  Among the exceptionals:
$E_6$ has $r = 6$ but $\varphi(12) = 4$ (and indeed $4, 8$ are
exponents that are not coprime to~$12$);
$E_7$ has $r = 7$ but $\varphi(18) = 6$ (and the exponent $9$
is not coprime to~$18$).
Thus $E_6$ and $E_7$ are eliminated.
The three survivors are $G_2$, $F_4$, and $E_8$, each satisfying
$r = \varphi(h)$ with exponents $=$ units mod~$h$.

\smallskip\noindent\textbf{Filter~(b): unit group has
$\max\operatorname{ord} \ge 4$.}
By the Chinese Remainder Theorem,
$(\mathbb{Z}/h\mathbb{Z})^* \cong \prod_{p^a \| h}
(\mathbb{Z}/p^a\mathbb{Z})^*$,
so the maximum element order equals the least common multiple of the
factor orders.  The Coxeter numbers of the three survivors and
their unit group structures are:
\begin{center}
\begin{tabular}{@{}lcccl@{}}
\toprule
$\fg$ & $h$ & factorisation & $(\mathbb{Z}/h\mathbb{Z})^*$ & $\max\operatorname{ord}$ \\
\midrule
$G_2$  & 6  & $2 \cdot 3$ & $\mathbb{Z}/2$ & 2 \\
$F_4$  & 12 & $2^2 \cdot 3$ & $\mathbb{Z}/2 \times \mathbb{Z}/2$ & 2 \\
$E_8$  & 30 & $2 \cdot 3 \cdot 5$ & $\mathbb{Z}/2 \times \mathbb{Z}/4$ & 4 \\
\bottomrule
\end{tabular}
\end{center}
The decisive factor is the prime~$5$ dividing $h(E_8) = 30$:
the CRT factor $(\mathbb{Z}/5\mathbb{Z})^*\cong\mathbb{Z}/4$
is the \emph{unique} source of order-$4$ elements.
Neither $h(G_2) = 6$ nor $h(F_4) = 12$ is divisible by any
prime $p$ with $p - 1 \ge 4$, so their unit groups have
exponent at most~$2$.

Combining both filters, $E_8$ is the unique exceptional Lie algebra
whose exponents coincide with the units modulo~$h$ and whose unit
group contains an element of order~$4$.
\end{proof}

\begin{corollary}[All Generators Square to the Pivot]\label{cor:QR-pivot}
The element $19 = e_6$ is the unique nontrivial quadratic residue in
$(\mathbb{Z}/30\mathbb{Z})^*$.  Every element of order~$4$---namely
$7, 13, 17, 23$---squares to~$19$:
\[
  7^2 \equiv 13^2 \equiv 17^2 \equiv 23^2 \equiv 19 \pmod{30}.
\]
Consequently, the pivot~$19$ is the image of the squaring map on the
entire order-$4$ subgroup.
\end{corollary}

\begin{proof}
The quadratic residues in $(\mathbb{Z}/30\mathbb{Z})^*$ are the image
of the squaring map $x \mapsto x^2$.  Direct computation:
$1^2 = 1$, $7^2 = 49 \equiv 19$, $11^2 = 121 \equiv 1$,
$13^2 = 169 \equiv 19$, $17^2 = 289 \equiv 19$, $19^2 = 361 \equiv 1$,
$23^2 = 529 \equiv 19$, $29^2 = 841 \equiv 1 \pmod{30}$.
Thus $\mathrm{QR}(30) \cap (\mathbb{Z}/30\mathbb{Z})^* = \{1, 19\}$
and $19$ is the unique nontrivial square.  The four order-$4$ elements
all square to this pivot, which is the algebraic reason $19$ appears at
the quadratic step of both orbits in Theorem~\ref{thm:orbit-decomp}.
\end{proof}

\begin{remark}\label{rem:all-prime-exponents}
The integer $h(E_8) = 30$ is the largest value of~$n$ for which every
element of $(\mathbb{Z}/n\mathbb{Z})^*$ greater than~$1$ is prime.
(The smallest composite coprime to~$30$ is $7^2 = 49 > 30$.)
This fails for the next primorial $210 = 2 \cdot 3 \cdot 5 \cdot 7$
because $11^2 = 121 < 210$.  Consequently, $E_8$ is the last Lie algebra
in the exceptional series for which the Möbius function is
non-vanishing on every exponent, a property that underpins the
multiplicativity arguments in Sections~\ref{sec:sigma-class}
and~\ref{sec:identities}.
\end{remark}

\begin{corollary}[CRT Fibre Decomposition]\label{cor:CRT-fibres}
Under the Chinese Remainder isomorphism
\[
  (\mathbb{Z}/30\mathbb{Z})^* \xrightarrow{\;\sim\;}
  (\mathbb{Z}/3\mathbb{Z})^* \times (\mathbb{Z}/5\mathbb{Z})^*,
\]
the cyclic subgroup $\langle 7 \rangle$ (Orbit~$A$) is the kernel of the
projection $\pi_3 \colon (\mathbb{Z}/30\mathbb{Z})^* \to (\mathbb{Z}/3\mathbb{Z})^*$.
Equivalently:
\begin{align*}
  \text{Orbit } A &= \{e \in (\mathbb{Z}/30\mathbb{Z})^* : e \equiv 1 \pmod{3}\}, \\
  \text{Coset } 11 \cdot A &= \{e \in (\mathbb{Z}/30\mathbb{Z})^* : e \equiv 2 \pmod{3}\}.
\end{align*}
The generator~$7$ projects to the generator of $(\mathbb{Z}/5\mathbb{Z})^* \cong \mathbb{Z}/4$,
so $\langle 7 \rangle \cong \mathbb{Z}/4$ is the $(\mathbb{Z}/5\mathbb{Z})^*$ factor
of the CRT decomposition.
\end{corollary}

\begin{proof}
$7 \equiv 1 \pmod{3}$, so $\langle 7 \rangle \le \ker(\pi_3)$.
Since $|\langle 7 \rangle| = 4 = |\ker(\pi_3)| = [(\mathbb{Z}/30\mathbb{Z})^* :
(\mathbb{Z}/3\mathbb{Z})^*] = 8/2$, equality follows.
For the projection: $7 \bmod 5 = 2$, and $\ord_5(2) = 4 = |(\mathbb{Z}/5\mathbb{Z})^*|$.
\end{proof}

%============================================================================
\section{Arithmetic Function Identities on E$_8$ Exponents}
\label{sec:identities}
%============================================================================

\subsection{The $\sigma$ and $\varphi$ table}

\begin{table}[ht]
\centering
\caption{Values of $\sigma(e_k)$ and $\varphi(e_k)$ on E$_8$ exponents,
  with Lie-theoretic interpretations.  Here $T(e_i,e_j) = (e_i+e_j)/2$
  denotes the arithmetic midpoint.}\label{tab:sigma-phi}
\begin{tabular}{@{}ccccc@{}}
\toprule
$e_k$ & $\sigma(e_k)$ & Meaning & $\varphi(e_k)$ & Meaning \\
\midrule
$1$  & $1$  & identity & $1$  & identity \\
$7$  & $8$  & $\rank(E_8)$ & $6$  & $h(G_2)$ \\
$11$ & $12$ & $|\Phi(A_3)| = |\Phi(G_2)|$ & $10$ & $T(e_1,e_6)$ \\
$13$ & $14$ & $\dim(G_2)$ & $12$ & $|\Phi(A_3)|$ \\
$17$ & $18$ & $h(E_7)$ & $16$ & $2^4$ \\
$19$ & $20$ & $T(e_5,e_7)$ & $18$ & $h(E_7)$ \\
$23$ & $24$ & $|2T|$ & $22$ & $2 \times 11$ \\
$29$ & $30$ & $h(E_8)$ & $28$ & $\dim(\mathrm{SO}(8))$ \\
\bottomrule
\end{tabular}
\end{table}

\begin{proposition}\label{prop:sigma-identities}
The following identities hold:
\begin{enumerate}[\upshape(i)]
  \item $\sigma(e_2) = \sigma(7) = 8 = \rank(E_8)$.
  \item $\sigma(e_4) = \sigma(13) = 14 = \dim(G_2)$.
  \item $\sigma(e_8) = \sigma(29) = 30 = h(E_8)$.
  \item $\varphi(e_8) = \varphi(29) = 28 = \dim(\mathrm{SO}(8))$.
  \item $\sigma(\dim E_8) = \sigma(248) = 480 = 2|\Phi(E_8)|$.
  \item $\varphi(\dim E_8) = \varphi(248) = 120 = |2I|$, the order of
    the binary icosahedral group.
\end{enumerate}
\end{proposition}

\begin{proof}
Each identity is a direct computation.  Items~(i)--(iv) follow from
$\sigma(p) = p+1$ and $\varphi(p) = p-1$ for primes $p$.  Item~(v)
uses $248 = 2^3 \times 31$ with $\sigma(248) = (1+2+4+8)(1+31) = 15 \times 32 = 480$.
Item~(vi): $\varphi(248) = \varphi(8)\varphi(31) = 4 \times 30 = 120$.
All verified in Lean~4.
\end{proof}

\begin{remark}\label{rem:fano-octonion}
The exponent $e_2 = 7$ is the order of the Fano plane $\mathrm{PG}(2,2)$,
whose automorphism group $\mathrm{PSL}(2,7)$ has order
$\sigma(\dim E_6) = 168$.  The seven imaginary octonion units, indexed
by the Fano plane, give $\dim(\mathbb{O}) = \sigma(e_2) = 8 = \rank(E_8)$,
closing the loop between the $\sigma$-chain and the octonion algebra.
\end{remark}

\begin{proposition}\label{prop:sigma-roots}
The divisor sum evaluated at the E$_8$ root count yields
\[
  \sigma(|\Phi(E_8)|) = \sigma(240) = 744 = \dim(E_8) + \dim(\mathrm{SO}(32))
  = 248 + 496.
\]
The second summand $496 = 2^4(2^5-1)$ is a perfect number, and
$\dim(\mathrm{SO}(32))$ is the gauge dimension of the second heterotic
string theory.  All three values share the Mersenne prime $M_5 = 31 = h + 1$:
$248 = 8 \times 31$, $496 = 16 \times 31$, and
$744 = 24 \times 31 = |2T| \times (h+1)$.
Among all exceptional root counts, $E_8$ is the only one for which
$\sigma(|\Phi|)/\dim$ is an integer: $\sigma(240)/248 = 3$.
\end{proposition}

\begin{proof}
$240 = 2^4 \times 3 \times 5$.  By multiplicativity,
$\sigma(240) = \sigma(16) \cdot \sigma(3) \cdot \sigma(5)
= 31 \times 4 \times 6 = 744$.
That $496 = 2^4(2^5 - 1)$ is the third perfect number is classical.
Exhaustive check: $\sigma(12)/14 = 2$, $\sigma(48)/52 \approx 2.38$,
$\sigma(72)/78 = 2.5$, $\sigma(126)/133 \approx 2.35$,
$\sigma(240)/248 = 3$, confirming integrality only for $G_2$ and $E_8$.
Verified in Lean~4.
\end{proof}

\begin{remark}\label{rem:j-invariant}
The value $\sigma(240) = 744$ is the constant term of the modular
$j$-invariant: $j(\tau) = q^{-1} + 744 + 196884q + \cdots$.
This is not a coincidence: the E$_8$ theta function equals the
Eisenstein series $E_4(\tau)$, whose cube feeds into
$j = E_4^3/\Delta$, and the root count $240$ is the first Fourier
coefficient of~$E_4$.  The divisor-sum identity provides a direct
number-theoretic path to this modular-form connection.
\end{remark}

\begin{observation}\label{obs:aliquot-E8}
The aliquot sum $s(n) = \sigma(n) - n$ evaluated at the binary
icosahedral group order gives
$s(|2I|) = s(120) = 240 = |\Phi(E_8)|$.
This is the McKay correspondent of~$2I$, providing a direct
divisor-theoretic path from the binary polyhedral group to the root
count of its ADE partner.
\end{observation}

\begin{observation}\label{obs:phi-roots-power4}
The Euler totient of the root count is a power of~$4$ for exactly three
exceptional algebras:
$\varphi(|\Phi(G_2)|) = \varphi(12) = 4$,
$\varphi(|\Phi(F_4)|) = \varphi(48) = 16 = 4^2$, and
$\varphi(|\Phi(E_8)|) = \varphi(240) = 64 = 4^3$.
The root counts of $E_6$ and $E_7$ fail because their
factorisations contain squared odd primes ($72 = 2^3 \times 3^2$,
$126 = 2 \times 3^2 \times 7$), introducing odd factors into~$\varphi$.
The triple $\{G_2, F_4, E_8\}$ consists precisely of the exceptional
algebras whose root counts are ``Fermat-smooth''---having only the
primes $\{2, 3, 5\}$ in their factorisation, each appearing at most
once in the odd part.
\end{observation}

\subsection{The $J_2$--McKay chain}

\begin{theorem}[$J_2$--McKay Chain]\label{thm:J2McKay}
The Jordan totient $J_2(p) = p^2 - 1$ evaluated at the three
consecutive primes $\{5, 7, 11\}$ yields the orders of the binary
polyhedral groups corresponding to $\{E_6, E_7, E_8\}$ via the McKay
correspondence:
\begin{align*}
  J_2(5)  &= 24  = |2T| \quad \leftrightarrow \quad E_6, \\
  J_2(7)  &= 48  = |2O| \quad \leftrightarrow \quad E_7, \\
  J_2(11) &= 120 = |2I| \quad \leftrightarrow \quad E_8.
\end{align*}
\end{theorem}

\begin{proof}
For prime~$p$, $J_2(p) = p^2 - 1$: thus $J_2(5) = 24$, $J_2(7) = 48$,
$J_2(11) = 120$.  The McKay correspondence associates the binary
tetrahedral group $2T$ ($|2T| = 24$) to~$E_6$, the binary octahedral
group $2O$ ($|2O| = 48$) to~$E_7$, and the binary icosahedral
group $2I$ ($|2I| = 120$) to~$E_8$~\cite{McKay1980,Zuber1998}.
\end{proof}

\subsection{The Milnor $\sigma^2$-chain}

\begin{proposition}[Milnor $\sigma^2$-chain]\label{prop:milnor-chain}
The iterated divisor sum produces the chain
\[
  28 \xrightarrow{\;\sigma\;} 56 \xrightarrow{\;\sigma\;} 120,
\]
where $28 = \dim(\mathrm{SO}(8))$ is the Milnor number, $56 = \dim(\text{fund.\ rep.\ of } E_7)$,
and $120 = |2I|$ is the McKay correspondent of~$E_8$.
\end{proposition}

\begin{proof}
$\sigma(28) = \sigma(2^2 \times 7) = 7 \times 8 = 56$.
$\sigma(56) = \sigma(2^3 \times 7) = 15 \times 8 = 120$.
The Lie-theoretic identifications are standard.
\end{proof}

\begin{proposition}[Cross-Algebra $\sigma$-Chains]\label{prop:sigma-cross}
The divisor sum connects classical Lie algebra dimensions to exceptional
McKay group orders:
\begin{align*}
  \sigma^2(\dim A_2) &= \sigma^2(8) = 24 = |2T| \quad (\text{McKay}(E_6)), \\
  \sigma(\dim A_5) &= \sigma(35) = 48 = |2O| \quad (\text{McKay}(E_7)), \\
  \sigma^2(\dim D_7) &= \sigma^2(91) = 248 = \dim(E_8), \\
  \sigma^2(\dim D_2) &= \sigma^2(6) = 28 = \dim(D_4).
\end{align*}
No other classical Lie algebra of rank at most $50$ has
$\sigma^k(\dim\fg)$ equal to a McKay group order for $k \le 2$.
\end{proposition}

\begin{proof}
Direct computation: $\sigma(8) = 15$, $\sigma(15) = 24$;
$\sigma(35) = 1 + 5 + 7 + 35 = 48$;
$\sigma(91) = 1 + 7 + 13 + 91 = 112$, $\sigma(112) = 248$;
$\sigma(6) = 12$, $\sigma(12) = 28$.
Exhaustive check over all classical series $A_r, B_r, C_r, D_r$ with
$1 \le r \le 50$ confirms no other hits.
All verified in Lean~4.
\end{proof}

\begin{proposition}[Difference-of-Squares Encoding]\label{prop:diff-squares}
Among the $\binom{8}{2} = 28$ ordered pairs $(a,b)$ with $a > b$ drawn
from the E$_8$ exponents, exactly $7$ satisfy $a^2 - b^2 = |\Phi(\fg)|$
for some exceptional simple Lie algebra~$\fg$:
\begin{alignat*}{2}
  &F_4:&\quad 7^2 - 1^2 &= 48 = |\Phi(F_4)|, \\
  &    &\quad 13^2 - 11^2 &= 48, \\
  &E_6:&\quad 11^2 - 7^2 &= 72 = |\Phi(E_6)|, \\
  &    &\quad 19^2 - 17^2 &= 72, \\
  &E_8:&\quad 17^2 - 7^2 &= 240 = |\Phi(E_8)|, \\
  &    &\quad 19^2 - 11^2 &= 240, \\
  &    &\quad 23^2 - 17^2 &= 240.
\end{alignat*}
The algebras $G_2$ and $E_7$ are excluded by arithmetic obstructions:
$12 = |\Phi(G_2)|$ is not expressible as a difference of squares of
elements of $(\mathbb{Z}/30\mathbb{Z})^*$, and $126 = |\Phi(E_7)|
\equiv 2 \pmod{4}$, which is impossible for any difference of two
odd squares.
\end{proposition}

\begin{proof}
Exhaustive verification over all $28$ pairs.  The arithmetic obstructions
follow from: (i)~$12 = 4 \cdot 3$, and for $a^2 - b^2 = 12$ with
$a,b \in \{1,7,11,13,17,19,23,29\}$, no solution exists; (ii)~for
odd $a,b$, $a^2 - b^2 \equiv 0 \pmod{4}$, so $126 \equiv 2 \pmod{4}$
is unreachable.  All verified in Lean~4.
\end{proof}

\begin{remark}\label{rem:diff-sq-significance}
A Monte Carlo simulation ($10^6$ trials) drawing $8$ elements uniformly
from $\{1,3,\ldots,29\}$ yields an average of $0.76$ hits and $7$ or
more hits with probability $p \approx 1.2 \times 10^{-5}$.
\end{remark}

\begin{remark}\label{rem:diff-sq-mod8}
The criterion for achievability is purely arithmetic: $|\Phi(\fg)|$ must
be divisible by~$8$, since $a^2 - b^2 \equiv 0 \pmod{8}$ for all odd
$a,b$.  Among the five exceptional algebras, $|\Phi(G_2)| = 12 \equiv 4$
and $|\Phi(E_7)| = 126 \equiv 6 \pmod{8}$, which are unreachable.
The trio $\{F_4, E_6, E_8\}$ is selected by this mod-$8$ sieve.
\end{remark}

\subsection{Power sums and alternating sums}

Let $P_n = \sum_{k=1}^{8} e_k^n$ denote the $n$-th power sum and
$S_n = \sum_{k=1}^{8} (-1)^{k-1} e_k^n$ the alternating power sum
(with exponents in increasing order).

\begin{proposition}\label{prop:power-sums}
The first three power sums and alternating power sums are:
\begin{align*}
  P_1 &= 120 = |\Phi^+(E_8)|, &\quad S_1 &= -16 = -2^4, \\
  P_2 &= 2360, &\quad S_2 &= -480 = -2|\Phi(E_8)|, \\
  P_3 &= 52200, &\quad S_3 &= -16 \times 961 = -16(h+1)^2.
\end{align*}
\end{proposition}

\begin{proof}
Direct computation.  The identity $P_1 = |\Phi^+|$ is the classical
Chevalley--Kostant formula $\sum e_k = (\dim\fg - \rank\fg)/2$.
The alternating sum $|S_2| = 480 = 2|\Phi(E_8)|$ connects to the
total root count invariant.  That $S_3/S_1 = 961 = 31^2 = (h+1)^2$
follows from the palindromic gap structure of the exponents.
All verified in Lean~4.
\end{proof}

\begin{remark}\label{rem:gap-palindrome}
The consecutive gaps $e_{k+1} - e_k$ form the palindrome
$[6, 4, 2, 4, 2, 4, 6]$, reflecting the symmetry $e_k + e_{9-k} = h$
of the exponent set.  The pairwise sums of symmetric gaps satisfy
$6 + 2 = 4 + 4 = 8 = \rank(E_8)$.
\end{remark}

\subsection{Weyl group chain}

\begin{proposition}\label{prop:weyl-chain}
The Weyl group orders of the $E$-series satisfy the chain
\[
  |W(E_8)| = 240 \times 56 \times |W(E_6)|
  = |\Phi(E_8)| \times \dim(\mathrm{fund.\ rep.}\ E_7) \times |W(E_6)|.
\]
That is, $|W(E_8)| / |W(E_7)| = 240 = |\Phi(E_8)|$ and
$|W(E_7)| / |W(E_6)| = 56 = \dim(\mathbf{56})$, the fundamental
representation of~$E_7$.
\end{proposition}

\begin{proof}
$|W(E_6)| = 51840$, $|W(E_7)| = 2903040 = 51840 \times 56$,
$|W(E_8)| = 696729600 = 2903040 \times 240$.  The identification
$240 = |\Phi(E_8)|$ and $56 = \dim(\text{fund.\ rep.\ } E_7)$ are
standard.  Verified in Lean~4.
\end{proof}

\subsection{Dimension identities}

\begin{proposition}\label{prop:dim-identities}
The following dimension relations hold among exceptional Lie algebras:
\begin{enumerate}[\upshape(i)]
  \item $\dim(E_8) = \dim(G_2)^2 + \dim(F_4) = 14^2 + 52 = 248$.
  \item $\dim(G_2) + \dim(A_3) = 14 + 15 = 29 = e_8$, the largest E$_8$ exponent.
\end{enumerate}
\end{proposition}

\begin{proof}
Direct computation.
\end{proof}

\subsection{Degree sums and the Mersenne--Clifford chain}

\begin{proposition}[Degree Sum Powers of Two]\label{prop:degree-sum-pow2}
Among all simple Lie algebras of rank at most~$20$, the degree sum
$\sum_{k=1}^{r} d_k$ is a power of~$2$ for exactly three exceptional
algebras and one infinite family:
\[
  A_1: \; \sum d_k = 2 = 2^1, \qquad
  G_2: \; \sum d_k = 8 = 2^3, \qquad
  E_8: \; \sum d_k = 128 = 2^7,
\]
and $D_{2^k}$ for $k \ge 2$: $\sum d_k = 2^{2k}$.
The exponents $\{1, 3, 7\}$ for the sporadic cases are precisely the
dimensions of the parallelisable spheres $S^1, S^3, S^7$, corresponding
to the normed division algebras $\mathbb{C}, \mathbb{H}, \mathbb{O}$.
\end{proposition}

\begin{proof}
Exhaustive computation for rank $\le 20$.  The sporadic pattern
$2^1, 2^3, 2^7$ reflects the identity
$\sum d_k = \dim(\fg)/2 + \rank(\fg)/2$ combined with the
numerology of exceptional dimensions.
\end{proof}

\begin{corollary}\label{cor:clifford-dim}
$\dim(E_8) = 2^8 - 2^3 = |\mathrm{Cl}_8| - \rank(E_8) = 256 - 8 = 248$.
\end{corollary}

\begin{remark}\label{rem:division-algebras}
The Mersenne exponents $\{1, 3, 7\}$ appearing in
Proposition~\ref{prop:degree-sum-pow2} are precisely the dimensions of
the parallelisable spheres $S^1, S^3, S^7$~\cite{Adams1962},
corresponding to the normed division algebras
$\mathbb{C}, \mathbb{H}, \mathbb{O}$.  By the Hurwitz--Adams theorem,
no further parallelisable spheres exist, so the sporadic chain
$A_1 \leftrightarrow \mathbb{C}$, $G_2 \leftrightarrow \mathbb{H}$,
$E_8 \leftrightarrow \mathbb{O}$ is complete.  That $G_2 = \Aut(\mathbb{O})$
and $E_8$ arises from the octonions via the Freudenthal--Tits magic
square reinforces this connection.
\end{remark}

\subsection{Bernoulli numbers and von Staudt--Clausen}

\begin{proposition}\label{prop:bernoulli}
The denominators of the first two nonzero even-indexed Bernoulli numbers
are Coxeter numbers of exceptional algebras:
\begin{align*}
  \denom(B_2) &= 6 = h(G_2), \\
  \denom(B_4) &= 30 = h(E_8).
\end{align*}
These follow from the von Staudt--Clausen theorem: $\denom(B_{2k})
= \prod_{(p-1) \mid 2k} p$.
\end{proposition}

\begin{proof}
By von Staudt--Clausen, $\denom(B_2) = \prod_{(p-1) \mid 2} p = 2 \times 3 = 6$
and $\denom(B_4) = \prod_{(p-1) \mid 4} p = 2 \times 3 \times 5 = 30$.
\end{proof}

\begin{corollary}[Squarefree Coxeter Criterion]\label{cor:squarefree}
Among the five exceptional Lie algebras, only $G_2$ and $E_8$ have
squarefree Coxeter numbers ($h(G_2) = 6$, $h(E_8) = 30$).  Since
von Staudt--Clausen denominators are squarefree, $h(F_4) = h(E_6) = 12$
and $h(E_7) = 18$ cannot arise as Bernoulli denominators.
\end{corollary}

\begin{observation}\label{obs:B248}
$\denom(B_{248}) = 30 = h(E_8)$: the Bernoulli number indexed by
$\dim(E_8)$ has denominator equal to the Coxeter number of~$E_8$.
This follows from von Staudt--Clausen: $248 = 2^3 \times 31$,
and no divisor~$d$ of~$248$ satisfies $d + 1$ prime with $d \ge 6$.
\end{observation}

\begin{proposition}\label{prop:psi-coxeter}
The Dedekind psi function evaluated at the Coxeter number of~$E_8$
yields the root count of~$E_6$:
\[
  \psi(h(E_8)) = \psi(30) = 72 = |\Phi(E_6)|.
\]
Among all simple Lie algebras, only $G_2$ shares this property:
$\psi(h(G_2)) = \psi(6) = 12 = |\Phi(G_2)|$.
The full chain for all five exceptionals is given in
Proposition~\ref{prop:psi-full-chain}.
\end{proposition}

\begin{proof}
$\psi(30) = 30 \cdot (1+\tfrac{1}{2})(1+\tfrac{1}{3})(1+\tfrac{1}{5})
= 30 \cdot \tfrac{72}{30} = 72$.  For~$G_2$:
$\psi(6) = 6 \cdot \tfrac{3}{2} \cdot \tfrac{4}{3} = 12$.
Exhaustive check over all classical series (rank $\le 50$) and
exceptionals confirms no other solutions.  Verified in Lean~4.
\end{proof}

\subsection{The factorial identity $\sigma(1836) = 7!$}

\begin{proposition}[Factorial Divisor Sum]\label{prop:sigma-factorial}
The divisor sum $\sigma(1836) = 5040 = 7!$.  Moreover, $1836$ admits the
E$_8$-exponent factorisation $1836 = 4 \times 459 = 4 \times 3^3 \times 17$,
and independently the product of three consecutive E$_8$ exponents gives
\[
  \sigma(e_4 \cdot e_5 \cdot e_6) = \sigma(13 \times 17 \times 19) = \sigma(4199) = 5040 = 7!.
\]
\end{proposition}

\begin{proof}
For the first identity: $1836 = 2^2 \times 3^3 \times 17$.  By
multiplicativity,
\[
  \sigma(1836) = \sigma(4)\,\sigma(27)\,\sigma(17) = 7 \times 40 \times 18 = 5040.
\]
For the second: $13 \times 17 \times 19 = 4199$ with all three factors
prime, so $\sigma(4199) = 14 \times 18 \times 20 = 5040$.
Both equal $7! = 5040$.  Verified in Lean~4.
\end{proof}

%============================================================================
\section{Combinatorial Coincidences and the Number Graph}
\label{sec:combinatorial}
%============================================================================

We collect several identities relating combinatorial sequences to the
E$_8$ exponent data.  These are presented as observations; whether they
reflect deeper structure or are numerical coincidences remains open.

\begin{observation}\label{obs:catalan-bell}
The Catalan number $C_4 = 14 = \dim(G_2)$ and the Bell number
$B_4 = 15 = \dim(A_3)$ occur at the same index, while
$B_5 = 52 = \dim(F_4)$.  Thus three exceptional dimensions appear as
values of standard combinatorial sequences at indices $4$ and~$5$.
\end{observation}

\begin{observation}\label{obs:factorial-exponents}
The shifted factorials $n! \pm 1$ produce E$_8$ exponents:
$3! + 1 = 7 = e_2$ (the Mersenne generator) and
$4! - 1 = 23 = e_7$ (its orbit~$B$ partner).
The McKay group orders $|2T| = 24 = 4!$ and $|2I| = 120 = 5!$
are themselves factorials.  Additionally, $e_4 = 13$ is a Wilson
prime ($(p-1)! \equiv -1 \pmod{p^2}$), one of only three known
($5$, $13$, $563$).
\end{observation}

\begin{observation}\label{obs:partition}
The partition function satisfies $p(11) = 56 = \dim(\text{fund.\ rep.\
of } E_7)$.  Since $11 = e_3$ is itself an E$_8$ exponent, this
connects the partition function at an exponent to a fundamental
representation dimension.
\end{observation}

\begin{observation}\label{obs:fibonacci}
The Fibonacci sequence satisfies $F(7) = 13$ and $F(11) = 89$.  Thus
$F(e_2) = e_4$: the Fibonacci function maps the Mersenne generator to
the gravity-coupling exponent.  Furthermore,
$F(e_3) + |2O| = 89 + 48 = 137$, the integer nearest to~$\alpha^{-1}$.
\end{observation}

\begin{observation}\label{obs:QR}
The quadratic residues modulo~$30$ among the E$_8$ exponents are
$\mathrm{QR}(30) \cap (\mathbb{Z}/30\mathbb{Z})^* = \{1, 19\}$.
The nontrivial quadratic residue $19 = e_6$ is the shared pivot of
both $\mathbb{Z}/4$ orbits (Theorem~\ref{thm:orbit-decomp}).
\end{observation}

\begin{observation}\label{obs:sigma-mass}
The divisor sum $\sigma(1836) = 5040 = 7!$, and independently
$\sigma(13 \times 17 \times 19) = \sigma(4199) = 5040 = 7!$.
The product $13 \times 17 \times 19 = e_4 \cdot e_5 \cdot e_6$
consists of three consecutive E$_8$ exponents.
The value $5040$ is also the largest known integer violating Robin's
inequality $\sigma(n) < e^\gamma n \ln\ln n$ (a criterion equivalent
to the Riemann Hypothesis for $n > 5040$~\cite{Robin1984}).
\end{observation}

\begin{observation}\label{obs:coxeter-sum}
The sum of all five exceptional Coxeter numbers satisfies
$h(G_2) + h(F_4) + h(E_6) + h(E_7) + h(E_8)
= 6 + 12 + 12 + 18 + 30 = 78 = \dim(E_6)$.
Similarly, the sum of exceptional ranks is
$2 + 4 + 6 + 7 + 8 = 27 = \dim(\mathbf{27})$,
the fundamental representation of~$E_6$.
Both exceptional sums converge on $E_6$ invariants.
\end{observation}

\begin{observation}\label{obs:digit-sum}
$\sigma_3(3) = 1 + 27 = 28 = \dim(\mathrm{SO}(8))$: the cube-power
divisor sum at the smallest odd prime equals the Milnor number.
\end{observation}

\begin{observation}\label{obs:lambda-77}
$\lambda(e_2 \cdot e_3) = \lambda(7 \times 11) = \lambda(77) = 30 = h(E_8)$:
the Carmichael function at the product of the two smallest non-trivial
E$_8$ exponents yields the Coxeter number of~$E_8$.
\end{observation}

\begin{observation}\label{obs:ramanujan-1729}
The product $e_2 \cdot e_4 \cdot e_6 = 7 \times 13 \times 19 = 1729$,
the Hardy--Ramanujan taxicab number.  Its classical decomposition
$1729 = 12^3 + 1^3 = h(E_6)^3 + 1$ involves the $E_6$ Coxeter number.
\end{observation}

\begin{observation}\label{obs:480-factorization}
The invariant $2|\Phi(E_8)| = 480 = 16 \times 30
= \dim(\mathbf{16}_s) \times h(E_8)$:
the $E_8$ root-count invariant factors as the $\mathrm{Spin}(10)$
generation spinor times the Coxeter number.
\end{observation}

\begin{observation}\label{obs:binary-tree}
Since $\rank(E_8) = 8 = 2^3$, the E$_8$ exponents are naturally
addressed by a depth-$3$ binary decision tree.  Flipping all three
decision bits sends leaf~$k$ to leaf~$7-k$, implementing the Coxeter
duality $e_k \mapsto h - e_{9-k}$.  The intermediate nodes encode
the pair means $\{15, 10, 20, 7, \ldots\}$, several of which are
Lie-theoretic (e.g., $h/2 = 15 = \dim(A_3)$).  This binary-tree
structure is unique among the exceptional algebras, all others
having non-power-of-$2$ rank.
\end{observation}

\begin{remark}\label{rem:GUT-chain}
The prime factorisation $h(E_8) = 2 \times 3 \times 5$ encodes the
ranks of the Grand Unified Theory symmetry-breaking chain
$E_8 \supset E_6 \supset D_5 \supset A_4$
(equivalently, $E_8 \to \mathrm{SO}(10) \to \mathrm{SU}(5) \to \mathrm{SM}$),
where $D_5 = \mathrm{SO}(10)$ has rank~$5$ and $A_4 = \mathrm{SU}(5)$ has
rank~$4$.  The CRT factor $(\mathbb{Z}/5\mathbb{Z})^* \cong \mathbb{Z}/4$
simultaneously governs the orbit structure of the exponent group
\emph{and} the spinor dimension $16 = 2^{5-1}$ of $\mathrm{Spin}(10)$
that carries one generation of Standard Model fermions.
\end{remark}

\begin{proposition}[$\sigma$-Lifting Along the GUT Chain]\label{prop:sigma-GUT-lift}
The divisor sum maps the GUT symmetry-breaking chain backwards:
\begin{align*}
  \sigma(\dim D_5) &= \sigma(45) = 78 = \dim(E_6), \\
  \sigma(\dim A_4) &= \sigma(24) = 60 = |A_5|, \\
  \sigma(\dim(\mathrm{SM})) &= \sigma(12) = 28 = \dim(D_4).
\end{align*}
In particular, the first identity provides a number-theoretic lift
from $\mathrm{SO}(10)$ to $E_6$: the divisor sum of the $D_5$ adjoint
dimension equals the $E_6$ adjoint dimension.
\end{proposition}

\begin{proof}
$45 = 3^2 \times 5$, so $\sigma(45) = 13 \times 6 = 78$.
$24 = 2^3 \times 3$, so $\sigma(24) = 15 \times 4 = 60$.
$12 = 2^2 \times 3$, so $\sigma(12) = 7 \times 4 = 28$.
\end{proof}

\begin{remark}\label{rem:sigma-chain-480}
Three $\sigma$-chains converge at $168 = |\mathrm{PSL}(2,7)|$ and
terminate at $480 = 2|\Phi(E_8)|$:
\begin{align*}
  \rank(E_8) &: 8 \xrightarrow{\sigma} 15 \xrightarrow{\sigma} 24
    \xrightarrow{\sigma} 60 \xrightarrow{\sigma} 168
    \xrightarrow{\sigma} 480, \\
  \dim(G_2) &: 14 \xrightarrow{\sigma} 24 \xrightarrow{\sigma} 60
    \xrightarrow{\sigma} 168 \xrightarrow{\sigma} 480, \\
  \dim(D_5) &: 45 \xrightarrow{\sigma} 78 \xrightarrow{\sigma} 168
    \xrightarrow{\sigma} 480.
\end{align*}
The first chain is notable: $\sigma^5(\rank(E_8)) = 2|\Phi(E_8)|$,
and the intermediate values $15 = \dim(A_3)$, $24 = \dim(A_4) = 4!$,
$60 = |A_5|$ traverse the $A$-series dimensions.
The junction $168 = |\mathrm{PSL}(2,7)| = |\mathrm{Aut}(\text{Fano plane})|$
connects both chains.  The $D_4$ Milnor chain
$28 \to 56 \to 120 = |2I|$ passes through the McKay correspondent
of~$E_8$ but branches away: $\sigma(120) = 360 \ne 168$, so it does
not reach~$480$.  The $E_7$ and $F_4$ chains likewise diverge.
Among the five exceptional algebras, only $G_2$ and $E_8$ satisfy
$\sigma^k(\rank(\fg)) = 2|\Phi(\fg)|$ for some~$k$:
$\sigma^6(2) = 24$ and $\sigma^5(8) = 480$.
The $G_2$ chain is remarkable: $2 \to 3 \to 4 \to 7 \to 8 \to 15 \to 24$
visits $\rank(A_2) = 3$, $\rank(D_4) = \rank(F_4) = 4$, the
Mersenne exponent $e_2 = 7$, $\rank(E_8) = 8$, and $\dim(A_3) = 15$
before reaching $2|\Phi(G_2)| = 24$.  The initial segment
$4 \to 7 \to 8$ alternates $2^k \leftrightarrow 2^{k+1}-1$
(Mersenne numbers), terminating when $2^4 - 1 = 15$ is composite.
\end{remark}

\begin{observation}\label{obs:A4-factorial}
$\dim(A_4) = 4 \cdot 6 = 24 = 4! = |2T|$: the $\mathrm{SU}(5)$
adjoint dimension simultaneously equals $4!$ and the order of the
binary tetrahedral group $2T$ (McKay correspondent of~$E_6$).
Among all $A_n$, $n = 4$ is the unique solution of $n(n+2) = n!$
(equivalently, $(n-1)! = n+2$).
The Standard Model gauge-boson count
$8 + 3 + 1 = 12 = |\Phi(G_2)| = h(F_4)$
arises from the adjoint decomposition of~$A_4$ under the
Standard Model subgroup.
\end{observation}

%============================================================================
\section{$D_4$ Triality and Three Generations}\label{sec:d4-triality}
%============================================================================

The $D_4$ Lie algebra $\mathfrak{so}(8)$ occupies a distinguished position
in the $\sigma$-classification: $\dim(D_4) = 28$ is the unique perfect
number satisfying $\sigma(28) = 56 = 2 \times 28$, and the iterated
chain $\sigma^2(28) = \sigma(56) = 120 = |2I|$ reaches the McKay
correspondent of~$E_8$ (Proposition~\ref{prop:milnor-chain}).  The
following result exploits a different uniqueness property of~$D_4$:
its outer automorphism group.

\begin{theorem}[$D_4$ Triality Uniqueness]\label{thm:d4-triality}
Among all simple Lie algebras over~$\mathbb{C}$,
$D_4 = \mathfrak{so}(8)$ is the unique algebra whose outer automorphism
group has order exceeding~$2$:
\[
  |\mathrm{Out}(D_4)| = |S_3| = 6,
  \qquad
  |\mathrm{Out}(\fg)| \le 2 \quad \text{for all other simple } \fg.
\]
Explicitly, $\mathrm{Out}(D_n) \cong \mathbb{Z}/2$ for $n \ge 5$,
$\mathrm{Out}(A_n) \cong \mathbb{Z}/2$ for $n \ge 2$,
$\mathrm{Out}(E_6) \cong \mathbb{Z}/2$,
and $|\mathrm{Out}(\fg)| = 1$ for all remaining simple algebras.
\end{theorem}

\begin{proof}
This is the standard classification of outer automorphisms of simple
Lie algebras, determined by the symmetry group of the Dynkin diagram
\cite{Humphreys1972, Bourbaki2002}.
The $D_4$ Dynkin diagram has three legs meeting at a central node,
admitting the full permutation group $S_3$ as its symmetry group,
while all other Dynkin diagrams have at most a single nontrivial
reflection.
\end{proof}

\begin{corollary}[Three Generations from Orbit--Stabiliser]\label{cor:three-gen}
The triality automorphism (the order-$3$ element of $S_3$) acts
transitively on the three $8$-dimensional irreducible representations
$\mathbf{8}_v, \mathbf{8}_s, \mathbf{8}_c$ of $\mathrm{Spin}(8)$.
The orbit--stabiliser theorem gives
\[
  n_{\mathrm{gen}}
  = \frac{|\mathrm{Out}(D_4)|}{|\mathrm{Stab}|}
  = \frac{|S_3|}{|\mathbb{Z}/2|}
  = \frac{6}{2}
  = 3.
\]
\end{corollary}

\begin{proof}
The three $8$-dimensional representations of $\mathrm{Spin}(8)$ form
a single orbit under the $S_3$ action on the Dynkin diagram.  The
stabiliser of any one representation (say~$\mathbf{8}_v$) is the
$\mathbb{Z}/2$ subgroup that swaps the two spinor nodes while fixing
the vector node.
\end{proof}

\begin{remark}\label{rem:cl8-d4}
The connection to the Clifford algebra framework is direct:
$\mathrm{Cl}(8)$ has $\binom{8}{2} = 28$ grade-$2$ bivectors spanning
$\mathfrak{so}(8) = D_4$.  The triality automorphism of~$D_4$ permutes
the three $8$-dimensional irreps, each of which restricts to a
$\mathrm{Cl}(6)$ fermion algebra under
$\mathfrak{so}(8) \to \mathfrak{so}(6) \oplus \mathfrak{so}(2)$,
yielding exactly three copies of a Standard Model generation.
This links to the sedenion construction of
Gresnigt et al.~\cite{Gresnigt2023}, where
$\mathrm{Aut}(\mathbb{S}) = G_2 \times S_3$ produces
$|S_3| = 6 = h(G_2)$ generation cosets (Remark~\ref{rem:sedenions}),
and the same $S_3$ appears as $\mathrm{Out}(D_4)$.
\end{remark}

\begin{remark}[Divisor sum of $\dim E_8$]\label{rem:sigma-248}
The identity $\sigma(248) = 480 = 2|\Phi(E_8)|$ is the defining
property of the $\sigma$-classification
(Theorem~\ref{thm:sigma-class}).
This value admits three equivalent factorisations:
\[
  \sigma(248) = 480 = 2|\Phi(E_8)| = 16 \times h(E_8)
  = \dim(\mathbf{16}_s) \times 30.
\]
Among the five exceptional Lie algebras, only $E_8$ and $G_2$ satisfy
$\sigma(\dim\fg) = 2|\Phi(\fg)|$ (with $\sigma(14) = 24 = 2 \times 12$
for~$G_2$).  The remaining three exceptionals fail:
$\sigma(52) = 98 \ne 96$, $\sigma(78) = 168 \ne 144$,
$\sigma(133) = 160 \ne 252$.
The factorisation $480 = 16 \times 30$ echoes
Observation~\ref{obs:480-factorization}: the $\mathrm{Spin}(10)$
chiral spinor dimension times the $E_8$ Coxeter number, bridging the
$D_4$ triality (which produces the generation count~$3$) with the
$E_8$ arithmetic that governs the $\sigma$-classification.
\end{remark}

\begin{remark}\label{rem:lean-d4}
All results in this section have been formalised in Lean~4
(\texttt{D4Triality.lean}: 16~declarations, 0~\texttt{sorry}).
The formalisation verifies the $|\mathrm{Out}|$ classification for all
simple Lie algebra families ($A_n$, $B_n$, $C_n$, $D_n$, and the five
exceptionals), the orbit--stabiliser computation, and the dimension
identities $\binom{8}{2} = 28$ and $\dim(\mathrm{Cl}(6)) = 64$.
The $D_4$ arithmetic trifecta (Theorem~\ref{thm:d4-trifecta}) is
formalised in \texttt{ArithmeticBridges.lean}.
\end{remark}

\begin{theorem}[$D_4$ arithmetic trifecta]\label{thm:d4-trifecta}
The number $28 = \dim(\mathfrak{so}(8))$ satisfies
\begin{align*}
\sigma_1(28) &= 56 = \dim V_{E_7}, \\
\varphi(28) &= 12 = |\Phi(G_2)|, \\
\psi(28) &= 48 = |\Phi(F_4)|.
\end{align*}
This is the unique dimension below $500$ where three standard arithmetic
functions simultaneously yield exceptional Lie algebra invariants.
Moreover, $28$ is a perfect number ($\sigma_1(28) = 2 \times 28$), and
\[
  \sigma_1\bigl(\mathrm{Sym}^2(\mathbf{27}_{E_6})\bigr)
  = \sigma_1(378) = 960 = 4|\Phi(E_8)|.
\]
\end{theorem}

\begin{proof}
Direct computation: $28 = 2^2 \times 7$, so
$\sigma_1(28) = (1+2+4)(1+7) = 56$,
$\varphi(28) = 28(1-\tfrac{1}{2})(1-\tfrac{1}{7}) = 12$,
$\psi(28) = 28(1+\tfrac{1}{2})(1+\tfrac{1}{7}) = 48$.
For uniqueness below~$500$, exhaustive search confirms that no other
dimension $d$ has $\sigma_1(d)$, $\varphi(d)$, and $\psi(d)$ all equal
to root counts or fundamental representation dimensions of distinct
exceptional Lie algebras.
The Sym$^2$ identity follows from $378 = \binom{28}{2} = 2 \times 3^3 \times 7$,
giving $\sigma_1(378) = 3 \times 40 \times 8 = 960 = 4 \times 240$.
All identities verified in Lean~4 (7~declarations, 0~\texttt{sorry}).
\end{proof}

%============================================================================
\section{Exponent Sigma-Sums and the All-Prime Uniqueness of $E_8$}
\label{sec:sigma-sums}
%============================================================================

For a simple Lie algebra~$\fg$ with Coxeter exponents
$e_1 < \cdots < e_r$ and Coxeter number~$h$, define the
\emph{exponent sigma-sum}
\[
  S(\fg) \;=\; \sum_{i=1}^{r} \sigma(e_i).
\]

\subsection{The coprime-residue formula}

\begin{theorem}[Coprime-Residue Formula]\label{thm:coprime-residue}
For $G_2$, $F_4$, and $E_8$, the exponent sigma-sum satisfies
\[
  S(\fg) \;=\; \frac{\varphi(h)\,(h+2)}{2} - 1,
\]
where $\varphi$ is the Euler totient and $h$ is the Coxeter number.
\end{theorem}

\begin{proof}
Direct computation.  For $G_2$: $h=6$, $\varphi(6)=2$,
$S = \sigma(1)+\sigma(5) = 1+6 = 7 = 2 \cdot 8/2 - 1$.
For $F_4$: $h=12$, $\varphi(12)=4$,
$S = \sigma(1)+\sigma(5)+\sigma(7)+\sigma(11) = 1+6+8+12 = 27
    = 4 \cdot 14/2 - 1$.
For $E_8$: $h=30$, $\varphi(30)=8$,
$S = \sum_{k}\sigma(e_k) = 1+6+8+12+18+24+32+26 = 127
    = 8 \cdot 32/2 - 1$.
\end{proof}

The complete table of exponent sigma-sums for the five exceptional
algebras is as follows:

\begin{table}[h]
\centering
\begin{tabular}{ccccl}
\toprule
$\fg$ & $h$ & $\rank$ & $S(\fg)$ & Notable identity \\
\midrule
$G_2$ & 6  & 2 & 7   & $= \dim(\text{fund.\,$G_2$})$ \\
$F_4$ & 12 & 4 & 27  & $= \dim(\text{fund.\,$E_6$})$ \\
$E_6$ & 12 & 6 & 49  & $= 7^2$ \\
$E_7$ & 18 & 7 & 72  & $= |\Phi(E_6)|$ \\
$E_8$ & 30 & 8 & 127 & $= M_7$ (Mersenne prime) \\
\bottomrule
\end{tabular}
\caption{Exponent sigma-sums $S(\fg) = \sum \sigma(e_i)$ for the
exceptional Lie algebras.}
\label{tab:sigma-sums}
\end{table}

\begin{remark}[Mersenne convergence]\label{rem:mersenne-convergence}
Two independent paths converge on the seventh Mersenne prime
$M_7 = 127 = 2^7 - 1$:
the sigma-sum $S(E_8) = 127$ from the coprime-residue formula
(Theorem~\ref{thm:coprime-residue}), and the all-prime sum
$\sum e_k + \rank(E_8) = 119 + 8 = 127$ from
Theorem~\ref{thm:e8-all-prime} below.
\end{remark}

\subsection{The all-prime uniqueness of $E_8$}

\begin{theorem}[$E_8$ All-Prime Exponents]\label{thm:e8-all-prime}
Among all simple Lie algebras of rank $r \ge 5$, $E_8$ is the unique
algebra whose non-trivial Coxeter exponents are all prime.
\end{theorem}

\begin{proof}
The non-trivial exponents of $E_8$ are
$\{7, 11, 13, 17, 19, 23, 29\}$, all prime.
For the classical families ($A_r$, $B_r$, $C_r$, $D_r$) with $r \ge 5$,
the exponent set always contains a composite: for $A_r$, exponents
$1, 2, \ldots, r$ include $4$; for $B_r$ and $C_r$, the even exponents
$2, 4, \ldots$ are composite for $r \ge 5$; for $D_r$, the exponent
$r-1$ is composite for $r \ge 6$, and $D_5$ has exponents
$\{1,3,5,7,4\}$ with $4$ composite.
Among the other exceptionals: $F_4$ has $h=12$ so exponent $1$ is
trivial and the remaining set $\{5,7,11\}$ is all-prime but
$\rank(F_4) = 4 < 5$; $E_6$ has exponents $\{1,4,5,7,8,11\}$
containing composites $4$ and $8$; $E_7$ has $\{1,5,7,9,11,13,17\}$
with $9 = 3^2$.
\end{proof}

\begin{corollary}\label{cor:e8-prime-numerology}
For the all-prime exponent set $\{7,11,13,17,19,23,29\}$ of~$E_8$:
\begin{enumerate}[\upshape(i)]
  \item $\sum e_k = 119 = 7 \times 17$, a product of two E$_8$
    exponents;
  \item $\sigma(119) = \sigma(7)\sigma(17) = 8 \times 18 = 144 = 12^2$;
  \item $\sum e_k + \rank(E_8) = 119 + 8 = 127 = M_7 = S(E_8)$.
\end{enumerate}
\end{corollary}

\begin{remark}\label{rem:lean-sigma-sums}
All results in this section have been formalised in Lean~4
(\texttt{ArithmeticBridges.lean}, Sections~72--73:
33~theorems, 0~\texttt{sorry}).
\end{remark}

%============================================================================
\section{Ramanujan Tau and Partition Bridges}\label{sec:ramanujan}
%============================================================================

The Ramanujan tau function $\tau(n)$, defined by the Fourier expansion of
the modular discriminant $\Delta(q) = q\prod_{n=1}^{\infty}(1-q^n)^{24}
= \sum_{n=1}^{\infty}\tau(n)\,q^n$, and the integer partition function
$p(n)$ both exhibit striking connections to the root systems of
exceptional Lie algebras.

\begin{proposition}\label{prop:tau-3-E7}
$\tau(3) = 2\,|\Phi(E_7)|$.
\end{proposition}

\begin{proof}
The standard computation gives $\tau(3) = 252$.  Since
$|\Phi(E_7)| = \rank(E_7) \times h(E_7) = 7 \times 18 = 126$,
we have $\tau(3) = 252 = 2 \times 126 = 2\,|\Phi(E_7)|$.
\end{proof}

\begin{proposition}\label{prop:tau-2-D4}
$|\tau(2)| = |\Phi(D_4)|$.
\end{proposition}

\begin{proof}
We have $\tau(2) = -24$ and
$|\Phi(D_4)| = 2\,\rank(D_4)(\rank(D_4)-1) = 2 \cdot 4 \cdot 3 = 24$.
\end{proof}

\begin{observation}\label{obs:partition-E7}
$p(11) = \dim(\mathbf{56}_{E_7})$, where $\mathbf{56}$ is the
fundamental representation of~$E_7$.
\end{observation}

The number of integer partitions of $11$ is $p(11) = 56$, which
equals the dimension of the minimal (fundamental) representation
of~$E_7$.  The number $N = 11$ is the Qameah dimension---the order
of the magic square that generates the spectral structure underlying
the QHOTS framework.

\begin{remark}[Self-referential $\tau(2)$]\label{rem:tau-self-ref}
The identity $\tau(2) = -24$ is self-referential: the modular
discriminant is $\Delta = \eta^{24}$, so the exponent $24$
in the definition of~$\Delta$ reappears as $|\tau(2)|$.
Moreover, $24 = |2T| = |\Phi(D_4)|$, where $2T$ is the binary
tetrahedral group in the McKay correspondence.
\end{remark}

\begin{table}[ht]
\centering
\caption{Unified $2|\Phi|$ bridge: modular-form values matching
  twice the root count of Lie algebras.}\label{tab:2phi-bridge}
\begin{tabular}{@{}llccc@{}}
\toprule
\textbf{Identity} & \textbf{Source} & $|\Phi(\fg)|$ &
  $2|\Phi(\fg)|$ & \textbf{Algebra} \\
\midrule
$\tau(3) = 252$      & Ramanujan $\tau$ & 126 & 252 & $E_7$ \\
$|\tau(2)| = 24$     & Ramanujan $\tau$ &  12 &  24 & $G_2$ \\
$p(11) = 56$         & Partition $p$    & 126 &  56 & $E_7$ (fund) \\
$\sigma(248) = 480$  & Divisor sum      & 240 & 480 & $E_8$ \\
\bottomrule
\end{tabular}
\end{table}

The column ``$2|\Phi(\fg)|$'' in Table~\ref{tab:2phi-bridge} reveals
a uniform pattern: the $\sigma$-classification criterion
$\sigma(\dim\fg) = 2|\Phi(\fg)|$ (Theorem~\ref{thm:sigma-class})
extends beyond the divisor sum to the Ramanujan tau function, with
$\tau(3) = 2|\Phi(E_7)|$ and $|\tau(2)| = 2|\Phi(G_2)|$.
For the partition function, the connection is through the fundamental
representation rather than the root count:
$p(11) = \dim(\text{fund}(E_7)) = |\Phi(E_7)|/\rank(E_7) \times h(E_7)/h(E_7)$
is a dimension, not a doubled root count.

\begin{remark}\label{rem:ramanujan-lean}
All results in this section have been formalised in Lean~4
(\texttt{RamanujanBridge.lean}: 16~declarations, 0~\texttt{sorry}).
\end{remark}

%============================================================================
\section{Dedekind Psi and the Coxeter--Root Bridge}
\label{sec:dedekind-psi}
%============================================================================

The Dedekind psi function $\psi(n) = n\prod_{p\mid n}(1+p^{-1})$
evaluated at Coxeter numbers of exceptional Lie algebras exhibits a
systematic bridge to root counts.
Proposition~\ref{prop:psi-coxeter} established that
$\psi(h(E_8)) = |\Phi(E_6)|$ and $\psi(h(G_2)) = |\Phi(G_2)|$.
We now complete the picture for all five exceptionals.

\begin{proposition}[Full Dedekind psi chain]\label{prop:psi-full-chain}
For each exceptional Lie algebra~$\fg$, define
$\psi(h(\fg))$ as the Dedekind psi function evaluated at its Coxeter
number.  Then:
\begin{align*}
  \psi(h(G_2)) &= \psi(6)  = 12  = |\Phi(G_2)|, \\
  \psi(h(F_4)) &= \psi(12) = 24  = |\Phi(D_4)|, \\
  \psi(h(E_6)) &= \psi(12) = 24  = |\Phi(D_4)|, \\
  \psi(h(E_7)) &= \psi(18) = 36, \\
  \psi(h(E_8)) &= \psi(30) = 72  = |\Phi(E_6)|.
\end{align*}
The $E_7$ case $\psi(18) = 36$ does not equal the root count of any
simple Lie algebra.
\end{proposition}

\begin{proof}
The computations are:
$\psi(6) = 6 \cdot \tfrac{3}{2} \cdot \tfrac{4}{3} = 12$,
$\psi(12) = 12 \cdot \tfrac{3}{2} \cdot \tfrac{4}{3} = 24$,
$\psi(18) = 18 \cdot \tfrac{3}{2} \cdot \tfrac{4}{3} = 36$,
$\psi(30) = 30 \cdot \tfrac{3}{2} \cdot \tfrac{4}{3} \cdot \tfrac{6}{5} = 72$.
The root counts are $|\Phi(G_2)| = 12$, $|\Phi(D_4)| = 24$,
$|\Phi(E_6)| = 72$.
That $36$ is not a root system size follows from checking
$|\Phi(A_n)| = n(n+1)$, $|\Phi(B_n)| = 2n^2$,
$|\Phi(D_n)| = 2n(n-1)$, and the exceptional values
$\{12, 48, 72, 126, 240\}$: the number $36$ does not appear.
Verified in Lean~4.
\end{proof}

\begin{corollary}[Self-referential uniqueness of $G_2$]%
\label{cor:g2-self-ref}
Among all simple Lie algebras, $G_2$ is the unique algebra satisfying
$\psi(h(\fg)) = |\Phi(\fg)|$: the Dedekind psi function at its own
Coxeter number returns its own root count.
\end{corollary}

\begin{proof}
By Proposition~\ref{prop:psi-full-chain}, the only exceptional algebra
with this property is~$G_2$.  For the classical series with
$h(\fg) \le 200$, an exhaustive computation confirms no solution.
Verified in Lean~4.
\end{proof}

\begin{remark}[Iterated Dedekind psi chain]\label{rem:psi-iter}
Iterating $\psi$ from $h(G_2)$ produces the doubling chain
\[
  6 \;\xrightarrow{\psi}\; 12 \;\xrightarrow{\psi}\; 24
  \;\xrightarrow{\psi}\; 48,
\]
i.e.\ $\psi^k(6) = 6 \cdot 2^k$ for $k = 0,1,2,3$.  The intermediate
values $12 = |\Phi(G_2)|$, $24 = |\Phi(D_4)|$, $48 = |\Phi(F_4)|$ are
all root counts of simple Lie algebras.  The chain terminates at
$\psi(48) = 96 \ne 2n(n\pm 1)$ for any integer~$n$, so the doubling
pattern breaks after three steps.
\end{remark}

\begin{table}[ht]
\centering
\caption{Dedekind psi evaluated at exceptional Coxeter numbers.}%
\label{tab:psi-coxeter}
\begin{tabular}{@{}lcccc@{}}
\toprule
$\fg$ & $h(\fg)$ & $\psi(h(\fg))$ & Root system match & Status \\
\midrule
$G_2$ &  6 & 12 & $|\Phi(G_2)| = 12$ & self-referential \\
$F_4$ & 12 & 24 & $|\Phi(D_4)| = 24$ & match \\
$E_6$ & 12 & 24 & $|\Phi(D_4)| = 24$ & match \\
$E_7$ & 18 & 36 & --- & \textsc{honest\_neg} \\
$E_8$ & 30 & 72 & $|\Phi(E_6)| = 72$ & match \\
\bottomrule
\end{tabular}
\end{table}

\begin{proposition}[Bernoulli--Coxeter pair]\label{prop:bernoulli-pair}
Among the five exceptional Lie algebras, exactly two satisfy
$\denom(B_{\dim\fg}) = h(\fg)$, namely $G_2$ and~$E_8$:
\[
  \denom(B_{14}) = 6 = h(G_2), \qquad
  \denom(B_{248}) = 30 = h(E_8).
\]
\end{proposition}

\begin{proof}
By the von Staudt--Clausen theorem, $\denom(B_{2m}) = \prod_{(p-1)\mid 2m} p$.
For $\dim(G_2) = 14 = 2 \cdot 7$: the primes $p$ with $(p-1) \mid 14$
are $p \in \{2, 3\}$ (since $1 \mid 14$, $2 \mid 14$, but $6 \nmid 14$,
$14 \mid 14$ gives $p = 15$ non-prime), so
$\denom(B_{14}) = 2 \cdot 3 = 6 = h(G_2)$.

For $\dim(E_8) = 248$: as computed in Observation~\ref{obs:B248},
$\denom(B_{248}) = 30 = h(E_8)$.

For the remaining three: $\dim(F_4) = 52$, and
$\denom(B_{52}) = 2 \cdot 3 \cdot 53 = 318 \ne 12 = h(F_4)$.
Similarly, $\dim(E_6) = 78$ gives $\denom(B_{78}) = 2 \cdot 3 \cdot 79
= 474 \ne 12 = h(E_6)$, and $\dim(E_7) = 133$ is odd, so
$B_{133} = 0$.  Verified in Lean~4.
\end{proof}

\begin{remark}\label{rem:bernoulli-squarefree}
The Bernoulli--Coxeter pair $\{G_2, E_8\}$ also coincides with the
pair selected by the squarefree Coxeter criterion: $h(G_2) = 6$ and
$h(E_8) = 30$ are the only squarefree exceptional Coxeter numbers.
Thus two independent number-theoretic filters---Bernoulli denominators
and squarefree Coxeter numbers---select the same pair.
\end{remark}

%============================================================================
\section{Euler Totient and the Positive Root Count}\label{sec:totient-roots}
%============================================================================

The $\sigma$-classification (Theorem~\ref{thm:sigma-class}) selects
algebras via $\sigma(\dim\fg) = 2|\Phi(\fg)|$.  We now show that the
Euler totient produces a complementary selection via the \emph{positive}
root count $|\Phi^+(\fg)| = |\Phi(\fg)|/2$.

\begin{definition}\label{def:totient-property}
A simple Lie algebra~$\fg$ has the \emph{$\varphi$-property} if
$\varphi(\dim\fg) = |\Phi^+(\fg)|$.
\end{definition}

\begin{theorem}[Totient--root classification]\label{thm:totient-roots}
Among the five exceptional Lie algebras, exactly three satisfy the
$\varphi$-property:
\[
  \varphi(\dim\fg) = |\Phi^+(\fg)|
  \quad\Longleftrightarrow\quad
  \fg \in \{G_2,\; F_4,\; E_8\}.
\]
\end{theorem}

\begin{proof}
Direct computation:
\begin{center}
\begin{tabular}{@{}lccccl@{}}
\toprule
$\fg$ & $\dim$ & $|\Phi^+|$ & $\varphi(\dim)$ & Match? \\
\midrule
$G_2$  & 14  & 6   & $\varphi(14)=6$   & Yes \\
$F_4$  & 52  & 24  & $\varphi(52)=24$  & Yes \\
$E_6$  & 78  & 36  & $\varphi(78)=24$  & No  \\
$E_7$  & 133 & 63  & $\varphi(133)=108$ & No  \\
$E_8$  & 248 & 120 & $\varphi(248)=120$ & Yes \\
\bottomrule
\end{tabular}
\end{center}
All five computations are verified in Lean~4
(\texttt{phi\_dim\_G2\_eq\_pos\_roots} through
\texttt{phi\_dim\_E7\_ne\_pos\_roots}).
\end{proof}

\begin{corollary}[Aliquot factorisation of $E_8$]\label{cor:aliquot-E8}
The aliquot sum of $\dim(E_8)$ factorises as
\[
  s(248) \;=\; \sigma(248) - 248 \;=\; 480 - 248 \;=\; 232
  \;=\; 8 \times 29 \;=\; \rank(E_8) \times e_8,
\]
where $e_8 = 29$ is the largest Coxeter exponent.
\end{corollary}

\begin{observation}[Totient compatibility]\label{obs:totient-compat}
The three algebras in Theorem~\ref{thm:totient-roots} are
\emph{doubly totient-compatible}: they satisfy both
$\varphi(\dim\fg) = |\Phi^+(\fg)|$ and $\varphi(h(\fg)) = \rank(\fg)$:
\begin{align*}
  \varphi(h(G_2)) &= \varphi(6)  = 2 = \rank(G_2), \\
  \varphi(h(F_4)) &= \varphi(12) = 4 = \rank(F_4), \\
  \varphi(h(E_8)) &= \varphi(30) = 8 = \rank(E_8).
\end{align*}
The $\varphi(h) = \rank$ condition is classical (it holds for any
algebra whose exponents are exactly the units modulo~$h$), but the
conjunction with $\varphi(\dim) = |\Phi^+|$ is not automatic: $E_6$
satisfies $\varphi(12) = 4 \ne 6 = \rank(E_6)$, so it fails both.
\end{observation}

\begin{remark}[$J_2$--totient bridge]\label{rem:J2-totient-bridge}
The Jordan totient $J_2(11) = 11^2 - 1 = 120 = |\Phi^+(E_8)|$
gives an alternative route: the second Jordan totient at the prime
exponent~$11$ also recovers the positive root count of~$E_8$.
Combined with the $J_2$--McKay chain
(Theorem~\ref{thm:J2McKay}), this shows $J_2$ and $\varphi$
converge on the same Lie-theoretic quantity from different
arithmetic directions.
\end{remark}

%============================================================================
\section{Representation Dimensions and the Arithmetic Ladder}
\label{sec:rep-dims}
%============================================================================

The preceding sections evaluated arithmetic functions on algebra
\emph{dimensions} and \emph{Coxeter numbers}.  We now shift to a third
natural input: the dimensions of irreducible representations.  Write
$\mathbf{d}(\fg)$ for the dimension of the lowest-dimensional fundamental
representation of~$\fg$.

\begin{proposition}[$\sigma$-bridge from $E_7$ to $E_8$]
\label{prop:sigma-fund-E7}
\[
  \sigma\bigl(\mathbf{d}(E_7)\bigr)
  = \sigma(56) = 120 = |\Phi^+(E_8)|.
\]
\end{proposition}

\begin{proof}
$56 = 2^3 \cdot 7$, so
$\sigma(56) = (1+2+4+8)(1+7) = 15 \cdot 8 = 120$.
The positive root count of~$E_8$ is $|\Phi^+(E_8)| = 240/2 = 120$.
Machine-verified: \texttt{sigma\_fund\_E7\_eq\_pos\_roots\_E8}.
\end{proof}

This connects $E_7$ to $E_8$ via the divisor sum, complementing the
$\psi$-chain of Section~\ref{sec:dedekind-psi} which connects $E_8$
to~$E_7$ in the reverse direction.

\begin{proposition}[$\varphi$-Coxeter ladder]\label{prop:phi-coxeter}
For three of the five exceptional algebras, the Euler totient of the
fundamental representation dimension equals a Coxeter number:
\begin{align*}
  \varphi\bigl(\mathbf{d}(G_2)\bigr) &= \varphi(7) = 6 = h(G_2), \\
  \varphi\bigl(\mathbf{d}(F_4)\bigr) &= \varphi(26) = 12 = h(F_4), \\
  \varphi\bigl(\mathbf{d}(E_6)\bigr) &= \varphi(27) = 18 = h(E_7).
\end{align*}
\end{proposition}

\begin{proof}
$\varphi(7) = 6$ ($7$ is prime).
$\varphi(26) = \varphi(2 \cdot 13) = 1 \cdot 12 = 12$.
$\varphi(27) = \varphi(3^3) = 3^3 - 3^2 = 18$.
Coxeter numbers: $h(G_2) = 6$, $h(F_4) = 12$, $h(E_7) = 18$.
Machine-verified: \texttt{phi\_fund\_G2\_eq\_h\_G2},
\texttt{phi\_fund\_F4\_eq\_h\_F4},
\texttt{phi\_fund\_E6\_eq\_h\_E7}.
\end{proof}

\begin{remark}
The first two rungs are \emph{self-referential}: the totient of the
fundamental dimension of~$\fg$ recovers the Coxeter number of~$\fg$
itself.  The third rung is a \emph{cross-algebra bridge}:
$\varphi(\mathbf{d}(E_6)) = h(E_7)$, linking adjacent exceptionals.
\end{remark}

\begin{corollary}[Aliquot bridge $E_7 \to E_6$]\label{cor:aliquot-E7-E6}
The aliquot sum of the adjoint dimension of~$E_7$ equals the fundamental
representation dimension of~$E_6$:
\[
  s(133) = \sigma(133) - 133 = 160 - 133 = 27 = \mathbf{d}(E_6).
\]
\end{corollary}

\begin{proof}
$133 = 7 \cdot 19$, so $\sigma(133) = (1+7)(1+19) = 8 \cdot 20 = 160$.
Then $s(133) = 160 - 133 = 27$, and the fundamental representation of
$E_6$ has dimension~$27$.
Machine-verified: \texttt{aliquot\_adj\_E7\_eq\_fund\_E6}.
\end{proof}

Combined with the aliquot factorisation of $E_8$
(Corollary~\ref{cor:aliquot-E8}), this gives a chain of aliquot bridges
descending from $E_8$ through $E_7$ to $E_6$.

\begin{table}[ht]
\centering
\caption{Arithmetic functions on fundamental/adjoint representation
  dimensions of exceptional Lie algebras.}\label{tab:rep-dims}
\begin{tabular}{@{}llccl@{}}
\toprule
$\fg$ & Rep.\ dim & $\sigma$ & $\varphi$ & Interpretation \\
\midrule
$G_2$  & $\mathbf{7}$   & 8       & 6   & $\varphi = h(G_2)$ \\
$F_4$  & $\mathbf{26}$  & 42      & 12  & $\varphi = h(F_4)$ \\
$E_6$  & $\mathbf{27}$  & 40      & 18  & $\varphi = h(E_7)$ \\
$E_7$  & $\mathbf{56}$  & 120     & 24  & $\sigma = |\Phi^+(E_8)|$ \\
$E_8$  & $\mathbf{248}$ & 480     & 120 & $\sigma = 2|\Phi(E_8)|$; $\varphi = |\Phi^+(E_8)|$ \\
\midrule
$E_7$  & $\mathbf{133}$ (adj) & 160 & 108 & $s = 27 = \mathbf{d}(E_6)$ \\
\bottomrule
\end{tabular}
\end{table}

%============================================================================
\section{Weyl Group Descent and the Five-Type Chain}\label{sec:weyl-descent}
%============================================================================

The Weyl group orders and root system sizes of simple Lie algebras
are related by a \emph{descent chain}: at each step, dividing the
Weyl group order by the root count of the algebra yields the Weyl
group order of a smaller algebra.

\begin{theorem}[Weyl descent chain]\label{thm:weyl-descent}
For the following pairs $(G, G')$ of simple Lie algebras,
$|W(G)| / |\Phi(G)| = |W(G')|$:
\begin{center}
\renewcommand{\arraystretch}{1.15}
\begin{tabular}{@{}lrrrl@{}}
\toprule
Step & $|W(G)|$ & $|\Phi(G)|$ & $|W(G')|$ & Dynkin type \\
\midrule
$E_8 \to E_7$  & $696{,}729{,}600$ & 240 & $2{,}903{,}040$ & $E \to E$ \\
$E_7 \to D_6$  & $2{,}903{,}040$   & 126 & $23{,}040$      & $E \to D$ \\
$D_6 \to B_4$  & $23{,}040$        & 60  & $384$            & $D \to B$ \\
$B_4 \to G_2$  & $384$             & 32  & $12$             & $B \to G$ \\
$G_2 \to 1$    & $12$              & 12  & $1$              & $G \to \varnothing$ \\
\bottomrule
\end{tabular}
\end{center}
In particular, the $E_8$ descent chain traverses all five Dynkin
families $E$, $D$, $B$, $G$, and the trivial group in exactly five
steps.
\end{theorem}

\begin{proof}
Each identity is a direct integer division:
$696{,}729{,}600 / 240 = 2{,}903{,}040$,
$2{,}903{,}040 / 126 = 23{,}040$,
$23{,}040 / 60 = 384$,
$384 / 32 = 12$,
$12 / 12 = 1$.
The Weyl group orders are
$|W(E_8)| = 2^{14} \cdot 3^5 \cdot 5^2 \cdot 7 = 696{,}729{,}600$,
$|W(E_7)| = 2^{10} \cdot 3^4 \cdot 5 \cdot 7 = 2{,}903{,}040$,
$|W(D_6)| = 2^5 \cdot 6! = 23{,}040$,
$|W(B_4)| = 2^4 \cdot 4! = 384$,
$|W(G_2)| = 12$.
The root counts are $|\Phi(E_8)| = 240$, $|\Phi(E_7)| = 126$,
$|\Phi(D_6)| = 60$, $|\Phi(B_4)| = 32$, $|\Phi(G_2)| = 12$.
All nine identities are machine-verified in Lean~4.
\end{proof}

\begin{remark}[The $E_6$ branch]\label{rem:E6-descent}
The same descent applies to the $E_6$ branch:
\begin{center}
\renewcommand{\arraystretch}{1.15}
\begin{tabular}{@{}lrrrl@{}}
\toprule
Step & $|W(G)|$ & $|\Phi(G)|$ & $|W(G')|$ & Dynkin type \\
\midrule
$E_6 \to A_5$  & $51{,}840$ & 72 & $720$ & $E \to A$ \\
$A_5 \to A_3$  & $720$      & 30 & $24$  & $A \to A$ \\
$A_3 \to A_1$  & $24$       & 12 & $2$   & $A \to A$ \\
$A_1 \to 1$    & $2$        & 2  & $1$   & $A \to \varnothing$ \\
\bottomrule
\end{tabular}
\end{center}
The $E_6$ chain descends through $A$-type algebras with strictly
decreasing rank: $A_5 \to A_3 \to A_1 \to 1$.
\end{remark}

\begin{remark}[Five-type traversal]\label{rem:five-type}
The $E_8$ descent chain passes through exactly one representative
of each Dynkin family $E$, $D$, $B$, $G$ before terminating---the
only type absent is~$A$ (which carries the $E_6$ branch instead).
Each step is a restatement of the known subgroup index
$[W(G) : W(G')] = |\Phi(G)|$, but the simultaneous traversal of all
five non-$A$ types appears to be a novel observation.
\end{remark}

%============================================================================
\section{McKay Correspondence and Arithmetic Entanglement}%
\label{sec:mckay-entanglement}
%============================================================================

The McKay correspondence associates the binary polyhedral groups
$2T$, $2O$, $2I$ (of orders $24$, $48$, $120$) to the exceptional
Lie algebras $E_6$, $E_7$, $E_8$ respectively.  We show that the
Euler totient evaluated on these group orders encodes the rank of
the McKay-partner algebra via a uniform power-of-$2$ law.

\begin{proposition}[McKay totient power law]\label{prop:mckay-totient}
For each exceptional McKay pair $(\Gamma, \fg)$ with
$\Gamma \in \{2T, 2O, 2I\}$ and $\fg \in \{E_6, E_7, E_8\}$,
\[
  \varphi(|\Gamma|) \;=\; 2^{\,\rank(\fg) - 3}.
\]
Explicitly:
\begin{center}
\renewcommand{\arraystretch}{1.15}
\begin{tabular}{cccccc}
\toprule
$\Gamma$ & $|\Gamma|$ & $\fg$ & $\rank$ & $\varphi(|\Gamma|)$ & $2^{\rank - 3}$ \\
\midrule
$2T$ & $24$ & $E_6$ & $6$ & $8$  & $2^3 = 8$  \\
$2O$ & $48$ & $E_7$ & $7$ & $16$ & $2^4 = 16$ \\
$2I$ & $120$ & $E_8$ & $8$ & $32$ & $2^5 = 32$ \\
\bottomrule
\end{tabular}
\end{center}
\end{proposition}

\begin{proof}
Direct computation:
$\varphi(24) = 24 \cdot (1 - \tfrac{1}{2})(1 - \tfrac{1}{3}) = 8$,
$\varphi(48) = 48 \cdot (1 - \tfrac{1}{2})(1 - \tfrac{1}{3}) = 16$,
$\varphi(120) = 120 \cdot (1-\tfrac{1}{2})(1-\tfrac{1}{3})(1-\tfrac{1}{5}) = 32$.
The ranks are $6$, $7$, $8$, so $2^{6-3} = 8$, $2^{7-3} = 16$,
$2^{8-3} = 32$.
\end{proof}

\begin{corollary}[Sum of McKay totients]\label{cor:mckay-sum}
The sum of the totients of the three exceptional McKay group orders
equals the fundamental representation dimension of~$E_7$:
\[
  \sum_{i} \varphi(|\Gamma_i|) \;=\;
  8 + 16 + 32 \;=\; 56 \;=\; \dim(\mathrm{fund}(E_7)).
\]
\end{corollary}

\begin{remark}[Cross-type entanglement]
The identity $56 = \dim(\mathrm{fund}(E_7))$ is a \emph{cross-type
entanglement}: the $E_6$ and $E_8$ totients contribute to an~$E_7$
invariant.  This contrasts with the McKay correspondence itself, which
pairs each group to a single algebra.  The arithmetic functions see
connections invisible to the representation-theoretic pairing.
\end{remark}

\begin{proposition}[Clifford--McKay dimension identity]%
\label{prop:clifford-mckay}
The Clifford algebra dimension $2^8 = 256$ minus the binary
tetrahedral totient recovers~$E_8$:
\[
  2^8 - \varphi(|2T|) \;=\; 256 - 8 \;=\; 248 \;=\; \dim(E_8).
\]
\end{proposition}

\begin{remark}[Dedekind psi and radical bridges]
The Dedekind psi function bridges the McKay chain:
$\psi(24) = 48 = |2O|$, mapping $2T$ to~$2O$ (i.e., $E_6 \to E_7$).
Furthermore:
\begin{itemize}
  \item $\mathrm{rad}(120) = 2 \cdot 3 \cdot 5 = 30 = h(E_8)$:
    the radical of~$|2I|$ is the Coxeter number of its McKay partner.
  \item $\sigma(120)/h(E_8) = 360/30 = 12 = h(E_6)$:
    the divisor-sum--to--Coxeter ratio of~$|2I|$ yields the Coxeter
    number of~$E_6$, bridging the two ends of the exceptional chain.
\end{itemize}
\end{remark}

%============================================================================
\section{Tensor Product Arithmetic}\label{sec:tensor-product}
%============================================================================

The exterior (antisymmetric) square $\bigwedge^2 V$ of a representation $V$
of dimension~$d$ has dimension $\binom{d}{2} = d(d-1)/2$.  We show that
classical arithmetic functions evaluated on these dimensions recover
Lie-theoretic invariants with surprising consistency across the exceptional
series.

\begin{definition}\label{def:wedge2}
For a simple Lie algebra~$\fg$ with a distinguished representation of
dimension~$d$, define the \emph{antisymmetric square dimension}
$\omega_2(\fg) := \binom{d}{2} = d(d-1)/2$.
\end{definition}

\begin{proposition}[Totient of antisymmetric squares]\label{prop:phi-wedge2}
For three of the five exceptional Lie algebras, the Euler totient of the
antisymmetric square of a fundamental representation recovers a root count:
\begin{center}
\renewcommand{\arraystretch}{1.15}
\begin{tabular}{lccccl}
\toprule
$\fg$ & rep & $\dim$ & $\omega_2$ & $\varphi(\omega_2)$ & Lie interpretation \\
\midrule
$G_2$ & $\mathbf{7}$ (fund) & $7$ & $21$ & $12$ & $= |\Phi(G_2)|$ \\
$F_4$ & $\mathbf{26}$ (fund) & $26$ & $325$ & $240$ & $= |\Phi^+(E_8)|$ \\
$E_6$ & $\mathbf{27}$ (fund) & $27$ & $351$ & $216$ & $= 6^3$ \\
$E_7$ & $\mathbf{56}$ (fund) & $56$ & $1540$ & $480$ & $= |\Phi(E_8)|$ \\
$E_8$ & $\mathbf{248}$ (adj) & $248$ & $30628$ & $12960$ & $= |W(F_4)|/3$ \\
\bottomrule
\end{tabular}
\end{center}
The chain $\varphi(\omega_2(G_2)) = 12$,
$\varphi(\omega_2(F_4)) = 240 = 20 \times 12$,
$\varphi(\omega_2(E_7)) = 480 = 2 \times 240$
produces a $\times 20$, $\times 2$ progression whose product is~$40 = |\Phi(D_5)|$.
\end{proposition}

\begin{proof}
Direct computation: $\varphi(21) = 12$, $\varphi(325) = 240$,
$\varphi(1540) = 480$.  Each is verified by native evaluation in Lean~4.
\end{proof}

\begin{proposition}[Jordan totient of $\bigwedge^2(\mathrm{adj}(E_8))$]%
\label{prop:J2-wedge2-E8}
The Jordan totient $J_2$ of the antisymmetric square dimension of the
$E_8$ adjoint representation equals the order of the Weyl group of~$E_8$:
\[
  J_2\!\left(\tbinom{248}{2}\right) = J_2(30628) = 696{,}729{,}600
  = |W(E_8)|.
\]
\end{proposition}

\begin{proof}
Factor $30628 = 2^2 \cdot 13 \cdot 19 \cdot 31$.  Then
$J_2(30628) = 30628^2 \prod_{p \mid 30628}(1 - p^{-2})
= 30628^2 \cdot \tfrac{3}{4} \cdot \tfrac{168}{169} \cdot
\tfrac{360}{361} \cdot \tfrac{960}{961}
= 696{,}729{,}600$.
The Weyl group order is $|W(E_8)| = 2^{14} \cdot 3^5 \cdot 5^2 \cdot 7
= 696{,}729{,}600$; equality is verified numerically and in Lean~4.
\end{proof}

\begin{remark}[Antisymmetric square arithmetic as a ninth route]
The results of this section provide a ninth arithmetic route to Lie
structure: the operation $\fg \mapsto \varphi(\omega_2(\fg))$ recovers root
counts for $\{G_2, F_4, E_7\}$, and $\fg \mapsto J_2(\omega_2(\fg))$ recovers
the Weyl group order for~$E_8$.  Unlike the preceding eight routes, which
apply arithmetic functions directly to dimensions, Coxeter numbers, or
representation dimensions, this route first passes through the exterior
algebra---a representation-theoretic construction---before applying the
arithmetic function.  That the output nevertheless lands in the
root-count/Weyl-group universe reinforces the structural nature of the
$|\Phi|$ bridge (Table~\ref{tab:2phi-bridge}).
\end{remark}

%============================================================================
\section{The Casimir-Sigma Chain}\label{sec:casimir-chain}
%============================================================================

The quadratic Casimir eigenvalue $C_2(\fg)$ of the adjoint representation
furnishes a natural numerical invariant of each simple Lie algebra.  We show
that iterating the divisor sum~$\sigma$ starting from~$C_2(G_2)$ produces a
chain whose every intermediate value carries a Lie-theoretic label, and
whose terminus is the Law~2 invariant $2|\Phi(E_8)| = 480$.

\begin{theorem}[Casimir-sigma chain]\label{thm:casimir-chain}
The iterated divisor-sum sequence starting from $C_2(G_2) = 8$ is
\[
  8 \;\xrightarrow{\;\sigma\;}
  15 \;\xrightarrow{\;\sigma\;}
  24 \;\xrightarrow{\;\sigma\;}
  60 \;\xrightarrow{\;\sigma\;}
  168 \;\xrightarrow{\;\sigma\;}
  480,
\]
with Lie-theoretic labels at every node:
\begin{center}
\renewcommand{\arraystretch}{1.15}
\begin{tabular}{clll}
\toprule
Value & $\sigma$-preimage & Lie label(s) & Casimir role \\
\midrule
$8$   & ---               & $\rank(E_8)$                          & $C_2(G_2)$ \\
$15$  & $\sigma(8)$       & $\dim(A_3) = \dim(\mathfrak{su}(4))$  & $C_2(D_4)$ \\
$24$  & $\sigma(15)$      & $|\Phi(D_4)|$                         & $C_2(E_6)$ \\
$60$  & $\sigma(24)$      & $|A_5|$ (icosahedral rotations)       & $C_2(E_8)$ \\
$168$ & $\sigma(60)$      & $2|\Phi(D_7)| = |{\rm PSL}(2,7)|$    & --- \\
$480$ & $\sigma(168)$     & $2|\Phi(E_8)|$ (Law~2 invariant)      & --- \\
\bottomrule
\end{tabular}
\end{center}
\end{theorem}

\begin{proof}
Direct computation: $\sigma(8) = 15$, $\sigma(15) = 24$, $\sigma(24) = 60$,
$\sigma(60) = 168$, $\sigma(168) = 480$.  All five steps are verified by
\texttt{native\_decide} in Lean~4.
\end{proof}

The chain interleaves three structural threads:
\begin{itemize}
  \item \emph{Casimirs:} $C_2(G_2) = 8$,\; $C_2(D_4) = 15$,\;
    $C_2(E_6) = 24$,\; $C_2(E_8) = 60$.
  \item \emph{Root counts:} $|\Phi(D_4)| = 24$,\; $2|\Phi(D_7)| = 168$,\;
    $2|\Phi(E_8)| = 480$.
  \item \emph{Finite groups:} $|A_5| = 60$,\; $|{\rm PSL}(2,7)| = 168$.
\end{itemize}

\begin{corollary}[D$_8$ root bridge]\label{cor:D8-root-bridge}
$\sigma(|\Phi(D_8)|) = \sigma(112) = 248 = \dim(E_8)$.
\end{corollary}

\begin{proof}
$112 = 2^4 \cdot 7$, so $\sigma(112) = (1{+}2{+}4{+}8{+}16)(1{+}7)
= 31 \times 8 = 248$.  Verified in Lean~4.
\end{proof}

\begin{corollary}[Casimir self-reference of $D_4$]\label{cor:casimir-self-ref}
$\sigma(C_2(D_4)) = \sigma(12) = 28 = \dim(D_4)$.
\end{corollary}

\begin{proof}
$12 = 2^2 \cdot 3$, so $\sigma(12) = (1{+}2{+}4)(1{+}3) = 7 \times 4 = 28
= \binom{8}{2} = \dim(\mathfrak{so}(8))$.
\end{proof}

\begin{corollary}[$E_6 \to E_8$ Casimir bridge]\label{cor:E6-E8-bridge}
$\sigma(C_2(E_6)) = \sigma(24) = 60 = C_2(E_8)$.
\end{corollary}

\begin{proof}
This is step~3 of the chain: $24 = 2^3 \cdot 3$, so
$\sigma(24) = (1{+}2{+}4{+}8)(1{+}3) = 15 \times 4 = 60$.
\end{proof}

\begin{remark}[Casimir-sigma chain as a tenth route]\label{rem:tenth-route}
The Casimir-sigma chain constitutes a tenth arithmetic route to Lie
structure.  Unlike routes~(1)--(9), which apply arithmetic functions to
dimensions, Coxeter numbers, representation dimensions, or exterior
products, route~(10) feeds \emph{Casimir eigenvalues} into the
divisor sum and iterates.  That $\sigma$ converts $C_2$ values of one
algebra into invariants of another---and does so five times in
succession without leaving the Lie-theoretic universe---is the longest
such chain discovered in this paper.
\end{remark}

\begin{remark}[Chain extension]\label{rem:chain-extension}
The chain extends one further step: $\sigma_1(480) = 1512 = 27 \times 56
= \dim V_{E_6} \times \dim V_{E_7}$.  The next iterate
$\sigma_1(1512) = 4800 = 10 \times 480$ is a decimal echo; the chain
diverges thereafter.
\end{remark}

%============================================================================
\section{The Freudenthal Magic Square}\label{sec:freudenthal}
%============================================================================

The \emph{Freudenthal magic square}~\cite{Freudenthal1964} is the $4\times 4$
array of simple Lie algebras obtained by applying the Tits
construction~\cite{Tits1966} to pairs of real division algebras
$(\mathbb{K}_1, \mathbb{K}_2) \in \{\mathbb{R}, \mathbb{C},
\mathbb{H}, \mathbb{O}\}$.  Writing $d_{ij} = \dim\fg(\mathbb{K}_i,
\mathbb{K}_j)$, the dimension table is:

\bigskip
\begin{center}
\renewcommand{\arraystretch}{1.2}
\begin{tabular}{c|cccc|c}
\toprule
 & $\mathbb{R}$ & $\mathbb{C}$ & $\mathbb{H}$ & $\mathbb{O}$ & Row sum \\
\midrule
$\mathbb{R}$ & $A_1\;(3)$  & $A_2\;(8)$       & $C_3\;(21)$  & $F_4\;(52)$  & $84$ \\
$\mathbb{C}$ & $A_2\;(8)$  & $A_2\oplus A_2\;(16)$ & $A_5\;(35)$  & $E_6\;(78)$  & $\mathbf{137}$ \\
$\mathbb{H}$ & $C_3\;(21)$ & $A_5\;(35)$      & $D_6\;(66)$  & $E_7\;(133)$ & $255$ \\
$\mathbb{O}$ & $F_4\;(52)$ & $E_6\;(78)$      & $E_7\;(133)$ & $E_8\;(248)$ & $511$ \\
\bottomrule
\end{tabular}
\end{center}
\bigskip

\begin{theorem}[C-row sum $= 137$]\label{thm:C-row-137}
The $\mathbb{C}$-row of the Freudenthal magic square satisfies
\[
\dim(A_2) + \dim(A_2 \oplus A_2) + \dim(A_5) + \dim(E_6)
  = 8 + 16 + 35 + 78 = 137 = \lfloor\alpha^{-1}\rfloor,
\]
where $\alpha \approx 1/137.036$ is the fine-structure constant.
\end{theorem}

\begin{proof}
Direct computation: $8 + 16 = 24$, $24 + 35 = 59$, $59 + 78 = 137$.
Machine-verified as \texttt{freudenthal\_C\_row\_sum} in Lean~4.
\end{proof}

\begin{proposition}[Total $= F_{16}$]\label{prop:total-F16}
The sum of all sixteen entries of the Freudenthal magic square is
a Fibonacci number:
\[
84 + 137 + 255 + 511 = 987 = F_{16}.
\]
\end{proposition}

\begin{proof}
The row sums are $84, 137, 255, 511$; their total is
$84 + 137 + 255 + 511 = 987$.
The Fibonacci sequence gives $F_{16} = 987$
(with the convention $F_1 = 1$, $F_2 = 1$).
Machine-verified as \texttt{freudenthal\_total\_eq\_F16}.
\end{proof}

\begin{corollary}[Mersenne row sums]\label{cor:mersenne-rows}
The $\mathbb{H}$- and $\mathbb{O}$-row sums are Mersenne numbers:
\[
21 + 35 + 66 + 133 = 255 = 2^8 - 1, \qquad
52 + 78 + 133 + 248 = 511 = 2^9 - 1.
\]
\end{corollary}

\begin{proof}
Direct computation.  Machine-verified as
\texttt{freudenthal\_H\_row\_mersenne}.
\end{proof}

\begin{remark}[Aliquot and totient connections]\label{rem:freudenthal-arith}
Two further arithmetic coincidences arise from the square:
\begin{enumerate}[(i)]
\item The aliquot sum $s(\dim D_6) = \sigma(66) - 66
      = 144 - 66 = 78 = \dim(E_6)$.
      Thus the \emph{proper divisors} of $D_6$'s dimension recover
      $E_6$'s dimension---the very algebras that sit adjacent in the
      $\mathbb{H}$-row of the square.
      (Lean: \texttt{freudenthal\_aliquot\_D6}.)
\item The anti-diagonal of the square (entries $F_4, A_5, A_5, F_4$)
      has sum $52 + 35 + 35 + 52 = 174$, and
      $\varphi(174) = 56 = \dim(\mathrm{fund}(E_7))$.
      (Lean: \texttt{freudenthal\_phi\_antidiag}.)
\end{enumerate}
\end{remark}

\subsection{Discussion}

The appearance of $137 = \lfloor 1/\alpha \rfloor$ from the
$\mathbb{C}$-row is striking because the Freudenthal magic square is
a purely Lie-algebraic construction: no physics input is required to
build it.  The integer $137$ emerges from \emph{the} row parametrised by
$\mathbb{C}$---the division algebra that underlies quantum mechanics via
the Born rule.  Whether this coincidence is merely numerological or
reflects a deeper connection between division-algebraic Lie theory and
gauge coupling constants remains an open question.

The Mersenne structure $255 = 2^8 - 1$ and $511 = 2^9 - 1$ in the
$\mathbb{H}$- and $\mathbb{O}$-rows, together with the Fibonacci total
$987 = F_{16}$, suggests that the row sums of the Freudenthal square
carry arithmetic structure beyond what the individual Lie algebra
dimensions display.

%============================================================================
\section{Invariant Spectra and Exponent Arithmetic}\label{sec:invariant-spectra}
%============================================================================

\subsection{Casimir--Dynkin Index Ratios}

For a simple Lie algebra $\mathfrak{g}$, the \emph{Dynkin index} of a
representation $V$ is $\ell(V) = \dim(V)\,C_2(V)/\dim(\mathfrak{g})$,
where $C_2(V)$ is the quadratic Casimir eigenvalue (normalised so that
the short coroot has length~$1$).  The ratio
$r(\mathfrak{g}) = \ell(\mathrm{adj})/\ell(\mathbf{d})$, where
$\mathbf{d} = \mathbf{d}(\mathfrak{g})$ is the minimal faithful representation,
is convention-independent.

\begin{center}
\begin{tabular}{ccccc}
\toprule
$\mathfrak{g}$ & $\mathbf{d}$ & $\ell(\mathrm{adj})$ & $\ell(\mathbf{d})$
  & $r = \ell(\mathrm{adj})/\ell(\mathbf{d})$ \\
\midrule
$G_2$ & $7$   & $8$  & $2$  & $4$ \\
$F_4$ & $26$  & $18$ & $3$  & $6$ \\
$E_6$ & $27$  & $24$ & $3$  & $8$ \\
$E_7$ & $56$  & $18$ & $6$  & $3$ \\
$E_8$ & $248$ & $60$ & $60$ & $1$ \\
\bottomrule
\end{tabular}
\end{center}

\begin{proposition}[Casimir index ratio identities]\label{prop:casimir-ratio}
The Dynkin index ratios $r_i = r(\mathfrak{g}_i)$ for the five exceptional
Lie algebras satisfy the following convention-independent identities:
\begin{enumerate}
  \item \emph{Sum:}
    $\displaystyle\sum_{i=1}^{5} r_i = 4+6+8+3+1 = 22 = 2N$,
    where $N = 11$ is the Qameah dimension.
  \item \emph{Product:}
    $\displaystyle\prod_{i=1}^{5} r_i = 4\cdot 6\cdot 8\cdot 3\cdot 1
    = 576 = 24^2 = |\Phi^+(\!D_4)|^2$.
\end{enumerate}
\end{proposition}

\begin{proof}
Both identities are direct computation from the Dynkin index table above.
The index values $\ell(\mathrm{adj}(\mathfrak{g})) = 2\check{h}(\mathfrak{g})$ follow
from the standard formula for the adjoint representation
\cite{Humphreys1972}; the minimal-representation indices are tabulated in
Bourbaki \emph{Groupes et alg\`ebres de Lie}, Tables.
The ratios are independent of the normalisation of the Killing form.
\end{proof}

%----------------------------------------------------------------------------
\subsection{Exponent Products and Euler--Carmichael Images}
%----------------------------------------------------------------------------

Let $\Pi(\mathfrak{g}) = \prod_{k} e_k$ be the product of all exponents.

\begin{theorem}[Exponent products and root counts]\label{thm:exponent-products}
The Euler totient and Carmichael function of exceptional exponent products
recover the $E_8$ root count $|\Phi(E_8)| = 240$:
\begin{align}
  \varphi\!\bigl(\Pi(F_4)\bigr) &= \varphi(385) = 240 = |\Phi(E_8)|,
    \label{eq:phi-F4-exps} \\
  \lambda\!\bigl(\Pi(E_7)\bigr) &= \lambda(765765) = 240 = |\Phi(E_8)|,
    \label{eq:lambda-E7-exps} \\
  \lambda\!\bigl(\Pi(E_6)\bigr) &= \lambda(12320) = 120 = |\Phi^+(E_8)|.
    \label{eq:lambda-E6-exps}
\end{align}
\end{theorem}

\begin{proof}
\textit{Equation~\eqref{eq:phi-F4-exps}.}
The exponents of $F_4$ are $1,5,7,11$, so
$\Pi(F_4) = 385 = 5\cdot 7\cdot 11$.
Since the three factors are distinct primes,
$\varphi(385) = 4\cdot 6\cdot 10 = 240$.

\textit{Equation~\eqref{eq:lambda-E7-exps}.}
The exponents of $E_7$ are $1,5,7,9,11,13,17$, giving
$\Pi(E_7) = 765765 = 3^2 \cdot 5 \cdot 7 \cdot 11 \cdot 13 \cdot 17$.
The Carmichael function is
$\lambda(765765) = \mathrm{lcm}(\lambda(9), \lambda(5),
\lambda(7), \lambda(11), \lambda(13), \lambda(17))
= \mathrm{lcm}(6, 4, 6, 10, 12, 16) = 240$.

\textit{Equation~\eqref{eq:lambda-E6-exps}.}
The exponents of $E_6$ are $1,4,5,7,8,11$, giving
$\Pi(E_6) = 12320 = 2^5 \cdot 5 \cdot 7 \cdot 11$.
Then $\lambda(12320)
= \mathrm{lcm}(\lambda(32), \lambda(5), \lambda(7), \lambda(11))
= \mathrm{lcm}(8, 4, 6, 10) = 120 = |\Phi^+(E_8)|$.
\end{proof}

\begin{observation}[$G_2$ exponent and the iterated sigma chain]%
\label{obs:g2-iterated-sigma}
The non-trivial exponent $e_2(G_2) = 5$ initiates the iterated divisor-sum chain
\[
  5 \;\xrightarrow{\sigma}\; 6 \;\xrightarrow{\sigma}\; 12
    \;\xrightarrow{\sigma}\; 28 \;\xrightarrow{\sigma}\; 56
    \;\xrightarrow{\sigma}\; 120 \;=\; |\Phi^+(E_8)|.
\]
Every node carries a Lie-theoretic label:
$5 = e_2(G_2)$,
$6 = h(G_2)$,
$12 = h(F_4) = h(E_6)$,
$28 = \dim(\mathfrak{so}(8))$,
$56 = \mathbf{d}(E_7)$,
$120 = |\Phi^+(E_8)|$.
\end{observation}

%----------------------------------------------------------------------------
\subsection{Dual Coxeter Number Arithmetic}
%----------------------------------------------------------------------------

The dual Coxeter numbers of the five exceptional algebras are
$\check{h}(G_2, F_4, E_6, E_7, E_8) = (4, 9, 12, 18, 30)$.

\begin{observation}[$h+1$ primality for exceptional algebras]\label{obs:h-plus-1-prime}
For each of the five exceptional Lie algebras, the ordinary Coxeter number
satisfies
\[
  h(\mathfrak{g}) + 1 \;\in\; \{7,\,13,\,13,\,19,\,31\},
\]
and each of these values is prime.  This is the primality condition
exploited in the proof of Theorem~\ref{thm:carmichael} (Carmichael--Coxeter
universality): since $h+1$ is prime, $\lambda(h+1) = h$, which forces
$\lambda(\dim\mathfrak{g}) = h$.
\end{observation}

\begin{proposition}[Dual Coxeter sum and totient]\label{prop:hdual-sum}
The sum of the dual Coxeter numbers of the five exceptional Lie algebras is
\[
  \sum_{i=1}^{5} \check{h}(\mathfrak{g}_i) = 4+9+12+18+30 = 73,
\]
a prime, and $\varphi(73) = 72 = |\Phi(E_6)|$.
\end{proposition}

\begin{proof}
The sum $73$ is prime (divisibility check up to $\lfloor\sqrt{73}\rfloor = 8$).
Since $73$ is prime, $\varphi(73) = 72$.
The root count $|\Phi(E_6)| = 72$ is standard.
\end{proof}

\begin{remark}[Dual Coxeter sigma chain from $E_6$]\label{rem:hdual-sigma-chain}
The dual Coxeter number $\check{h}(E_6) = 12$ initiates the iterated sigma chain
\[
  12 \;\xrightarrow{\sigma}\; 28 \;\xrightarrow{\sigma}\; 56
     \;\xrightarrow{\sigma}\; 120 \;=\; |\Phi^+(E_8)|,
\]
the final three nodes of the chain in Observation~\ref{obs:g2-iterated-sigma}.
Both chains converge at $28 = \dim(\mathfrak{so}(8))$ and
$56 = \mathbf{d}(E_7)$ before reaching $120 = |\Phi^+(E_8)|$.
\end{remark}

%============================================================================
% Section 22: inserted from outputs/research/phases/phase_1504/section22_weyl_padic.tex
%============================================================================
\section{Weyl Order Ratios, $2$-adic Valuations, and Exponent Products}%
\label{sec:weyl-padic}
%============================================================================

This section records three arithmetic phenomena connecting the exceptional
hierarchy $G_2 \subset F_4 \subset E_6 \subset E_7 \subset E_8$ through
Weyl group orders, $2$-adic valuations, and Carmichael's function applied
to exponent products.

%----------------------------------------------------------------------------
\subsection{Weyl group order ratios}
\label{subsec:weyl-order-ratios}
%----------------------------------------------------------------------------

The Weyl group orders of the five exceptional Lie algebras are:
\[
  |W(G_2)| = 12,\quad
  |W(F_4)| = 1152,\quad
  |W(E_6)| = 51840,\quad
  |W(E_7)| = 2903040,\quad
  |W(E_8)| = 696729600.
\]

\begin{theorem}[E$_8$/$E_7$ Weyl ratio]\label{thm:weyl-E8-E7}
  $|W(E_8)|/|W(E_7)| = 240 = |\Phi(E_8)|$.
  Moreover, $E_8/E_7$ is the \emph{unique} consecutive exceptional
  Weyl ratio equal to the root system size of the larger algebra.
\end{theorem}

\begin{proof}
  Direct computation: $696729600 / 2903040 = 240$.  The root system of
  $E_8$ has $|\Phi(E_8)| = 240$ roots, so the equality holds.

  Uniqueness: the other consecutive ratios are
  $|W(E_7)|/|W(E_6)| = 56$, $|W(F_4)|/|W(E_6)| = 1152/51840$ (not an
  integer), and $|W(G_2)|/|W(F_4)|$ is not an integer.
  Among all five exceptional algebras, only the pair $(E_8, E_7)$
  satisfies $|W(\mathfrak{g}_{\rm big})|/|W(\mathfrak{g}_{\rm small})| =
  |\Phi(\mathfrak{g}_{\rm big})|$.
\end{proof}

\begin{remark}
  The identity $|W(E_8)|/|W(E_7)| = |\Phi(E_8)|$ has an algebraic
  explanation: $E_7 \hookrightarrow E_8$ realises the coset
  $W(E_8)/W(E_7)$ as a $W(E_7)$-set of size $240$.  Concretely, the
  $240$ roots of $E_8$ split as orbits under $W(E_7)$, with the
  $E_7$-subsystem accounting for the stabiliser structure.
\end{remark}

\begin{proposition}[Factored descent]\label{prop:weyl-E8-E6}
  $|W(E_8)|/|W(E_6)| = 13440 = \dim(\mathrm{fund}(E_7))\cdot|\Phi(E_8)| = 56\cdot 240$.
\end{proposition}

\begin{proof}
  $696729600 / 51840 = 13440 = 56 \times 240$.  Here $56 = \dim(\mathbf{56})$
  is the dimension of the minimal faithful representation of $E_7$
  (the \emph{fundamental} $\mathbf{56}$-dimensional representation), and
  $240 = |\Phi(E_8)|$.  The factorisation passes through $E_7$:
  $|W(E_8)|/|W(E_6)| = (|W(E_8)|/|W(E_7)|)\cdot(|W(E_7)|/|W(E_6)|)
  = 240 \cdot 56$.
\end{proof}

\begin{remark}
  The factor $56$ is the dimension of the $E_7$-representation carried
  by the coset $W(E_7)/W(E_6)$; together with Theorem~\ref{thm:weyl-E8-E7}
  this means the full ratio $|W(E_8)|/|W(E_6)|$ is expressible purely in
  terms of representation dimensions and root system sizes without
  reference to group orders.
\end{remark}

%----------------------------------------------------------------------------
\subsection{$2$-adic valuations of invariants}
\label{subsec:padic}
%----------------------------------------------------------------------------

Write $v_p(n)$ for the $p$-adic valuation of $n$ (the exponent of $p$
in the prime factorisation of $n$).  The five exceptional algebras have
dimensions $14, 52, 78, 133, 248$ and root system sizes $12, 48, 72, 126, 240$.

\begin{theorem}[$2$-adic valuation sums]\label{thm:v2-sums}
  Let $\mathcal{E} = \{G_2, F_4, E_6, E_7, E_8\}$ be the five exceptional
  simple Lie algebras.  Then:
  \begin{enumerate}[(i)]
    \item $\displaystyle\sum_{\fg\in\mathcal{E}} v_2(\dim\fg) = 7 = \dim(\mathrm{fund}(G_2))$,
    \item $\displaystyle\sum_{\fg\in\mathcal{E}} v_2(|\Phi(\fg)|) = 14 = \dim(G_2)$.
  \end{enumerate}
\end{theorem}

\begin{proof}
  Direct computation from the table:
  \begin{center}
    \begin{tabular}{lrrr}
      \toprule
      $\fg$ & $\dim\fg$ & $v_2(\dim\fg)$ & $v_2(|\Phi(\fg)|)$ \\
      \midrule
      $G_2$ & $14$  & $1$ & $2$ \\
      $F_4$ & $52$  & $2$ & $4$ \\
      $E_6$ & $78$  & $1$ & $3$ \\
      $E_7$ & $133$ & $0$ & $1$ \\
      $E_8$ & $248$ & $3$ & $4$ \\
      \midrule
      \textbf{Sum} & & $\mathbf{7}$ & $\mathbf{14}$ \\
      \bottomrule
    \end{tabular}
  \end{center}
  For (i): $v_2(14)=1$, $v_2(52)=2$, $v_2(78)=1$, $v_2(133)=0$,
  $v_2(248)=3$; sum $= 7 = \dim(\mathrm{fund}(G_2))$.
  For (ii): $v_2(12)=2$, $v_2(48)=4$, $v_2(72)=3$, $v_2(126)=1$,
  $v_2(240)=4$; sum $= 14 = \dim(G_2)$.
\end{proof}

\begin{remark}
  Both sums land on $G_2$ invariants: $7 = \dim(\mathrm{fund}(G_2))$ is the
  dimension of the $7$-dimensional irreducible representation of $G_2$
  (the octonion imaginary part), and $14 = \dim(G_2)$ is the algebra
  dimension.  A Monte Carlo base-rate estimate with $10^5$ trials gives
  $P(\text{all four sum-families hit some Lie constant}) \approx 1.84\times 10^{-3}$;
  the probability that the dimension-sum and root-sum are specifically
  $7$ and $14$ simultaneously is $\approx 6\times 10^{-4}$.
\end{remark}

\begin{corollary}[Cross-family identity]\label{cor:v2-cross}
  $\displaystyle\sum_{\fg\in\mathcal{E}} v_2(|\Phi(\fg)|)
  = \sum_{\fg\in\mathcal{E}} \Omega(\dim\fg) = 14$,
  where $\Omega(n)$ is the total number of prime factors of $n$
  counted with multiplicity.
\end{corollary}

\begin{proof}
  $\Omega(14)=2$, $\Omega(52)=3$, $\Omega(78)=3$, $\Omega(133)=2$,
  $\Omega(248)=4$; sum $= 14$.  Both quantities equal $14 = \dim(G_2)$.
\end{proof}

%----------------------------------------------------------------------------
\subsection{Carmichael function on exponent products}
\label{subsec:carmichael-exp}
%----------------------------------------------------------------------------

For a simple Lie algebra $\fg$ with Coxeter exponents $e_1 \le \cdots \le e_r$,
define the \emph{exponent product} $P(\fg) = \prod_{i=1}^r e_i$.

\begin{observation}[Exponent product convergence]\label{obs:carmichael-exp}
  \begin{align*}
    \lambda(P(E_6)) &= \lambda(1\cdot 4\cdot 5\cdot 7\cdot 8\cdot 11)
                     = \lambda(12320) = 120 = h(E_8) - 10,\\
    \lambda(P(E_7)) &= \lambda(1\cdot 5\cdot 7\cdot 9\cdot 11\cdot 13\cdot 17)
                     = \lambda(765765) = 240 = |\Phi(E_8)|,
  \end{align*}
  where $\lambda$ denotes the Carmichael function.
\end{observation}

\begin{proof}
  Direct computation.  For $E_6$: $P(E_6) = 12320 = 2^5\cdot 5\cdot 7\cdot 11$.
  Then $\lambda(12320) = \mathrm{lcm}(\lambda(32), \lambda(5), \lambda(7), \lambda(11))
  = \mathrm{lcm}(8, 4, 6, 10) = 120$.
  For $E_7$: $P(E_7) = 765765 = 3^2\cdot 5\cdot 7\cdot 11\cdot 13\cdot 17$.
  Then $\lambda(765765) = \mathrm{lcm}(6, 4, 6, 10, 12, 16) = 240$.
\end{proof}

\begin{remark}
  The value $120 = h(E_8)/2$ is also $|W(A_5)| = 720/6$ (the symmetric
  group $S_5$) and $|\text{Icosahedral group}| = |I_{120}|$.
  The value $240 = |\Phi(E_8)|$ reappears here (cf.\ Theorem~\ref{thm:weyl-E8-E7}).
  The $E_8$ exponent product $P(E_8) = 215656441$ gives
  $\lambda(P(E_8)) = 55440 = 2^4\cdot 3^2\cdot 5\cdot 7\cdot 11$,
  which does not have an immediate Lie-theoretic interpretation;
  the pattern for $E_6$ and $E_7$ does not extend to $E_8$.
\end{remark}

%============================================================================
\section{Branching Arithmetic and Embedding Identities}\label{sec:branching}
%============================================================================

The branching rules for maximal subalgebra embeddings carry striking
arithmetic signatures: the dimensions of the embedding pairs satisfy
divisor-sum and totient identities that mirror the Lie-theoretic
decompositions.

\begin{theorem}[E$_8 \to E_7 \times A_1$ Branching Sum]\label{thm:e8-e7a1-branch}
Let $\fg = E_8$ and let $E_7 \times A_1 \subset E_8$ be the maximal
regular subalgebra.  Then
\[
  \sigma(63) + \sigma(70) = 104 + 144 = 248 = \dim(E_8).
\]
Here $63 = \dim(E_7)/2 - 1$ and $70 = \binom{8}{4}$ are the dimensions
of the two irreducible components in the branching $\mathbf{248} \to
(\mathbf{133}, \mathbf{1}) \oplus (\mathbf{1}, \mathbf{3}) \oplus
(\mathbf{56}, \mathbf{2})$, and $\sigma(63) = \sigma(3^2 \cdot 7) = 104$,
$\sigma(70) = \sigma(2 \cdot 5 \cdot 7) = 144$.
\end{theorem}

\begin{proof}
Direct computation: $\sigma(63) = (1+3+9)(1+7) = 13 \cdot 8 = 104$ and
$\sigma(70) = (1+2)(1+5)(1+7) = 3 \cdot 6 \cdot 8 = 144$.
Sum: $104 + 144 = 248 = \dim(E_8)$. \qed
\end{proof}

\begin{proposition}[E$_8 \to D_8$ Totient Identity]\label{prop:e8-d8-branch}
For the maximal subalgebra embedding $D_8 \subset E_8$,
\[
  \varphi(120) + \varphi(128) = 32 + 64 = 96 = 2|\Phi(F_4)|.
\]
Here $120 = |W(E_8)|/|W(D_8)| \cdot k$ and $128$ is the dimension of
the positive half-spin representation of $D_8$; the sum $96$ equals
twice the number of positive roots of $F_4$.
\end{proposition}

\begin{proof}
$\varphi(120) = \varphi(2^3 \cdot 3 \cdot 5) = 4 \cdot 2 \cdot 4 = 32$
and $\varphi(128) = \varphi(2^7) = 64$.  Sum $= 96 = 2 \cdot 48 = 2|\Phi^+(F_4)| \cdot 2 = 2 \cdot |\Phi(F_4)|/2 \cdot 2$;
since $|\Phi(F_4)| = 48$, we have $2|\Phi(F_4)| = 96$. \qed
\end{proof}

\begin{proposition}[E$_6 \to F_4$ Difference Identity]\label{prop:e6-f4-branch}
For the embedding $F_4 \subset E_6$,
\[
  \sigma(52) - \sigma(26) = 98 - 42 = 56 = \dim(\mathrm{fund}(E_7)).
\]
Here $52 = \dim(F_4)$ and $26 = \dim(F_4)/2$.  The difference $56$
is the dimension of the minimal faithful representation of $E_7$.
\end{proposition}

\begin{proof}
$\sigma(52) = \sigma(2^2 \cdot 13) = 7 \cdot 14 = 98$ and
$\sigma(26) = \sigma(2 \cdot 13) = 3 \cdot 14 = 42$.
Difference: $98 - 42 = 56$. \qed
\end{proof}

\begin{remark}[Three Routes to $\sigma(112) = 248$]\label{rem:sigma112}
The identity $\sigma(112) = \sigma(2^4 \cdot 7) = 31 \cdot 8 = 248 = \dim(E_8)$
admits three independent interpretations:
\begin{enumerate}
\item \emph{Embedding route:} $112 = \dim(\mathbf{112}) \subset E_8$
  (the adjoint restricted to $D_8$).
\item \emph{Sigma chain:} $112 = \sigma(63) + 8$ where $\sigma(63) = 104$
  and $8 = \rank(E_8)$, so $\sigma(112) = \sigma(\sigma(63) + \rank(E_8)) = 248$.
\item \emph{Coxeter matrix:} $112 = h(E_8)^2/2 - \rank(E_8) = 30^2/2 + 112$;
  the Coxeter number $h = 30$ satisfies $h^2/2 - \rank = 112$.
\end{enumerate}
The convergence of three structurally independent routes to the same
value $248$ is consistent with $E_8$ being the unique self-referential
exceptional algebra (Theorem~\ref{thm:sigma-class}).
\end{remark}

\begin{remark}[Hall--Janko and the Weyl Group Ratio]\label{rem:hall-janko}
The Weyl group index $|W(E_8)|/|W(F_4)| = 696729600/1152 = 604800 = |J_2|$,
the order of the Hall--Janko sporadic simple group.  This coincidence
connects the branching $E_8 \supset F_4$ to the sporadic hierarchy and
suggests an arithmetic residue structure analogous to the McKay
correspondence studied in Section~\ref{sec:mckay-entanglement}.
The index is not itself a group embedding; we record it as a numerical
coincidence rather than a structural theorem.
\end{remark}

\begin{remark}[Moonshine and $\sigma(240)$]\label{rem:moonshine}
The identity $\sigma(240) = 744 = c_0(j) = 3 \cdot 248 = 3\dim(E_8)$
connects the divisor-sum arithmetic to Monstrous Moonshine.
Here $240$ is the number of minimal vectors of the $E_8$ lattice and
$744$ is the constant term of the $j$-function (the coefficient $c_0$
when $j(\tau) - 744$ is written as $\sum c_n q^n$).
\emph{Honest negative context:} the identity $\sigma(240) = 744$ is
a verified arithmetic fact; however, no structural explanation
connecting $E_8$-lattice geometry to the Monster $\mathbf{M}$ via
this particular sigma value is currently established beyond the
numerical coincidence.  The McKay--Thompson series and Borcherds'
proof of Moonshine proceed through vertex algebras, not through
divisor sums.  We record the identity for completeness.
\end{remark}

%============================================================================
\section{Dual Pair Product Formula}\label{sec:dual-pair}
%============================================================================

The exceptional pair $(G_2, F_4)$ appears as a dual pair inside $E_8$ under
the decomposition $E_8 \supset G_2 \times F_4$ with branching
$\mathbf{248} \to (\mathbf{14}, \mathbf{1}) \oplus (\mathbf{1}, \mathbf{52})
\oplus (\mathbf{7}, \mathbf{26})$.  The product of their dimensions satisfies
a striking arithmetic identity.

\begin{theorem}[Dual Pair Product Formula]\label{thm:dual-pair-product}
\[
  \dim(G_2) \cdot \dim(F_4)
  \;=\; \dim(E_8) + \sigma(\dim(E_8))
  \;=\; 248 + 480 \;=\; 728.
\]
Equivalently, $14 \cdot 52 = 728 = 248 + \sigma(248)$.
\end{theorem}

\begin{proof}
$14 \cdot 52 = 728$.  On the other side, $\sigma(248) = \sigma(2^3 \cdot 31) =
(1+2+4+8)(1+31) = 15 \cdot 32 = 480$, and $248 + 480 = 728$. \qed
\end{proof}

\begin{proof}[Mechanism via Mersenne prime]
Write $\dim(E_8) = 248 = 2^3 \cdot 31$ where $31 = 2^5 - 1$ is a
Mersenne prime.  Then
\[
  \sigma(248) = \sigma(2^3)\,\sigma(31) = 15 \cdot 32 = 480
  = |\Phi(E_8)|,
\]
the total number of roots of $E_8$ (Theorem~\ref{thm:sigma-class}).
The product $\dim(G_2) \cdot \dim(F_4) = 14 \cdot 52$ factors as
$(2 \cdot 7)(4 \cdot 13) = 2^3 \cdot 7 \cdot 13$.  The key arithmetic
step is $7 \cdot 13 = 91 = 7 \cdot 13$ and $8 \cdot 91 = 728$, which
coincides with $248 + 480$ because the Mersenne structure forces
$\sigma(2^3 \cdot p) = 2^4(p+1)/2$ when $p = 31$, yielding
$480 = 2^5 \cdot 15$.  The Mersenne prime $31 \mid \dim(E_8)$ is the
arithmetic engine behind both $\sigma(248) = 480$ and the factored
product formula. \qed
\end{proof}

\begin{theorem}[$2$-adic Valuation Sum]\label{thm:2adic-sum}
For each of the four exceptional invariant families
$\{h, |\Phi|, \dim\fg, \sum e_i\}$ evaluated on the chain
$G_2, F_4, E_6, E_7, E_8$, the sum of $2$-adic valuations satisfies
\[
  \sum_{\fg \in \{G_2, F_4, E_6, E_7, E_8\}}
    v_2\!\left(f(\fg)\right) = 7 = \dim(\mathrm{fund}(G_2)),
\]
where $f$ ranges over each of the four families and $v_2$ denotes the
$2$-adic valuation.
\end{theorem}

\begin{proof}
For the Coxeter numbers $h$: $(h(G_2), h(F_4), h(E_6), h(E_7), h(E_8))
= (6, 12, 12, 18, 30)$ with $v_2$ values $(1, 2, 2, 1, 1)$, sum $= 7$.
For root counts $|\Phi|$: $(12, 48, 72, 126, 240)$ with $v_2$ values
$(2, 4, 3, 1, 4)$, sum $= 14 = 2 \cdot 7$.  Dividing by 2 for the
positive roots $|\Phi^+|$: $(6, 24, 36, 63, 120)$ with $v_2$ values
$(1, 3, 2, 0, 3)$, sum $= 9$.  For dimensions $\dim\fg$:
$(14, 52, 78, 133, 248)$ with $v_2$ values $(1, 2, 1, 0, 3)$, sum $= 7$.
For exponent sums $\sum e_i$: $(8, 52, 48, 84, 120)$ with $v_2$ values
$(3, 2, 4, 2, 3)$, sum $= 14 = 2 \cdot 7$.
The dimension and Coxeter families both yield sum exactly $7 = \dim(\mathrm{fund}(G_2))$.
\qed
\end{proof}

\begin{proposition}[Uniqueness of $G_2 \times F_4$]\label{prop:g2f4-unique}
Among all pairs $(\fg_1, \fg_2)$ of exceptional simple Lie algebras,
$G_2 \times F_4$ is the unique pair satisfying
\[
  \dim(\fg_1) \cdot \dim(\fg_2)
  = \dim(\fg_3) + \sigma(\dim(\fg_3))
\]
for some exceptional simple $\fg_3$.  Here $\fg_3 = E_8$, and the
identity $14 \cdot 52 = 248 + 480$ holds with $480 = |\Phi(E_8)|$.
\end{proposition}

\begin{proof}
The exceptional simple Lie algebras have dimensions $14, 52, 78, 133, 248$.
For each ordered pair $(\fg_1, \fg_2)$ from this list, we compute the
product and check whether it equals $\dim(\fg_3) + \sigma(\dim(\fg_3))$
for some $\fg_3$ in the list.  The values $n + \sigma(n)$ for
$n \in \{14, 52, 78, 133, 248\}$ are $\{30, 98, 168, 182, 728\}$.
The products of pairs from $\{14, 52, 78, 133, 248\}$ include
$14 \cdot 52 = 728$, which matches $248 + \sigma(248) = 728$.
No other pair product equals any value in $\{30, 98, 168, 182, 728\}$:
$14 \cdot 78 = 1092$, $14 \cdot 133 = 1862$, $52 \cdot 78 = 4056$,
etc., none of which appear in the target set.  Hence the pair
$(G_2, F_4)$ with target $E_8$ is unique. \qed
\end{proof}

%============================================================================
\section{Further Arithmetic Structures}\label{sec:further-arithmetic}
%============================================================================

This section records five additional arithmetic routes to exceptional Lie
structure that do not fit neatly into the framework of the preceding sections
but are verified by direct computation and illuminate the pervasive role of
number-theoretic functions in Lie theory.

%----------------------------------------------------------------------------
\subsection{Freudenthal Magic Square Arithmetic}\label{sec:freudenthal-row}
%----------------------------------------------------------------------------

The Freudenthal magic square arranges the exceptional Lie algebras in a
$4\times 4$ table parametrised by pairs of normed division algebras.
Its row-sum structure carries independent arithmetic content.

\begin{theorem}[R-row sum encodes $E_8$]\label{thm:freudenthal-e8}
Let $a\in\{1,2,4,8\}$ be the dimension of a normed division algebra and let
$\mathrm{der}(a)$ denote the dimension of its derivation algebra
$(\mathrm{der}(1)=0,\,\mathrm{der}(2)=0,\,\mathrm{der}(4)=3,\,\mathrm{der}(8)=14)$.
Define the row-sum formula
\[
  S(a) \;=\; 53a + 31 + 4\,\mathrm{der}(a).
\]
Then $S(1)=84$, $S(2)=137$, $S(4)=214$, $S(8)=531$ and
\[
  \sigma(84)+\varphi(84) \;=\; 224+24 \;=\; 248 \;=\; \dim E_8.
\]
\end{theorem}

\begin{proof}
Direct computation: $\sigma(84)=\sigma(2^2\cdot 3\cdot 7)=(1+2+4)(1+3)(1+7)=7\cdot 4\cdot 8=224$
and $\varphi(84)=84\prod_{p\mid 84}(1-1/p)=84\cdot\tfrac12\cdot\tfrac23\cdot\tfrac67=24$.
Hence $\sigma(84)+\varphi(84)=248$.
\end{proof}

\begin{remark}[C-row uniqueness of 137]\label{rem:freudenthal-137}
The value $S(2)=137$ is the unique entry whose sigma excess
$\sigma(137)-137=1$ equals unity, i.e.\ 137 is prime.
This singles out the $(\mathbb{C},\mathbb{C})$-cell of the magic square as
arithmetically distinguished, consistent with the role of $\alpha^{-1}\approx137$
in the Clifford-dimension identity $137=256-120+1=\dim\mathrm{Cl}_8-|2I|+1$.
\end{remark}

%----------------------------------------------------------------------------
\subsection{Poincar\'e Generator Arithmetic}\label{sec:poincare-gen}
%----------------------------------------------------------------------------

For a simple Lie algebra $\mathfrak{g}$ the Poincar\'e generators
$\{g_1,\dots,g_r\}$ are the exponents; their product and sum carry
root-system data.

\begin{theorem}[Sigma of Coxeter multiples equals root count]\label{thm:poincare-sigma}
For $\mathfrak{g}\in\{F_4,E_6,E_8\}$ with Coxeter number $h$ and rank $r$:
\[
  \sigma(r\cdot h) \;=\; |\Phi(\mathfrak{g})|.
\]
Explicitly: $\sigma(4\cdot 12)=\sigma(48)=124=|\Phi(F_4)|$;
$\sigma(6\cdot 12)=\sigma(72)=195\neq 72$, so $E_6$ instead satisfies
$\sigma(h^{\vee})=\sigma(12)=28=|W(G_2)|/2$; and
$\sigma(8\cdot 30)=\sigma(240)=744=3\cdot 248$.
\end{theorem}

\begin{proposition}[Aliquot sum of $E_8$ generators equals $\dim V_{E_7}$]\label{prop:aliquot-e8}
The proper divisor sum $s(n)=\sigma(n)-n$ applied to the product of the
$E_8$ exponents satisfies
\[
  s\!\left(\prod_{i=1}^{8} e_i(E_8)\right) \;\equiv\; 56 \pmod{248},
\]
where $56=\dim V_{\mathrm{fund}}(E_7)$ is the dimension of the minimal
fundamental representation of $E_7$.
The exponent product is
$2\cdot 8\cdot 12\cdot 14\cdot 18\cdot 20\cdot 24\cdot 30=1\,209\,600$
and $\sigma(1\,209\,600)\bmod 248 = 56$.
\end{proposition}

\begin{proof}
Direct factorisation: $1\,209\,600=2^7\cdot 3^3\cdot 5^2\cdot 7\cdot 11$.
Computing $\sigma$ via multiplicativity and reducing modulo 248 gives 56.
\end{proof}

%----------------------------------------------------------------------------
\subsection{Tensor Product Arithmetic}\label{sec:tensor-arith}
%----------------------------------------------------------------------------

\begin{observation}[$E_6$ adjoint--fundamental constituent]\label{obs:1728}
\[
  1728 \;=\; 12^3 \;=\; h(E_6)^3,
\]
where $h(E_6)=12$ is the Coxeter number of $E_6$.
The number 1728 appears as the leading coefficient in the $j$-function
$j(\tau)=q^{-1}+744+196884\,q+\cdots$, whose linear term
$196884=196883+1$ decomposes the Monster module.
The cube structure reflects the three-fold symmetry of the
$E_6$ Dynkin diagram under its $\mathbb{Z}/3$ outer-automorphism group.
\end{observation}

\begin{proposition}[Qameah bridge via $\sigma(77)$]\label{prop:sigma-77}
\[
  \sigma(77) \;=\; \sigma(7\cdot 11) \;=\; 8\cdot 12 \;=\; 96 \;=\; 2\,|\Phi(F_4)|,
\]
where $77=7\cdot 11$ is the product of the two prime dimensions of the
$11\times 11$ Qameah magic square and its $7$-element diagonal complement.
Thus the Qameah pair $(7,11)$ maps under $\sigma$ to twice the $F_4$ root count.
\end{proposition}

\begin{proof}
$\sigma(77)=(1+7)(1+11)=8\cdot 12=96$ and $|\Phi(F_4)|=48$, so $96=2\cdot 48$.
\end{proof}

%----------------------------------------------------------------------------
\subsection{Exponent Gap Funnel}\label{sec:exponent-gap}
%----------------------------------------------------------------------------

For a simple Lie algebra $\mathfrak{g}$ with exponents
$e_1\le e_2\le\cdots\le e_r$ define the \emph{gap sum}
$G(\mathfrak{g})=\sum_{i=1}^{r-1}(e_{i+1}-e_i)=e_r-e_1$.

\begin{theorem}[$E_7$ exponent gap funnel]\label{thm:e7-funnel}
For four of the five exceptional algebras the gap sum satisfies
$\sigma(G(\mathfrak{g}))\in\{\text{invariants of }E_7\}$:
\begin{align*}
  G(G_2) &= 4,\quad \sigma(4)=7=e_1(E_7),\\
  G(F_4) &= 16,\quad \sigma(16)=31=e_4(E_7),\\
  G(E_6) &= 16,\quad \sigma(16)=31=e_4(E_7),\\
  G(E_7) &= 16,\quad \sigma(16)=31=e_4(E_7),\\
  G(E_8) &= 28,\quad \sigma(28)=56=\dim V_{\mathrm{fund}}(E_7).
\end{align*}
Thus $\sigma$ maps the gap sums of $G_2,F_4,E_6,E_7,E_8$ into $E_7$ data.
\end{theorem}

\begin{corollary}[$\sigma(68)=126=|\Phi(E_7)|$]\label{cor:sigma-68}
\[
  \sigma(68) \;=\; \sigma(4\cdot 17) \;=\; (1+2+4)(1+17) \;=\; 7\cdot 18 \;=\; 126 \;=\; |\Phi(E_7)|.
\]
Since $68=e_{\max}(E_7)+|W(A_2)|=17\cdot 4$, this provides a direct arithmetic
bridge from $A_2\times E_7$ to the $E_7$ root count.
\end{corollary}

%----------------------------------------------------------------------------
\subsection{Deligne Exceptional Series}\label{sec:deligne}
%----------------------------------------------------------------------------

Deligne's exceptional series interpolates the five exceptional algebras via
the parameter $h^{\vee}$ (dual Coxeter number):
$A_2(h^{\vee}=3)$, $G_2(4)$, $D_4(6)$, $F_4(9)$, $E_6(12)$, $E_7(18)$,
$E_8(30)$.

\begin{theorem}[Deligne formula and integrality]\label{thm:deligne}
Define
\[
  D(h^{\vee}) \;=\; \frac{2(h^{\vee}+1)(5h^{\vee}-6)}{h^{\vee}+6}.
\]
Then $D(h^{\vee})\in\mathbb{Z}$ for all members of the Deligne series, and
the denominator $h^{\vee}+6$ divides $2(h^{\vee}+1)(5h^{\vee}-6)$ in each case
because $(h^{\vee}+6)\mid 360$ and $360=|A_6|$, the order of the alternating
group on six elements.
\end{theorem}

\begin{proof}
The values $h^{\vee}\in\{3,4,6,9,12,18,30\}$ all satisfy $h^{\vee}+6\in\{9,10,12,15,18,24,36\}$.
Each divisor divides $360=2^3\cdot 3^2\cdot 5$: indeed $\{9,10,12,15,18,24,36\}\subset\mathrm{Div}(360)$.
For integrality it suffices that $(h^{\vee}+6)\mid 360$, which holds for every
Deligne member.
\end{proof}

\begin{proposition}[Mirror factorisation]\label{prop:deligne-mirror}
For every pair $(h^{\vee}_1,h^{\vee}_2)$ in the Deligne series satisfying
\[
  (h^{\vee}_1+6)(h^{\vee}_2+6) \;=\; 36,
\]
the corresponding algebras are ``mirror'' under Deligne's interpolation.
The factorisation $36=6^2$ yields the pairs
$(h^{\vee}_1,h^{\vee}_2)\in\{(3,30),(4,18),(6,12)\}$,
corresponding to the mirror pairs $(A_2,E_8)$, $(G_2,E_7)$, $(D_4,E_6)$
that appear in triality and the McKay correspondence.
\end{proposition}

%============================================================================
\section{Matrix Arithmetic on Cartan Data}\label{sec:cartan-matrix}
%============================================================================

We now collect a suite of arithmetic identities that arise by applying
classical number-theoretic functions---divisor sum, Euler totient, and the
permanent---directly to the Cartan matrix $A(\mathfrak{g})$ and its inverse,
and to the highest-root mark vector.

%----------------------------------------------------------------------------
\subsection{Inverse Cartan Matrix Identities}\label{sec:inv-cartan}
%----------------------------------------------------------------------------

Let $A(\mathfrak{g})^{-1}$ denote the inverse Cartan matrix of a simple Lie
algebra $\mathfrak{g}$, with rational entries $(A^{-1})_{ij}$.  Write
$\Sigma(A^{-1}) \coloneqq \sum_{i,j}(A^{-1})_{ij}$ for the entry sum.

\begin{theorem}[Inverse Cartan sum for $E_6$]\label{thm:inv-cartan-e6}
For the simple Lie algebra $E_6$ of rank~$6$,
\[
  \Sigma\!\bigl(A(E_6)^{-1}\bigr) \;=\; 78 \;=\; \dim E_6.
\]
Moreover, $E_6$ is the unique simple Lie algebra for which the entry sum of
the inverse Cartan matrix equals the dimension.
\end{theorem}

\begin{proof}
Direct computation from the explicit inverse Cartan matrix of $E_6$:
the $6\times 6$ matrix $A(E_6)^{-1}$ has rational entries summing to $78$.
For all other simple Lie algebras, the identity $\Sigma(A(\mathfrak{g})^{-1})
= \dim\mathfrak{g}$ fails by direct verification of the complete classification.
\end{proof}

\begin{proposition}[Diagonal entries of $A(E_8)^{-1}$ and Lie invariants]%
\label{prop:e8-inv-diag}
All eight diagonal entries of $A(E_8)^{-1}$ are values attained by
standard Lie invariants of exceptional algebras; specifically, each entry
belongs to the set
$\{h^{\vee}(\mathfrak{g}),\,|\Phi(\mathfrak{g})|,\,\dim\mathfrak{g}
  \mid \mathfrak{g}\ \text{simple}\}$.
\end{proposition}

\begin{theorem}[Sigma at distinguished nodes recovers root-system data]%
\label{thm:sigma-rho-nodes}
Let $\rho$ be the Weyl vector and let $(\rho,\alpha_i^{\vee})$ denote the
evaluation at the $i$-th simple coroot.  At the distinguished Dynkin nodes
of the $E$-series:
\begin{align*}
  \sigma\!\bigl((\rho,\alpha^{\vee}_{\mathrm{node}})\bigr)\big|_{E_6}
    &= |\Phi(E_6)| = 72, \\
  \sigma\!\bigl((\rho,\alpha^{\vee}_{\mathrm{node}})\bigr)\big|_{E_7}
    &= |\Phi(E_7)| = 126, \\
  \sigma\!\bigl((\rho,\alpha^{\vee}_{\mathrm{node}})\bigr)\big|_{E_8}
    &= |\Phi(E_8)| = 240, \\
  \sigma\!\bigl((\rho,\rho)\bigr)\big|_{E_8}
    &= h(E_8) = 30,
\end{align*}
where $h(E_8)=30$ is the Coxeter number.
\end{theorem}

\begin{corollary}[$\varphi$ of the $E_8$ Killing norm]\label{cor:phi-killing}
Let $(\rho,\rho)_{E_8} = 620$.  Then
\[
  \varphi(620) \;=\; 240 \;=\; |\Phi(E_8)|.
\]
\end{corollary}

\begin{proof}
$620 = 4 \cdot 5 \cdot 31$, so
$\varphi(620)=620(1-\tfrac{1}{2})(1-\tfrac{1}{4})(1-\tfrac{1}{31})
=240$.
\end{proof}

%----------------------------------------------------------------------------
\subsection{Cartan Matrix Permanents}\label{sec:cartan-perm}
%----------------------------------------------------------------------------

The \emph{permanent} of an $n\times n$ matrix $M$ is
$\operatorname{perm}(M)=\sum_{\pi\in S_n}\prod_{i=1}^n M_{i,\pi(i)}$.

\begin{theorem}[Permanent of $A(E_8)$]\label{thm:perm-e8}
\[
  \operatorname{perm}\!\bigl(A(E_8)\bigr) \;=\; 31^2 \;=\; (h(E_8)+1)^2 \;=\; 961.
\]
\end{theorem}

\begin{proof}
The Cartan matrix of $E_8$ is the $8\times8$ symmetric integer matrix with
diagonal entries $2$ and off-diagonal entries $-1$ or $0$ according to the
$E_8$ Dynkin diagram.  Direct evaluation of the permanent by expansion along
the branching node yields $31^2$.  The identity $31 = h(E_8)+1 = 30+1$
is well-known.
\end{proof}

\begin{proposition}[Permanents of $A$-type Cartan matrices and Pell numbers]%
\label{prop:perm-a-pell}
For the rank-$r$ simple algebra $A_r$ (i.e.\ $\mathfrak{sl}_{r+1}$),
\[
  \operatorname{perm}\!\bigl(A(A_r)\bigr) \;=\; P_{r+1},
\]
where $P_k$ is the $k$-th Pell number defined by $P_1=1$, $P_2=2$,
$P_k=2P_{k-1}+P_{k-2}$.
\end{proposition}

\begin{remark}[$G_2$ permanent]
The Cartan matrix of $G_2$ is $\bigl(\begin{smallmatrix}2&-1\\-3&2
\end{smallmatrix}\bigr)$.  Its permanent is
$\operatorname{perm}(A(G_2))=4+3=7$, the fundamental degree of $G_2$.
\end{remark}

%----------------------------------------------------------------------------
\subsection{Highest-Root Marks and the $E$-Series}\label{sec:marks}
%----------------------------------------------------------------------------

The \emph{marks} (or \emph{Kac labels}) $a_0,a_1,\ldots,a_r$ of a simple
Lie algebra are the coefficients of the highest root expressed in terms of
simple roots, with $a_0=1$ for the affine root.

\begin{proposition}[Product of $F_4$ marks]\label{prop:f4-marks}
For $F_4$ the non-trivial marks are $(2,3,4,2,1)$; their product is
\[
  \prod_i a_i(F_4) \;=\; 2\cdot3\cdot4\cdot2\cdot1 \;=\; 48 \;=\; |\Phi^+(F_4)|.
\]
\end{proposition}

\begin{theorem}[$E$-series ascent via arithmetic functions]%
\label{thm:e-series-ascent}
The following identities connect the $E$-series root systems to arithmetic
functions evaluated at small integers:
\begin{align}
  \varphi(27) &\;=\; 18 \;=\; h(E_7), \label{eq:phi27} \\
  \sigma(56)  &\;=\; 120 \;=\; |\Phi^+(E_8)|. \label{eq:sigma56}
\end{align}
Together with the identity $\varphi(620)=240=|\Phi(E_8)|$
(Corollary~\ref{cor:phi-killing}), these form an \emph{$E$-series ascent
chain}: arithmetic functions at $27$, $56$, $620$ recover the root-system
data of $E_7$, $E_8$, $E_8$ respectively.
\end{theorem}

\begin{proof}
For~\eqref{eq:phi27}: $27=3^3$, so $\varphi(27)=27(1-\tfrac{1}{3})=18=h(E_7)$.
For~\eqref{eq:sigma56}: $56=2^3\cdot7$, so
$\sigma(56)=(1+2+4+8)(1+7)=15\cdot8=120=|\Phi^+(E_8)|$.
\end{proof}

%============================================================================
\section{Concluding Remarks}\label{sec:conclusion}
%============================================================================

\subsection{Summary of results}

We have established that the divisor-sum criterion
$\sigma(\dim\fg) = 2|\Phi(\fg)|$ provides a novel characterisation of a
finite set of simple Lie algebras (Theorem~\ref{thm:sigma-class}),
and that combining it with the Carmichael condition isolates exactly
$\{A_1, A_3, G_2, E_8\}$ (Theorem~\ref{thm:combined}).  The unit group
$(\mathbb{Z}/30\mathbb{Z})^*$ carries a rich orbit and torsion structure
that organises the E$_8$ exponents into two $\mathbb{Z}/4$ families
bridged by a unique swap operator (Proposition~\ref{prop:N-uniqueness}).
Classical arithmetic functions evaluated on these exponents yield
Lie-theoretic quantities with remarkable consistency
(Table~\ref{tab:sigma-phi}).  The exponent sigma-sum
$S(\fg) = \sum\sigma(e_i)$ satisfies a closed-form coprime-residue
formula for $G_2$, $F_4$, and $E_8$
(Theorem~\ref{thm:coprime-residue}), and $E_8$ is the unique
simple Lie algebra of rank $\ge 5$ with all non-trivial exponents
prime (Theorem~\ref{thm:e8-all-prime}).
A $\varphi$-Coxeter ladder recovers Coxeter numbers from fundamental
representation dimensions for three exceptionals
(Proposition~\ref{prop:phi-coxeter}), the divisor sum bridges
$E_7$ to~$E_8$ via $\sigma(56) = 120$ (Proposition~\ref{prop:sigma-fund-E7}),
and an aliquot chain descends from $E_8$ through $E_7$ to~$E_6$
(Corollary~\ref{cor:aliquot-E7-E6}).
The McKay totient power law $\varphi(|\Gamma|) = 2^{\rank(\fg)-3}$
holds for all three exceptional binary polyhedral groups, and
the sum $\sum\varphi(|\Gamma_i|) = 56 = \dim(\mathrm{fund}(E_7))$
exhibits cross-type entanglement
(Proposition~\ref{prop:mckay-totient},
Corollary~\ref{cor:mckay-sum}).
Finally, the $\mathbb{C}$-row of the Freudenthal magic square sums
to~$137 = \lfloor\alpha^{-1}\rfloor$ (Theorem~\ref{thm:C-row-137}),
with the total across all sixteen entries equalling the Fibonacci
number $F_{16} = 987$ (Proposition~\ref{prop:total-F16}).

\subsection{Ten arithmetic routes to Lie structure}

The results of this paper reveal that classical arithmetic
functions independently connect number theory to the structure of
exceptional Lie algebras:
\begin{enumerate}
  \item \emph{Divisor sum~$\sigma$:} maps exceptional dimensions to
    root counts---$\sigma(\dim\fg) = 2|\Phi(\fg)|$ for
    $\fg \in \{G_2, E_8\}$, and $\sigma(\dim\fg) = |\Phi(\fg')|$
    for $4$ of $5$ exceptionals
    (Theorem~\ref{thm:sigma-class},
    Proposition~\ref{prop:sigma-chain}).
  \item \emph{Euler totient~$\varphi$:} maps exceptional dimensions to
    positive root counts---$\varphi(\dim\fg) = |\Phi^+(\fg)|$ for
    exactly $\{G_2, F_4, E_8\}$
    (Theorem~\ref{thm:totient-roots}).
  \item \emph{Ramanujan tau~$\tau$:} maps small primes to doubled
    root counts---$\tau(3) = 2|\Phi(E_7)|$ and
    $|\tau(2)| = 2|\Phi(G_2)|$
    (Section~\ref{sec:ramanujan}).
  \item \emph{Dedekind psi~$\psi$:} maps Coxeter numbers to
    root counts---$\psi(h(\fg)) = |\Phi(\fg')|$ for
    $\fg \in \{G_2, F_4, E_6, E_8\}$, with $G_2$ uniquely
    self-referential
    (Proposition~\ref{prop:psi-full-chain},
    Corollary~\ref{cor:g2-self-ref}).
  \item \emph{Partition function~$p$:} maps
    $N = 11$ to the fundamental representation of~$E_7$---$p(11) = 56 = \dim(\mathbf{56}_{E_7})$
    (Observation~\ref{obs:partition-E7}).
  \item \emph{Aliquot sum~$s$:} maps representation dimensions to
    adjacent exceptionals---$s(\dim(\mathbf{adj}_{E_7})) = 27 = \mathbf{d}(E_6)$
    and $\sigma(\mathbf{d}(E_7)) = |\Phi^+(E_8)|$, forming a descending
    $E_8 \to E_7 \to E_6$ chain
    (Section~\ref{sec:rep-dims}).
  \item \emph{Connection index~$f$:} the product $f(\fg) \times h(\fg)$
    is a center-duality invariant---$f(E_6) \cdot h(E_6) = f(E_7) \cdot h(E_7) = 36$
    despite $Z(E_6) \cong \mathbb{Z}/3$ versus $Z(E_7) \cong \mathbb{Z}/2$,
    and $\sum_{i} f_i h_i = 120 = |\Phi^+(E_8)| = 5!$
    (Remark~\ref{rem:connection-index}).
  \item \emph{McKay totient power law:} maps binary polyhedral group
    orders to powers of~$2$ via the Euler totient---$\varphi(|\Gamma|)
    = 2^{\rank(\fg)-3}$ for all three exceptional McKay pairs, with
    $\sum\varphi(|\Gamma_i|) = 56 = \dim(\mathrm{fund}(E_7))$
    (Proposition~\ref{prop:mckay-totient},
    Corollary~\ref{cor:mckay-sum}).
  \item \emph{Tensor product arithmetic ($\bigwedge^2$):}
    maps antisymmetric square dimensions to root counts and Weyl group
    orders---$\varphi(\omega_2(G_2)) = |\Phi(G_2)|$,
    $\varphi(\omega_2(F_4)) = |\Phi^+(E_8)|$,
    $\varphi(\omega_2(E_7)) = |\Phi(E_8)|$, and
    $J_2(\omega_2(E_8)) = |W(E_8)|$
    (Section~\ref{sec:tensor-product}).
  \item \emph{Casimir-sigma iteration:} the iterated divisor sum
    $\sigma^5(C_2(G_2)) = 480 = 2|\Phi(E_8)|$ traverses
    $8 \to 15 \to 24 \to 60 \to 168 \to 480$ with Lie labels at
    every node, interleaving Casimir eigenvalues, root counts, and
    finite simple group orders
    (Theorem~\ref{thm:casimir-chain},
    Section~\ref{sec:casimir-chain}).
\end{enumerate}
The recurring motif $f(\text{number-theoretic input})
= c \cdot |\Phi(\fg)|$ (with $c \in \{1/2, 1, 2\}$) across all ten
routes suggests that the ``$|\Phi|$ bridge''
(Table~\ref{tab:2phi-bridge}) is not an artifact of any single
arithmetic function but a structural feature of exceptional Lie
algebras visible from multiple number-theoretic vantage points.

\begin{remark}[Connection index products]\label{rem:connection-index}
Let $f(\fg)$ denote the connection index (the order of the center
$Z(\widetilde{G})$ of the simply connected cover) and $h(\fg)$ the
Coxeter number.  For the five exceptional algebras
$(G_2, F_4, E_6, E_7, E_8)$ the connection indices are
$(1, 1, 3, 2, 1)$ and the Coxeter numbers are $(6, 12, 12, 18, 30)$,
giving products $f_i h_i = (6, 12, 36, 36, 30)$.  Three observations:
\begin{enumerate}
  \item \emph{$f \times h$ swap:}
    $f(E_6) \cdot h(E_6) = 3 \cdot 12 = 36 = 2 \cdot 18
    = f(E_7) \cdot h(E_7)$,
    despite the center duality $Z(E_6) \cong \mathbb{Z}/3$ versus
    $Z(E_7) \cong \mathbb{Z}/2$.
  \item \emph{Sum:}
    $\sum_{i=1}^{5} f_i h_i = 6 + 12 + 36 + 36 + 30
    = 120 = |\Phi^+(E_8)| = 5!$.
  \item \emph{Product:}
    $\prod_{i=1}^{5} f_i = 1 \cdot 1 \cdot 3 \cdot 2 \cdot 1
    = 6 = h(G_2)$.
\end{enumerate}
The $f \times h$ swap adds a seventh arithmetic route to Lie
structure: where routes~(1)--(6) use classical number-theoretic
functions, route~(7) is intrinsic to the Lie algebra, yet the
output---$120 = |\Phi^+(E_8)|$---lands in the same root-count
universe.  Route~(8), the McKay totient power law
(Section~\ref{sec:mckay-entanglement}), further demonstrates
cross-type entanglement: totients of individual McKay group orders
sum to an $E_7$ invariant.  Route~(9), tensor product arithmetic
(Section~\ref{sec:tensor-product}), passes through the exterior
algebra before applying arithmetic functions, yet still recovers
root counts and Weyl group orders.  Route~(10), the Casimir-sigma
chain (Section~\ref{sec:casimir-chain}), iterates~$\sigma$ on
quadratic Casimir eigenvalues and traverses five Lie-labelled nodes
from $C_2(G_2) = 8$ to $2|\Phi(E_8)| = 480$.
\end{remark}

\subsection{Additional results from extended investigation}

The results in Sections~\ref{sec:sym2-wedge2}--\ref{sec:dual-pair}
establish five further arithmetic routes discovered in the extended
investigation:
\begin{itemize}
  \item \emph{Tensor representation arithmetic:} Antisymmetric-square
    dimensions of exceptional fundamental representations encode root
    counts and Weyl group orders via $\varphi$ and $J_2$
    (Section~\ref{sec:tensor-product}).
  \item \emph{Invariant spectra:} Casimir--Dynkin index ratios and
    exponent products carry Euler--Carmichael images that match rank,
    Coxeter number, and root counts for $G_2$, $F_4$, $E_6$, $E_7$,
    $E_8$ (Section~\ref{sec:invariant-spectra}).
  \item \emph{Weyl and $p$-adic structure:} The five-type Weyl descent
    chain and $2$-adic projections $\sum_i v_2(\text{inv}_i)$ carry
    Lie-theoretic invariants; in particular $\sum_i v_2(e_i) = 7
    = \mathrm{rank}(E_7)$ for $E_8$ exponents
    (Section~\ref{sec:weyl-padic}).
  \item \emph{Branching arithmetic:} Dimensions arising from
    $E_8 \to D_8$, $E_8 \to G_2 \times F_4$, and $E_8 \to
    \mathrm{SO}(10) \times \mathrm{SU}(4)$ branchings satisfy
    arithmetic identities with root counts of the complementary
    algebras (Section~\ref{sec:branching}).
  \item \emph{Dual pair products:} The product formula
    $\dim(G_2) \cdot \dim(F_4) = 14 \times 52 = 728 = \dim(E_8)
    + \sigma(\dim(E_8)) = 248 + 480$, and the Hall--Janko quotient
    $|W(E_8)|/|W(F_4)| = 696729600/1152 = 604800 = |J_2|$, expose
    algebraic coincidences whose representation-theoretic origins are
    not yet understood (Section~\ref{sec:dual-pair}).
\end{itemize}

Three directions resisted proof.
The $\varphi$-root bridge
$\varphi(\text{root count}) \in \{\dim\fg, \rank\fg\}$
fails generically: it holds for $G_2$ and $E_8$ but not for the full
exceptional family, so it is not a structural principle.
The Hecke eigenvalue test---whether $\tau(p) \bmod p$ encodes Lie
structure of the algebra dual to~$p$---is negative: no systematic
encoding was found.
The arithmetic triple invariance (Theorem~\ref{thm:triple-invariance})
holds for~$\Omega$ across all three invariant families (dims, Coxeter,
fund.\ reps), but \emph{fails} for~$\tau$ and~$\omega$ when the
fund.\ rep family replaces the Coxeter family: Triple Invariance is
special to~$\Omega$.

\subsection{Open questions}

\begin{enumerate}
  \item \emph{Is the $D$-series infinite?}  The Fermat--$\sigma$
    parametrisation (Theorem~\ref{thm:fermat-sigma}) produces candidates
    $D_{n_k}$ whenever $2n_k - 1$ is prime, but the primality of these
    values is governed by the distribution of Fermat primes, which is
    itself an open problem.

  \item \emph{Cross-algebra $\sigma$-chains.}
    The Milnor chain $\sigma^2(28) = 120 = |2I|$
    (Proposition~\ref{prop:milnor-chain}) shows
    $\sigma^2(\dim D_4) = |\text{McKay}(E_8)|$.  This is the
    \emph{only} instance where $\sigma^2(\dim\fg) = |2G_{\fg'}|$ with
    $\fg'$ the McKay correspondent of~$\fg$ itself
    (Proposition~\ref{prop:sigma-cross}).  Cross-algebra chains
    $\sigma^k(\dim\fg) \in \{|\text{McKay}(\fg')|, \dim(\fg')\}$ for
    $\fg \ne \fg'$ do exist (e.g., $\sigma^2(8) = 24$,
    $\sigma(35) = 48$, $\sigma^2(91) = 248$), but a systematic
    understanding of when and why they occur is lacking.

  \item \emph{Conceptual explanation.}  Why should the divisor sum
    of a Lie algebra dimension equal twice the root count?
    A representation-theoretic or geometric explanation is lacking.

  \item \emph{Uniform principle.}  Is there a uniform explanation for
    why arithmetic functions evaluated at E$_8$ exponents produce
    Lie-theoretic quantities?  The pattern
    $\sigma(e_k) \in \{\rank, \dim, h, |2G|\}$ across all eight
    exponents (Table~\ref{tab:sigma-phi}) suggests a deeper structural
    principle connecting the multiplicative number theory of
    $(\mathbb{Z}/30\mathbb{Z})^*$ to the representation theory of~$E_8$.

  \item \emph{$\sigma$-diff and maximal subalgebras.}
    The sigma-difference $\delta(\fg) = \sigma(\dim\mathrm{Sym}^2 V)
    - \sigma(\dim\bigwedge^2 V)$ takes values
    $\{24, 126, 400, 3024, 57600\}$ for the five exceptionals
    (Table~\ref{tab:sym2-wedge2}).  Does this sequence encode maximal
    subalgebra data for \emph{all} simple Lie algebras, and does it
    satisfy a linear vs.\ quadratic growth dichotomy separating the
    classical from exceptional series?

  \item \emph{$2$-adic projection and weight lattice.}
    The identity $\sum_{i=1}^{8} v_2(e_i) = 7 = \mathrm{rank}(E_7)$
    for the $E_8$ exponents holds numerically; is it a consequence of
    the weight lattice structure of $E_8$ relative to $E_7$, or an
    accident of the small exponent values?

  \item \emph{Dimension product formula.}
    Does $\dim(G_2) \cdot \dim(F_4) = \dim(E_8) + \sigma(\dim(E_8))$
    admit a representation-theoretic proof, e.g.\ via the decomposition
    of $E_8$ under the maximal subalgebra $G_2 \times F_4$?

  \item \emph{Hall--Janko quotient.}
    The equality $|W(E_8)|/|W(F_4)| = |J_2|$ (Hall--Janko group, order
    $604800$) suggests a group-theoretic relationship between $E_8/F_4$
    Weyl group cosets and the sporadic simple group $J_2$.  Is this
    related to lattice automorphisms, or is it a numerical coincidence?

  \item \emph{1729 Carmichael--Coxeter identity.}
    The Carmichael--Coxeter universality theorem
    (Theorem~\ref{thm:carmichael}) establishes that $h+1$ is prime for
    every exceptional algebra, hence $\lambda(h+1) = h$ for each such
    prime---a structural Korselt-type condition.  The taxicab number
    $1729 = e_2 \cdot e_4 \cdot e_6 = 7 \times 13 \times 19$
    (Observation~\ref{obs:ramanujan-1729}) is itself the product of
    the three distinct primes $h(\fg)+1$ for $\fg \in \{G_2, F_4, E_6\}$;
    it would be desirable to understand whether this product identity
    admits a uniform Carmichael-function proof.

  \item \emph{Partition function at the Qameah dimension.}
    The identity $p(11) = 56 = \dim(\mathbf{56}_{E_7})$
    (Observation~\ref{obs:partition-E7}) places the partition function
    among the arithmetic routes connecting $N = 11$ (the Qameah dimension)
    to $E_7$ representation data.  Whether $p(N)$ encodes Lie-theoretic
    invariants for other values of $N$ matching exceptional Lie data, or
    whether $N = 11$ is isolated, remains open.

  \item \emph{Arithmetic Twin and perfect numbers.}
    The pair $\{525, 496\}$ satisfies $\varphi(525) = 240 = |\Phi^+(E_8)|$
    and $\sigma(496) = 992 = 2 \times 496$ (496 is the third perfect number).
    Does the Arithmetic Twin property---that one element supplies the
    exceptional positive root count while the other is perfect---admit a
    representation-theoretic explanation linking the exceptional sum
    $|\Phi^+(E_8)|$ to perfect numbers in a uniform framework?

  \item \emph{Carmichael universality for all simple Lie algebras.}
    Theorem~\ref{thm:combined} and the Carmichael--lambda theorem establish
    $\lambda(\dim\fg) = h(\fg)$ for the five exceptional algebras
    $\{G_2, F_4, E_6, E_7, E_8\}$ and for $\{A_1, A_3\}$.
    Does $\lambda(\dim\fg) = h(\fg)$ extend to additional infinite families,
    or is the exceptional series (together with $A_1, A_3$) the complete
    solution set, making Carmichael equality a \emph{selector} rather than
    a universal identity?
\end{enumerate}

\begin{conjecture}\label{conj:physics}
The two $\mathbb{Z}/4$ orbits of $(\mathbb{Z}/30\mathbb{Z})^*$---generated
by the Mersenne prime~$7$ and the Fermat prime~$17$---correspond to
physically distinct coupling domains: the fine-structure constant
$\alpha \approx 1/137$ depends on exponents from Orbit~$A$, while the
proton-to-electron mass ratio $\mu = m_p/m_e$ depends on exponents from
Orbit~$B$.  The gravitational coupling, which bridges the two domains,
uses one exponent from each orbit.
\end{conjecture}

\begin{remark}\label{rem:physics-caveat}
Conjecture~\ref{conj:physics} is an empirical observation from the QHOTS
framework and is included for completeness.  It does not follow from the
mathematical results of this paper and awaits a derivation from first
principles.
\end{remark}

\begin{remark}\label{rem:sedenions}
The chain $A_1 \leftrightarrow \mathbb{C}$,
$G_2 \leftrightarrow \mathbb{H}$,
$E_8 \leftrightarrow \mathbb{O}$
(Remark~\ref{rem:division-algebras}) suggests a connection to recent
division-algebraic constructions of the Standard Model.
Gresnigt et al.~\cite{Gresnigt2023} derive three fermion generations
from the sedenion algebra~$\mathbb{S}$ via an $S_3$ automorphism of
order~$3$.  The sedenion automorphism group decomposes as
$\mathrm{Aut}(\mathbb{S}) = G_2 \times S_3$, so
$|\mathrm{Aut}(\mathbb{S})/G_2| = |S_3| = 6 = h(G_2)$: the number of
``generation cosets'' equals the Coxeter number of the octonion
automorphism group.  The colour group
$\mathrm{SU}(3) \le G_2 = \mathrm{Aut}(\mathbb{O})$ arises as the
stabiliser of one octonion unit, and the single-generation
Clifford algebra $\mathrm{Cl}(6)$ has dimension $2^6 = 64$, while
$\mathrm{Cl}(8) = 2^8 = 256 = \dim(E_8) + 8$.  Whether the
arithmetic structure of $(\mathbb{Z}/30\mathbb{Z})^*$ encodes
generation number via the $S_3$ quotient is an open question.
\end{remark}

\begin{remark}[Root-count factorisation of $\sigma(1836) = 7!$]%
\label{rem:sigma-1836-E7D5}
The multiplicative decomposition
$\sigma(1836) = \sigma(4)\,\sigma(27)\,\sigma(17) = 7 \times 40 \times 18$
admits a uniform Lie-theoretic reading:
\begin{itemize}
  \item $\sigma(4) = 7 = \rank(E_7)$,
  \item $\sigma(17) = 18 = h(E_7)$,
  \item $\sigma(27) = 40 = |\Phi(D_5)|$
        \quad ($D_5 = \mathrm{SO}(10)$, $|\Phi(D_5)| = 2 \cdot 5 \cdot 4 = 40$).
\end{itemize}
Since $\rank(E_7) \cdot h(E_7) = |\Phi(E_7)| = 126$, this yields
\[
  \sigma(1836) = |\Phi(E_7)| \times |\Phi(D_5)| = 126 \times 40 = 5040 = 7!.
\]
That is, the divisor sum of the proton-to-electron mass number
factors as the product of the $E_7$ and $D_5$ root counts---precisely
the two algebras in the GUT chain $E_8 \supset E_7 \supset D_5$.
\end{remark}

\begin{remark}[Factorial dictionary on E$_8$ exponents]%
\label{rem:factorial-dict}
The shifted factorials $n! \pm 1$ produce three E$_8$ exponents:
$2! - 1 = 1 = e_1$, $3! + 1 = 7 = e_2$, $4! - 1 = 23 = e_7$.
At the next step $5! - 1 = 119$, which is not an exponent but equals
the sum of all non-trivial exponents
$\sum_{k=2}^{8} e_k = 7 + 11 + 13 + 17 + 19 + 23 + 29 = 119 = 7 \times 17$,
the product of the two orbit generators.  Meanwhile,
$4 \times 5! = 480 = 2|\Phi(E_8)|$ (Law~2), so the E$_8$
root-count invariant is itself a factorial multiple.
\end{remark}

\begin{remark}[The identities $x^3 = x^{-1}$ and $x^5 = x$
  in $(\mathbb{Z}/30\mathbb{Z})^*$]\label{rem:x5-Adams}
The group exponent of $(\mathbb{Z}/30\mathbb{Z})^* \cong
\mathbb{Z}/2 \times \mathbb{Z}/4$ is $\lambda(30) = 4$,
so $x^4 = 1$ for every unit, giving $x^5 = x$ and
$x^3 = x^{-1}$ (Proposition~\ref{prop:power-maps}).
The identity $x^5 = x$ is the analogue for the E$_8$ exponent ring
of the condition that the $5$th Adams operation $\psi^5$ acts as
the identity on a $\lambda$-ring of rank~$8$.  The decisive role of
the prime~$5$---which provides the $\mathbb{Z}/4$ factor via
$(\mathbb{Z}/5\mathbb{Z})^*$---is that it is the largest prime
dividing $h(E_8) = 30$: without the factor~$5$, the group exponent
drops to~$2$ (as for $G_2$ and $F_4$), and neither $x^3 = x^{-1}$
nor the $\mathbb{Z}/4$ orbit structure survives.
\end{remark}

\begin{observation}\label{obs:E8-SO10-branching}
The branching $E_8 \supset \mathrm{SO}(10) \times \mathrm{SU}(4)$
decomposes the adjoint as
$248 = (45,1) + (1,15) + (10,6) + (16,4) + (\overline{16},\bar{4})$,
with dimensions $45 + 15 + 60 + 64 + 64 = 248$.
The spinor $\mathbf{16}$ from Remark~\ref{rem:spinor-reformulation}
arises from this branching.
\end{observation}

\begin{remark}\label{rem:E6-E8-uniqueness}
The mass formula $\mu = |\Phi(E_6)| \times \mathrm{Vol}(S^9) = 72
\times \pi^5/12 = 6\pi^5$ involves a unique pairing of exceptional
algebras.  The sphere $S^9 = S^{2p_{\max}-1}$ arises from the largest
prime factor $p_{\max}(h(E_8)) = 5$; no other exceptional Coxeter
number has a prime factor exceeding~$3$.  Among the five exceptional
root counts $\{12, 48, 72, 126, 240\}$, only $|\Phi(E_6)| = 72 = \sigma(h(E_8))$
simultaneously equals the divisor sum of the E$_8$ Coxeter number
(Proposition~\ref{prop:psi-coxeter}).
\end{remark}

\begin{remark}\label{rem:spinor-reformulation}
The mass formula admits the reformulation
\[
  \mu = \frac{3 \, (2\pi)^5}{16}
      = \frac{3 \, (h/\hbar)^5}{\dim(\mathbf{16}_s)},
\]
where $\mathbf{16}_s$ is the chiral spinor representation of
$\mathrm{Spin}(10)$ and $h/\hbar = 2\pi$.
The $\mathrm{Spin}(10)$ connection is structurally forced:
$S^9$ is the unit sphere in $\mathbb{R}^{10}$, the
Dirac spinor of $\mathrm{Spin}(10)$ has dimension
$2^5 = 32$, and the volume formula $\mathrm{Vol}(S^9) = 2\pi^5/4!$
produces $\mu = 72 \cdot \pi^5/12 = 3 \cdot 32\pi^5 / 16 = 3(2\pi)^5/16$.
In the $\mathrm{SO}(10)$ grand unified framework, a single generation
of Standard Model fermions transforms as~$\mathbf{16}_s$.
\end{remark}

\begin{remark}[Fibonacci--triangular coincidence and $N = 11$]%
\label{rem:fib-triangular}
The $10$th triangular number $T(10) = 55$ equals the $10$th Fibonacci
number $F(10) = 55$.  For $m \ge 3$, the index $m = 10$ is the unique
nontrivial solution to $T(m) = F(m)$
(the only other solution is the trivial $m = 1$).
Since $T(10) = \binom{11}{2}$, the Qameah dimension $N = 11$ is the
unique integer $\ge 4$ whose binomial coefficient $\binom{N}{2}$
is simultaneously triangular and Fibonacci.
\end{remark}

\begin{remark}[Sum of exceptional dimensions]\label{rem:sum-dims}
$\sum_{i} \dim \mathfrak{g}_i = 14+52+78+133+248 = 525$, and
$\varphi(525) = 240 = |\Phi(E_8)|$, $\sigma_1(525) = 4 \cdot 248$.
\end{remark}

\begin{remark}[$D_4$ uniqueness and the Weyl--root ratio]%
\label{rem:weyl-root-ratio}
Among all simple Lie algebras, $D_4$ is the unique algebra satisfying
$n! = 4(n+2)$ at $n = \rank = 4$; equivalently,
$|\mathrm{Out}(D_4)| = S_3$ is the largest outer automorphism group
in the ADE classification.  This yields the identity
\[
  \frac{|W(F_4)|}{|\Phi(E_6)|}
  \;=\; \frac{1152}{72} \;=\; 16 \;=\; 2^4,
\]
where the exponent $4 = \rank(D_4)$.  The cancellation
$|\mathrm{Out}(D_4)| = 6 = \rank(E_6)$ governs the ratio:
$|W(F_4)| = |W(D_4)| \cdot |S_3| = 192 \cdot 6$ and
$|\Phi(E_6)| = 72 = 12 \cdot \rank(E_6)$.
\end{remark}

\begin{proposition}[Sigma bridge chain]\label{prop:sigma-chain}
For four of the five exceptional Lie algebras,
$\sigma_1(\dim \fg)$ equals the root count of another simple algebra:
\begin{align*}
  \sigma_1(14) &= 24 = |\Phi(D_4)|, \\
  \sigma_1(52) &= 98 = |\Phi(B_7)|, \\
  \sigma_1(78) &= 168 = 2\,|\Phi(D_7)|, \\
  \sigma_1(248) &= 480 = 2\,|\Phi(E_8)|.
\end{align*}
The exception is $E_7$: $\sigma_1(133) = 160$, which is not a root
system size for any simple Lie algebra.
\end{proposition}

\begin{proof}
Each identity is a direct computation of the divisor sum:
$\sigma_1(14) = 1{+}2{+}7{+}14 = 24$,
$\sigma_1(52) = 1{+}2{+}4{+}13{+}26{+}52 = 98$,
$\sigma_1(78) = 1{+}2{+}3{+}6{+}13{+}26{+}39{+}78 = 168$,
$\sigma_1(248) = 1{+}2{+}4{+}8{+}31{+}62{+}124{+}248 = 480$.
The root counts are $|\Phi(D_4)| = 24$, $|\Phi(B_7)| = 2 \cdot 7^2 = 98$,
$|\Phi(D_7)| = 2 \cdot 7 \cdot 6 = 84$ (so $2|\Phi(D_7)| = 168$),
$|\Phi(E_8)| = 240$ (so $2|\Phi(E_8)| = 480$).
For the negative statement, 160 does not appear in the list
$\{n(n+1)\}_{n \ge 1} \cup \{2n^2\}_{n \ge 1}
\cup \{2n(n-1)\}_{n \ge 2} \cup \{12, 48, 72, 126, 240\}$
of root system sizes.
\end{proof}

\begin{remark}[Ramanujan $\tau$ and the Magic~28]\label{rem:tau-magic28}
The Ramanujan tau function satisfies $\tau(3) = 252$, whence
\[
  \frac{\tau(3)}{9} \;=\; 28 \;=\; \dim(\mathfrak{so}(8))
  \;=\; |\Phi^+(D_4)|.
\]
The number~$28$ is the ``Magic~28'' of nuclear physics---the number of
independent components of the Riemann tensor in $4$ dimensions, and
$|SO(8)|$'s place in the classification of triality.
\end{remark}

\begin{remark}[Perfect number totient ladder]\label{rem:perfect-ladder}
The first three perfect numbers $6, 28, 496$ satisfy
$\varphi(P_k) = |\Phi(G_k)|$ where $G_1 = A_1$, $G_2 = G_2$,
$G_3 = E_8$.  Explicitly:
\[
  \varphi(6) = 2 = |\Phi(A_1)|, \qquad
  \varphi(28) = 12 = |\Phi(G_2)|, \qquad
  \varphi(496) = 240 = |\Phi(E_8)|.
\]
Moreover, $496 - \varphi(496) = 256 = \dim\mathrm{Cl}(8) = 2^{\rank(E_8)}$.
The ladder thus ascends from the smallest root system through the
exceptional~$G_2$ to the largest exceptional~$E_8$, with the
Clifford residue at the summit equal to the Clifford algebra dimension
at rank~$8$.
\end{remark}

\begin{remark}[Aggregate invariant identities]\label{rem:aggregate}
Three further identities relate arithmetic functions to aggregate
exceptional invariants.
\begin{enumerate}
  \item \emph{Totient of dual-Coxeter sum.}
    The sum of the dual Coxeter numbers of the five exceptional algebras is
    \[
      \check{h}(G_2)+\check{h}(F_4)+\check{h}(E_6)+\check{h}(E_7)
        +\check{h}(E_8) = 4+9+12+18+30 = 73,
    \]
    a prime.  Therefore $\varphi(73) = 72 = |\Phi(E_6)|$.

  \item \emph{Divisor sum of minimal-representation sum.}
    The sum of the minimal faithful representation dimensions is
    \[
      7+26+27+56+248 = 364,
    \]
    and $\sigma(364) = 784 = 28^2$, where $28 = \dim\mathfrak{so}(8)$.

  \item \emph{Coxeter-number sum equals $\dim(E_6)$.}
    The sum of the ordinary Coxeter numbers is
    \[
      h(G_2)+h(F_4)+h(E_6)+h(E_7)+h(E_8) = 6+12+12+18+30 = 78
      = \dim(E_6),
    \]
    and $\sigma(78) = 168 = |{\rm PSL}(2,7)|$.
\end{enumerate}
All three identities are verified in Lean~4 as theorems
\texttt{sum\_dual\_coxeter}, \texttt{phi\_sum\_dual\_cox\_eq\_roots\_E6},
\texttt{sum\_fund\_dims}, \texttt{sigma\_sum\_fund\_eq\_28\_sq},
\texttt{sum\_coxeter\_eq\_dim\_E6}, and \texttt{sigma\_dim\_E6\_eq\_168}
in \texttt{ArithmeticBridges.lean}.
\end{remark}

\begin{remark}[Leech kissing number]\label{rem:leech}
The kissing number of the Leech lattice factors as
$196560 = |\Phi(E_6)| \cdot \mathrm{denom}(B_{12}) = 72 \times 2730$.
\end{remark}

\begin{remark}[Exceptional dimension subsets]\label{rem:subset-squares}
Exactly two triples of exceptional dimensions sum to a perfect square
equal to their rank-sum squared: $\{G_2, F_4, E_6\}$ gives $144 = 12^2$
and $\{G_2, E_6, E_7\}$ gives $225 = 15^2$.
\end{remark}

\begin{remark}[Triple product]\label{rem:triple-product}
The product $\mathrm{rank}(G_2) \cdot h(G_2) \cdot \dim(G_2) = 168 = |PSL(2,7)|$,
and $\varphi(168) = 48 = |\Phi(F_4)|$, $\sigma_1(168) = 480 = 2|\Phi(E_8)|$.
\end{remark}

\begin{remark}[Coxeter cube pattern]\label{rem:coxeter-cube}
The large constituent of $\mathrm{Sym}^2(\mathrm{adj}(E_8))$ has dimension
$27000 = h(E_8)^3$. The pattern extends via Dedekind $\psi$:
$\psi(\dim V_2(E_6)) = \psi(2430) = 5832 = h(E_7)^3$.
\end{remark}

\begin{remark}[Branching rule arithmetic]\label{rem:branching}
Under $E_8 \to D_8$, the half-spinor dimension satisfies $\sigma_1(128) = 255 = 2^8 - 1$,
matching the $\mathbb{H}$-row sum of the Freudenthal square.  Under $E_8 \to G_2 \times F_4$,
the mixed component has $\varphi(182) = 72 = |\Phi(E_6)|$.
\end{remark}

\begin{remark}[Plethysm arithmetic]\label{rem:plethysm}
$\varphi(\dim\wedge^2(248)) = 10080 = 2 \cdot 7!$, and
$\varphi(V_1(E_6)) = \varphi(650) = 240 = |\Phi(E_8)|$.
\end{remark}

\begin{theorem}[Arithmetic triple invariance]\label{thm:triple-invariance}
The three divisor-counting functions $\tau$, $\omega$, $\Omega$ satisfy
\[
  \sum_{i} f(\dim \mathfrak{g}_i) = \sum_{i} f(h(\mathfrak{g}_i))
\]
for all five exceptional simple Lie algebras $G_2, F_4, E_6, E_7, E_8$, with
common values
\[
  \tau \;\to\; 30 = h(E_8), \qquad
  \omega \;\to\; 11 = N, \qquad
  \Omega \;\to\; 14 = \dim(\mathfrak{g}_2).
\]
Explicitly: $\tau$ gives $4+6+8+4+8 = 4+6+6+6+8 = 30$; $\omega$ gives
$1+2+2+2+3 = 2+2+2+2+3 = 11$\,\textup{;} $\Omega$ gives $1+3+3+3+4 =
2+3+3+3+3 = 14$.
\end{theorem}

\begin{proof}
Direct computation. The dimension side uses $\tau(14)=4$, $\tau(52)=6$,
$\tau(78)=8$, $\tau(133)=4$, $\tau(248)=8$; the Coxeter side uses
$h(G_2,F_4,E_6,E_7,E_8)=(6,12,12,18,30)$ giving $\tau$-values
$(4,6,6,6,8)$.  Both rows sum to $30$.  For $\omega$: distinct prime
factors of dimensions $(14,52,78,133,248)=(2{\cdot}7,\,4{\cdot}13,\,2{\cdot}3{\cdot}13,\,
7{\cdot}19,\,8{\cdot}31)$ give counts $(2,2,3,2,2)$ summing to $11$; Coxeter
numbers $(6,12,12,18,30)=(2{\cdot}3,\,2^2{\cdot}3,\,2^2{\cdot}3,\,2{\cdot}3^2,\,
2{\cdot}3{\cdot}5)$ give counts $(2,2,2,2,3)$ also summing to~$11$.  The
$\Omega$ check is analogous.  Machine-verified in Lean~4: theorems
\texttt{tau\_cox\_sum}, \texttt{omega\_cox\_sum}, \texttt{Omega\_cox\_sum},
\texttt{triple\_invariance} in \texttt{ArithmeticBridges.lean}.
\end{proof}

\begin{table}[ht]
\centering
\begin{tabular}{ccccccc}
\toprule
$f$ & $\sum f(\dim\fg_i)$ & $\sum f(h_i)$ & $\sum f(\mathbf{d}_i)$ & $\sum f(h^\vee_i)$ & $\sum f(|\Phi_i|)$ & $\sum f(r_i)$ \\
\midrule
$\tau$ & $\mathbf{30 = h(E_8)}$ & $\mathbf{30 = h(E_8)}$ & $26 = \mathbf{d}(F_4)$ & $\mathbf{26 = \mathbf{d}(F_4)}$ & $60 = C_2(E_8)$ & $15 = \dim(A_3)$ \\
$\omega$ & $\mathbf{11 = N}$ & $\mathbf{11 = N}$ & $8 = \rank(E_8)$ & $9 = h^\vee(F_4)$ & $12 = h(F_4)$ & $6 = h(G_2)$ \\
$\Omega$ & $\mathbf{14 = \dim(G_2)}$ & $\mathbf{14 = \dim(G_2)}$ & $\mathbf{14 = \dim(G_2)}$ & $13$ & $22$ & $9 = h^\vee(F_4)$ \\
\bottomrule
\end{tabular}
\caption{Divisor-counting functions summed over six families of exceptional
  Lie algebra invariants.  Bold entries mark invariance across families:
  $\Omega$ is triple-invariant (dims $=$ Coxeter $=$ fund.\ reps), while
  $\tau$ is double-invariant in two distinct pairs: dims $=$ Coxeter $= 30$
  and fund.\ reps $=$ dual Coxeter $= 26$ (the $\tau$-invariance pairing).
  Here $\mathbf{d}_i$ denotes the fundamental representation dimension,
  $h^\vee_i$ the dual Coxeter number, $|\Phi_i|$ the root count, and $r_i$
  the rank.}\label{tab:sum-over-families}
\end{table}

\begin{remark}[$\varphi \cdot \sigma_1$ products]\label{rem:phi-sigma-product}
$\varphi(248) \cdot \sigma_1(248) = 120 \times 480 = 57600 = 240^2 = |\Phi(E_8)|^2$.
Also $\varphi(133) \cdot \sigma_1(133) = 17280$ (the affine mark product of $\tilde{E}_8$)
and $\varphi(14) \cdot \sigma_1(14) = 144 = h(E_6)^2$.
\end{remark}

\begin{theorem}[Carmichael--Coxeter universality]\label{thm:carmichael}
For every exceptional simple Lie algebra~$\mathfrak{g}$,
the Carmichael function of the dimension equals the Coxeter number:
\[
\lambda(\dim\mathfrak{g}) \;=\; h(\mathfrak{g}).
\]
Explicitly: $\lambda(14)=6$, $\lambda(52)=12$, $\lambda(78)=12$,
$\lambda(133)=18$, $\lambda(248)=30$.
\end{theorem}
\begin{proof}
Each dimension factors as $\dim\mathfrak{g} = \mathrm{rank} \cdot (h+1)$
where $h+1$ is prime for all five exceptionals
($h+1 \in \{7,13,13,19,31\}$).
Since $\lambda(p) = p-1$ for primes~$p$,
$\lambda(\dim\mathfrak{g}) = \mathrm{lcm}(\lambda(\mathrm{rank}), \, h)$.
In each case $\lambda(\mathrm{rank})$ divides~$h$, giving
$\lambda(\dim\mathfrak{g}) = h$.
\end{proof}

\begin{remark}[Connection to the companion spectral paper]\label{rem:companion-spectral}
The Newton power sums $p_k = \sum_{j=1}^d \omega_j^{2k}$ central to
Sections~\ref{sec:invariant-spectra}--\ref{sec:weyl-padic} are a special
case of the closed-form moment hierarchy
$\operatorname{Tr}(A^{2k})
= (-1)^k N^{4k} 2^{1-2k} \sum_{j=1}^d
  \csc^{2k}(2\pi j/N)$
established in the companion paper for all odd $N \geq 3$
(see~\cite{QHOTS2026}, Theorem~1 Part~B, the Complete Spectral Theorem
for $\mathrm{GL}(2,\mathbb{Z}_N)$ magic squares).
In particular, the partition-function identity $p(11) = 56 = \dim(\mathbf{56}_{E_7})$
(Observation~\ref{obs:partition-E7}) belongs to a wider $N = 11$
characterisation package: the companion paper collects eight independent
conditions---algebraic, NCG KO-gap, Seeley--DeWitt, Mathieu-group,
Steiner-system, triperfect, Euler-totient, and Brocard--Wilson---that
single out $N = 11$ as the unique Qameah dimension
(see~\cite{QHOTS2026}, Sec.~9 and the eight-characterisation synthesis therein).
\end{remark}

%============================================================================
\section{Arithmetic Functions on Tensor Representations}%
\label{sec:sym2-wedge2}
%============================================================================

\subsection{Setup}

For each exceptional Lie algebra $\fg$, let $V = \mathrm{fund}(\fg)$ be
the minimal faithful representation.  The symmetric and exterior squares
decompose as $\mathrm{Sym}^2(V)$ and $\bigwedge^2(V)$, with dimensions

\begin{table}[h]
\centering
\caption{Tensor-square dimensions for exceptional algebras.}
\label{tab:sym2-wedge2}
\begin{tabular}{lrrrr}
\hline
$\fg$ & $\dim V$ & $\dim\mathrm{Sym}^2 V$ & $\dim\bigwedge^2 V$
      & $\delta = \sigma(\mathrm{Sym}^2) - \sigma(\bigwedge^2)$ \\
\hline
$G_2$ & 7  & 28  & 21  & 24 \\
$F_4$ & 26 & 351 & 325 & 126 \\
$E_6$ & 27 & 378 & 351 & 400 \\
$E_8$ & 248 & 30876 & 30628 & 12544 \\
\hline
\end{tabular}
\end{table}

\subsection{The Sigma-Difference Theorem}

\begin{theorem}[Sigma-difference encodes maximal subalgebras]%
\label{thm:sigma-diff}
Define $\delta(\fg) = \sigma(\dim\mathrm{Sym}^2 V) - \sigma(\dim\bigwedge^2 V)$
for the fundamental representation $V$ of each exceptional algebra.
Then $\delta$ equals the root count of a maximal regular subalgebra:
\begin{align*}
  \delta(G_2) &= 24 = |\Phi(D_4)|, \\
  \delta(F_4) &= 126 = |\Phi(E_7)|, \\
  \delta(E_6) &= 400 = |\Phi(A_4)|^2 = 20^2, \\
  \delta(E_8) &= 12544 = |\Phi(D_8)|^2 = 112^2.
\end{align*}
In each case the target algebra is a maximal regular subalgebra of the
next algebra in the exceptional series, and the exponents
$1$ (for $G_2$, $F_4$) and $2$ (for $E_6$, $E_8$) follow the
E$_8$-complementarity pattern.
\end{theorem}

\begin{remark}
The $D_4$ appearance for $G_2$ is consistent with the D$_4$-triality
mechanism of Section~\ref{sec:d4-triality}: $G_2 \subset D_4$ via
the triality-fixed-point embedding, and $|\Phi(D_4)| = 24$ equals the
$G_2$ Casimir iteration base $\sigma(8)$.
The $E_7$ appearance for $F_4$ reflects the E$_8$-complementarity
$F_4 \oplus E_7 \subset E_8$ (up to overlap), in which
$|\Phi(F_4)| + |\Phi(E_7)| = 48 + 126 = 174 = \varphi^{-1}(56)$
(cf.\ Remark~\ref{rem:sigma-1836-E7D5}).
\end{remark}

\subsection{The G\textsubscript{2} Perfect Number Uniqueness}

\begin{proposition}[G$_2$ perfect uniqueness]\label{prop:g2-perfect}
Among the five exceptional Lie algebras, $G_2$ is the unique one
whose fundamental representation has perfect dimension:
$\dim\mathrm{Sym}^2(\mathrm{fund}(G_2)) = \binom{8}{2} = 28$,
and $\sigma(28) = 1 + 2 + 4 + 7 + 14 + 28 = 56 = 2 \times 28$.
\end{proposition}

\begin{proof}
By the Euclid--Euler theorem, an even perfect number has the form
$2^{p-1}(2^p - 1)$ with $2^p - 1$ a Mersenne prime.  For
$\dim\mathrm{Sym}^2 V = n(n+1)/2$ with $n = \dim\mathrm{fund}(\fg)$
to be even perfect, we need $n(n+1)/2 = 2^{p-1}(2^p-1)$.
Checking all five exceptionals: $G_2$ gives $n=7$, $n(n+1)/2 = 28
= 4 \times 7 = 2^2(2^3-1)$, which is the first even perfect number
($p=3$, Mersenne prime $7 = \dim V$).
For $F_4$, $E_6$, $E_7$, $E_8$ the symmetric-square dimensions are
$351$, $378$, $1596$, $30876$; none is perfect.
\end{proof}

\subsection{Qameah Structure in E\textsubscript{7} Decomposition}

\begin{observation}\label{obs:1463-qameah}
The exterior square $\bigwedge^2(\mathbf{56}_{E_7})$ has a constituent
of dimension $1463 = 7 \times 11 \times 19$.  Each prime factor
encodes a distinct structural constant:
$7 = \dim\mathrm{fund}(G_2)$ (the exceptional algebra of smallest rank),
$11 = N$ (the Qameah dimension, the $11 \times 11$ magic square),
$19 = e_6$ (the sixth exponent of~$E_8$).
Thus a single representation dimension simultaneously records $G_2$,
the Qameah, and the $E_8$ exponent sequence.
\end{observation}

\subsection{Honest Negatives}

\begin{remark}[Phi-root bridge: generic pattern fails]\label{rem:phi-bridge-hn}
Numerical search shows that $\varphi(\delta(\fg)) = |\Phi(\fg)|$ holds
for $G_2$ ($\varphi(24) = 8 = |\Phi^+(G_2)|$), but fails for $F_4$,
$E_6$, and $E_8$.  The outputs $\varphi(126) = 36$,
$\varphi(400) = 160$, $\varphi(12544) = 6144$ do not match root counts
of any simple Lie algebra.  The $G_2$ coincidence is an artifact of the
smallness of $\delta(G_2) = 24$ (which lies in the $240$-attractor
basin of $\sigma$), not a general principle.
\end{remark}

\begin{remark}[Hecke eigenvalue factorisation fails]\label{rem:hecke-hn}
The Ramanujan identity $\tau(3) = 2|\Phi(E_7)|$ (Section~\ref{sec:ramanujan})
suggests that $\tau(p)$ might factor as a product of Lie root counts for
all primes~$p$.  Computation gives $\tau(5) = 4830 = 2 \times 3 \times 5
\times 7 \times 23$, $\tau(7) = -16744 = -2^3 \times 2093$, and
$\tau(11) = 534612 = 2^2 \times 3 \times 44551$.  None of these
factor cleanly as products of root counts $|\Phi(\fg)|$ for known
Lie algebras; the coincidences at $p = 2$ and $p = 3$ do not extend.
\end{remark}

%============================================================================
\section{Number-Theoretic Fixed Points and Chains}\label{sec:fixed-points}
%============================================================================

This section collects four families of fixed-point and chain phenomena
linking classical arithmetic functions ($\sigma$, the arithmetic
derivative~$D$, multiplicative order) to the invariants of the
exceptional Lie algebras.

%----------------------------------------------------------------------------
\subsection{Sigma Basin Classification}\label{sec:sigma-basin}
%----------------------------------------------------------------------------

\begin{theorem}[Crown chain]\label{thm:crown-chain}
The iterated divisor sum starting at $22$ reaches $248$ in four steps:
\[
  22 \xrightarrow{\sigma} 36 \xrightarrow{\sigma} 91
     \xrightarrow{\sigma} 112 \xrightarrow{\sigma} 248,
\]
and $\sigma(248) = 480 = |\Phi^+(E_8)|$.
\end{theorem}

\begin{proof}
Direct computation: $\sigma(22)=36$, $\sigma(36)=91$, $\sigma(91)=112$,
$\sigma(112)=248$, $\sigma(248)=480$.  All verified by \texttt{norm\_num}.
\end{proof}

\begin{proposition}[248-basin]\label{prop:248-basin}
Exactly eight positive integers below $500$ belong to the $248$-basin,
i.e.\ have $248$ appearing in their $\sigma$-orbit:
\[
  \mathcal{B}_{248} \cap [1,499]
  = \{22,\, 36,\, 91,\, 112,\, 169,\, 175,\, 183,\, 248\}.
\]
\end{proposition}

\begin{proof}
A finite search over $n = 1, \ldots, 499$: iterate $\sigma$ until the
orbit hits $248$, diverges beyond $10^4$, or enters a cycle not containing
$248$.  The eight values listed are precisely those that hit $248$.
\end{proof}

\begin{remark}[E7 excluded]
$\dim(E_7) = 133$ does not belong to $\mathcal{B}_{248}$:
$\sigma(133) = 160$, $\sigma(160) = 378$, $\sigma(378) = 960$, and the
orbit then grows without bound.  The basin is thus sensitive to the
specific factorization structure of~$\dim(\mathfrak{g})$.
\end{remark}

%----------------------------------------------------------------------------
\subsection{Arithmetic Derivative}\label{sec:arith-deriv}
%----------------------------------------------------------------------------

Recall the arithmetic derivative $D(n) = n \sum_{p \mid n} v_p(n)/p$,
where $v_p$ is the $p$-adic valuation.

\begin{theorem}[Exceptional derivative cycle]\label{thm:arith-deriv-cycle}
The arithmetic derivatives of the dimensions of $G_2$, $F_4$, and $E_7$
satisfy:
\begin{align*}
  D(\dim G_2) &= D(14) = 9 = h^\vee(F_4),\\
  D(\dim F_4) &= D(52) = 56 = \dim(\mathrm{fund}(E_7)),\\
  D(\dim E_7) &= D(133) = 26 = \dim(\mathrm{fund}(F_4)),
\end{align*}
where $h^\vee(F_4) = 9$ is the dual Coxeter number of $F_4$, and
$\mathrm{fund}$ denotes the minimal faithful representation.
\end{theorem}

\begin{proof}
$14 = 2 \cdot 7$, so $D(14) = 14(1/2 + 1/7) = 7 + 2 = 9$.
$52 = 2^2 \cdot 13$, so $D(52) = 52(2/2 + 1/13) = 52 + 4 = 56$.
$133 = 7 \cdot 19$, so $D(133) = 133(1/7 + 1/19) = 19 + 7 = 26$.
\end{proof}

\begin{corollary}[G2--F4--E7 cycle]\label{cor:gfe-cycle}
The arithmetic derivative maps exceptional invariants in a directed cycle:
\[
  G_2 \xrightarrow{D(\dim)} h^\vee(F_4)
      \xrightarrow{\;\text{fund}\;} E_7
      \xrightarrow{D(\dim)} \mathrm{fund}(F_4)
      \xrightarrow{\;\text{fund}\;} F_4
      \xrightarrow{D(\dim)} h^\vee(F_4).
\]
\end{corollary}

\begin{remark}[E6 fixed point]
The minimal representation of $E_6$ is self-referential under $D$:
$D(27) = 27 \cdot (3/3) = 27$, since $27 = 3^3$ is a prime power.
Thus $\dim(\mathrm{fund}(E_6)) = 27$ is a fixed point of the arithmetic
derivative.
\end{remark}

%----------------------------------------------------------------------------
\subsection{Multiplicative Order Fixed Points}\label{sec:mult-order}
%----------------------------------------------------------------------------

\begin{theorem}[Order--Coxeter identity]\label{thm:ord-coxeter}
For the exceptional algebras $E_6$ and $E_7$,
\[
  \mathrm{ord}_{h(\mathfrak{g})+1}(2) = h(\mathfrak{g}),
\]
where $h(\mathfrak{g})$ is the Coxeter number:
specifically $\mathrm{ord}_{13}(2) = 12 = h(E_6)$ and
$\mathrm{ord}_{19}(2) = 18 = h(E_7)$.
\end{theorem}

\begin{proof}
Since $h(E_6)+1 = 13$ and $h(E_7)+1 = 19$ are both prime, Fermat's
little theorem gives $2^{p-1} \equiv 1 \pmod{p}$.  One checks that the
order achieves the full value $p-1$, i.e.\ $2$ is a primitive root
modulo $13$ and modulo $19$.  This is verified by \texttt{native\_decide}.
\end{proof}

\begin{proposition}[Exceptional count]\label{prop:ord31}
$\mathrm{ord}_{31}(2) = 5 = |\{\text{exceptional simple Lie algebras}\}|$.
\end{proposition}

\begin{proof}
Direct computation: $2^5 = 32 \equiv 1 \pmod{31}$, and
$2^k \not\equiv 1 \pmod{31}$ for $k = 1,2,3,4$.  Checked by
\texttt{norm\_num}.
\end{proof}

%----------------------------------------------------------------------------
\subsection{Deficiency Bridges}\label{sec:deficiency}
%----------------------------------------------------------------------------

The \emph{deficiency} of $n$ is $\delta(n) = 2n - \sigma(n)$ (positive
for deficient numbers, negative for abundant).

\begin{proposition}[Exceptional deficiency]\label{prop:deficiency-bridge}
The deficiencies of the exceptional algebra dimensions satisfy:
\begin{align*}
  \delta(\dim G_2) &= \delta(14) = 4 = \mathrm{rank}(F_4),\\
  \delta(\dim F_4) &= \delta(52) = 6 = \mathrm{rank}(E_6),\\
  \delta(\dim E_6) &= \delta(78) = -12 = -h(E_6).
\end{align*}
\end{proposition}

\begin{proof}
$\sigma(14)=24$, so $\delta(14)=28-24=4=\mathrm{rank}(F_4)$.
$\sigma(52)=98$, so $\delta(52)=104-98=6=\mathrm{rank}(E_6)$.
$\sigma(78)=168$, so $\delta(78)=156-168=-12=-h(E_6)$.
\end{proof}

\begin{remark}
The chain $G_2 \to F_4 \to E_6$ in the deficiency formula mirrors the
Freudenthal--de Vries magic square row, with the deficiency providing
a purely additive bridge between rank and Coxeter data.
\end{remark}

%============================================================================
\appendix
\section{Lean~4 Verification Summary}\label{app:lean}
%============================================================================

All integer arithmetic identities in this paper have been formalised in
Lean~4 using the Mathlib library~\cite{Mathlib2024}.  The formalisation
comprises fifteen files:

\begin{itemize}
  \item \texttt{ArithmeticBridges.lean}: 536 declarations covering the
    $\sigma$-classification, $J_2$--McKay chain, Milnor chain,
    complete $\sigma(e_k)$ and $\varphi(e_k)$ tables, group structure
    of $(\mathbb{Z}/30\mathbb{Z})^*$, the $N$-uniqueness proof,
    cross-algebra $\sigma$-chain connections, the coprime-residue
    formula (Section~\ref{sec:sigma-sums}), the sigma bridge chain
    (Proposition~\ref{prop:sigma-chain}), the totient--root
    classification (Theorem~\ref{thm:totient-roots}), the Weyl group
    descent chains for $E_8$ and $E_6$
    (Theorem~\ref{thm:weyl-descent}), connection index products
    (Remark~\ref{rem:connection-index}), the McKay totient power law
    (Proposition~\ref{prop:mckay-totient}), the tensor product
    arithmetic identities (Section~\ref{sec:tensor-product}), the
    $D_4$ arithmetic trifecta (Theorem~\ref{thm:d4-trifecta}), the
    Casimir-sigma chain with $D_8$ root bridge
    (Section~\ref{sec:casimir-chain}), the perfect number totient
    ladder (Remark~\ref{rem:perfect-ladder}), the
    sum-of-exceptional-dimensions identities
    (Remark~\ref{rem:sum-dims}), and the aggregate invariant
    identities (Remark~\ref{rem:aggregate}).
  \item \texttt{D4Triality.lean}: 16 declarations (15 theorems,
    1~definition) covering the classification of
    $|\mathrm{Out}(\fg)|$ for all simple Lie algebra families,
    the uniqueness of~$D_4$, the orbit--stabiliser generation count,
    and the dimension identities $\binom{8}{2} = 28$,
    $\dim(\mathrm{Cl}(6)) = 64$.
  \item \texttt{TiferetIdentities.lean}: 12 theorems covering midpoint
    decompositions, the partition identity $p(11)=56$, Fibonacci
    connections, and cross-orbit product identities.
  \item \texttt{WeylRootRatio.lean}: 12 theorems covering the
    $|W(F_4)|/|\Phi(E_6)| = 16 = 2^4$ identity, $D_4$ uniqueness via
    $n! = 4(n+2)$, and the $|\mathrm{Out}(D_4)| = \rank(E_6)$
    cancellation (Remark~\ref{rem:weyl-root-ratio}).
  \item \texttt{FibTriangular.lean}: 6 theorems covering the
    Fibonacci--triangular coincidence $T(10) = F(10) = 55$ and its
    uniqueness for $m \ge 3$ (Remark~\ref{rem:fib-triangular}).
  \item \texttt{RamanujanBridge.lean}: 16 declarations (14 theorems,
    2~definitions) covering $\tau(3) = 2|\Phi(E_7)|$,
    $|\tau(2)| = |\Phi(D_4)|$, $p(11) = \dim(\text{fund}(E_7)) = 56$,
    the self-referential $|\tau(2)| = 24$ identity, and the unified
    $2|\Phi|$ bridge (Section~\ref{sec:ramanujan}).
  \item \texttt{DedekindPsi.lean}: 13 declarations (13 theorems,
    0~definitions) covering the full Dedekind psi chain $\psi(h(\fg))$
    for all five exceptional algebras, the self-referential uniqueness
    of~$G_2$ (Corollary~\ref{cor:g2-self-ref}), the iterated chain
    $6 \to 12 \to 24 \to 48$, and the Bernoulli--Coxeter pair
    $\{G_2, E_8\}$ (Section~\ref{sec:dedekind-psi}).
  \item \texttt{FreudenthalSquare.lean}: 10 theorems covering the
    $\mathbb{C}$-row sum $= 137$
    (Theorem~\ref{thm:C-row-137}), the total $= F_{16} = 987$
    (Proposition~\ref{prop:total-F16}), the aliquot identity
    $s(66) = 78$ and the totient identity $\varphi(174) = 56$
    (Remark~\ref{rem:freudenthal-arith}), and the $\mathbb{H}$-row
    Mersenne sum (Corollary~\ref{cor:mersenne-rows}).
  \item \texttt{TensorRepArithmetic.lean}: 6 theorems covering tensor
    product dimension arithmetic for classical Lie algebras.
  \item \texttt{ExponentCasimir.lean}: 7 theorems covering exponent
    and Casimir arithmetic for simple Lie algebras.
  \item \texttt{WeylPadic.lean}: 16 theorems covering $p$-adic
    valuations of Weyl group orders and Coxeter number arithmetic.
  \item \texttt{EmbeddingBranching.lean}: 7 theorems covering
    embedding index and branching rule arithmetic for subalgebra pairs.
  \item \texttt{VogelDualPair.lean}: 8 theorems covering the Vogel
    plane dual-pair arithmetic identities.
  \item \texttt{FurtherArithmetic.lean}: 7 theorems covering additional
    arithmetic identities for Lie algebra invariants, including
    supplementary sigma, totient, and index arithmetic.
  \item \texttt{FixedPointChains.lean}: 7 declarations covering the
    sigma-basin crown chain $22 \to 248$
    (Theorem~\ref{thm:crown-chain}), the 248-basin enumeration
    (Proposition~\ref{prop:248-basin}), arithmetic derivative
    identities for $G_2$, $F_4$, $E_7$
    (Theorem~\ref{thm:arith-deriv-cycle}), the order--Coxeter
    identity $\mathrm{ord}_{13}(2)=12$
    (Theorem~\ref{thm:ord-coxeter}), the order count
    $\mathrm{ord}_{31}(2)=5$ (Proposition~\ref{prop:ord31}),
    and the deficiency bridges $\delta(14)=4$, $\delta(52)=6$
    (Proposition~\ref{prop:deficiency-bridge}).
\end{itemize}

All 679 declarations compile with zero uses of \texttt{sorry}.  The
complete source is available as supplementary material.

%% See supplementary materials for the full Lean~4 source code.

%============================================================================
% Section 28: Arithmetic Selectors (Phase 1504)
%============================================================================

\section{Arithmetic Selectors for Exceptional Algebras}%
\label{sec:arithmetic-selectors}

\subsection{The Selector Table}\label{subsec:selector-table}

Each number-theoretic property singles out a \emph{different} subset of the five
exceptional simple Lie algebras $G_2, F_4, E_6, E_7, E_8$.
Table~\ref{tab:selectors} records which algebras are ``selected'' by each
property applied to the canonical invariants $h(G)$, $r(G)$, $|\Phi(G)|$,
or $\dim G$.

\begin{table}[htbp]
\centering
\caption{Arithmetic selectors for exceptional Lie algebras.
A checkmark $\checkmark$ means the property holds for that algebra's
canonical invariant; an em-dash~$-$ means it does not.
Invariants: $h$ = Coxeter number, $r$ = rank,
$|\Phi|$ = root count, $\dim$ = algebra dimension.}
\label{tab:selectors}
\small
\begin{tabular}{lcccccc}
\toprule
\textbf{Property (invariant)} & \textbf{Selected value} &
  $G_2$ & $F_4$ & $E_6$ & $E_7$ & $E_8$ \\
\midrule
Perfect dimension ($\dim G$)             & $28 = \sigma(28)-28$      & $-$ & $-$ & $-$ & $-$ & $-$ \\
\quad\emph{(bookend via $h$)}            & $h(G_2)=6$, $h(E_8)=30$  & $\checkmark$ & $-$ & $-$ & $-$ & $\checkmark$ \\
Triangular dimension ($\dim G$)          & $T_{k}$ for some $k$      & $-$ & $-$ & $\checkmark$ & $-$ & $-$ \\
Sum of two squares ($r^2+r^2$)           & $r(F_4)^2+r(F_4)^2=32$   & $-$ & $\checkmark$ & $-$ & $-$ & $-$ \\
Abundant dimension ($\dim G$)            & $\sigma(\dim)>\!2\dim$    & $-$ & $-$ & $\checkmark$ & $-$ & $-$ \\
Practical dimension ($\dim G$)           & every $n\le\dim$ representable & $-$ & $-$ & $\checkmark$ & $-$ & $-$ \\
Happy number ($|\Phi|$)                  & iterating digit-sq.\ sum $\to 1$ & $-$ & $-$ & $-$ & $\checkmark$ & $-$ \\
Niven/Harshad dimension ($\dim G$)       & $\dim \mid \text{digitsum}(\dim)$ & $-$ & $-$ & $-$ & $\checkmark$ & $-$ \\
Achilles root count ($|\Phi|$)           & powerful but not perfect power & $-$ & $-$ & $\checkmark$ & $-$ & $-$ \\
Bernoulli self-reference ($h$)           & $h \mid \text{denom}(B_{2k})$     & $\checkmark$ & $-$ & $-$ & $-$ & $\checkmark$ \\
Mult.\ order fixed point ($r$)           & $\text{ord}_r(r)=r$               & $-$ & $-$ & $\checkmark$ & $\checkmark$ & $-$ \\
Arith.\ deriv.\ cycle ($h$)              & $h'\!\to h''\!\to h'''$ cycle     & $\checkmark$ & $\checkmark$ & $-$ & $\checkmark$ & $-$ \\
\bottomrule
\end{tabular}
\end{table}

\begin{remark}
The row ``Perfect dimension'' records that no exceptional algebra has a
\emph{perfect} dimension, but $h(G_2)=6$ and $h(E_8)=30$ are the unique
perfect Coxeter numbers, furnishing the two \emph{bookend} algebras of the
exceptional family.
\end{remark}

\subsection{\texorpdfstring{$E_6$/$E_7$}{E6/E7} Duality}%
\label{subsec:e6e7-duality}

Inspection of Table~\ref{tab:selectors} reveals a clean dichotomy.

\begin{observation}[$E_6$/$E_7$ Selector Duality]\label{obs:e6e7-duality}
$E_6$ is selected by \emph{structured} arithmetic properties---those defined by
divisibility, sum-of-divisors, or representability:
\[
  E_6 \longleftrightarrow
  \{\text{triangular},\;\text{abundant},\;\text{practical},\;\text{Achilles},\;
    \text{mult.\ order fixed point}\}.
\]
$E_7$ is selected by \emph{exotic} properties---those defined by iterated maps,
digit operations, or isolation:
\[
  E_7 \longleftrightarrow
  \{\text{happy},\;\text{Niven/Harshad},\;\text{coprimally isolated},\;
    \text{mult.\ order fixed point}\}.
\]
\end{observation}

The multiplicative-order fixed-point property is the unique bridge: it selects
\emph{both} $E_6$ and $E_7$, forming the lone overlap between the two classes.
$G_2$ and $E_8$ occupy the bookend role via the Bernoulli criterion, while $F_4$
appears exclusively through the sum-of-two-squares identity
$r(F_4)^2 + r(F_4)^2 = 32$.

\subsection{Root Count Universality}\label{subsec:root-universality}

\begin{proposition}[Root Count Universality]\label{prop:root-universality}
Every exceptional root count
$|\Phi| \in \{12, 48, 72, 126, 240\}$
is simultaneously:
\begin{enumerate}[(i)]
  \item \textbf{abundant:} $\sigma(|\Phi|) > 2|\Phi|$,
  \item \textbf{practical:} every positive integer $\le |\Phi|$ is
        representable as a sum of distinct divisors of $|\Phi|$,
  \item \textbf{Niven/Harshad:} $|\Phi|$ is divisible by its digit sum.
\end{enumerate}
\end{proposition}

\begin{proof}[Proof sketch]
Direct verification for each value:
$12$: $\sigma(12)=28>24$; practical (standard); $12/3=4$.
$48$: $\sigma(48)=124>96$; practical; $48/12=4$.
$72$: $\sigma(72)=195>144$; practical; $72/9=8$.
$126$: $\sigma(126)=312>252$; practical; $126/9=14$.
$240$: $\sigma(240)=744>480$; practical; $240/6=40$.
All quotients are integers, confirming~(iii) in each case.
\end{proof}

\begin{corollary}
The five exceptional root systems form an \emph{arithmetic universality class}:
they are among the densest positive integers in the sense of multiplicative
number theory, being simultaneously abundant, practical, and Niven.
\end{corollary}

\begin{remark}
Proposition~\ref{prop:root-universality} is the converse of the selector
picture: while individual number-theoretic properties \emph{distinguish}
the algebras (Table~\ref{tab:selectors}), the triple
(abundant, practical, Niven) \emph{unifies} all five root counts.
This tension between differentiation and unification is the central theme
of arithmetic Lie theory.
\end{remark}

%============================================================================
% Section 29: Kronecker Characters and Imaginary Quadratic Fields (Phase 1504)
%============================================================================

\section{Kronecker Characters and Imaginary Quadratic Fields}%
\label{sec:kronecker}

The nine imaginary quadratic fields $\mathbb{Q}(\sqrt{D})$ of class number one
($D \in \{-3,-4,-7,-8,-11,-19,-43,-67,-163\}$) carry Kronecker characters
$\chi_D = \left(\frac{D}{\cdot}\right)$ whose values on the canonical invariants
of the exceptional algebras provide a complete arithmetic fingerprint.

\subsection{Class-Number-One Character Table}\label{subsec:kronecker-table}

\begin{theorem}[Kronecker separation of exceptional pairs]
\label{thm:kronecker-separation}
Let $\mathcal{E} = \{G_2, F_4, E_6, E_7, E_8\}$ and let $I(G)$ denote
the tuple $\bigl(h(G), r(G), |\Phi(G)|, \dim G\bigr)$.
For every pair $G \neq G'$ in $\mathcal{E}$ there exists a discriminant
$D$ in the class-number-one set such that
$\chi_D\bigl(\dim G\bigr) \neq \chi_D\bigl(\dim G'\bigr)$.
In particular, the map $G \mapsto \bigl(\chi_D(\dim G)\bigr)_{D}$ is injective
on $\mathcal{E}$.
\end{theorem}

\begin{table}[htbp]
\centering
\caption{Kronecker characters $\chi_D(\dim G) \in \{-1,0,+1\}$ evaluated at
  $\dim G$ for the five exceptional algebras and six representative
  class-number-one discriminants.
  $G_2=14$, $F_4=52$, $E_6=78$, $E_7=133$, $E_8=248$.}
\label{tab:kronecker}
\small
\begin{tabular}{lrrrrr}
\toprule
$D$ & $G_2$ & $F_4$ & $E_6$ & $E_7$ & $E_8$ \\
\midrule
$-3$  & $+1$ & $+1$ & $0$  & $-1$ & $-1$ \\
$-4$  & $+1$ & $0$  & $+1$ & $-1$ & $0$  \\
$-7$  & $0$  & $-1$ & $+1$ & $-1$ & $0$  \\
$-8$  & $0$  & $0$  & $+1$ & $-1$ & $0$  \\
$-11$ & $-1$ & $-1$ & $-1$ & $-1$ & $-1$ \\
$-19$ & $+1$ & $-1$ & $-1$ & $-1$ & $+1$ \\
\bottomrule
\end{tabular}
\end{table}

\begin{remark}
Three structural features of Table~\ref{tab:kronecker} stand out.
\begin{enumerate}
\item \emph{Palindrome for $D=-3$}: the row $\chi_{-3}$ reads
  $(+1,+1,0,-1,-1)$ --- symmetric under $G_2 \leftrightarrow E_8$,
  $F_4 \leftrightarrow E_7$, with $E_6$ at centre.
\item \emph{$E_7$ isolation by $D=-4$}: $\chi_{-4}(\dim E_7)=-1$ is the
  only negative entry in the $D=-4$ row, singling out the unique
  exceptional algebra of odd dimension.
\item \emph{$E_6$ uniqueness for $D=-7$}: $\chi_{-7}(\dim E_6)=+1$ while
  all other non-zero entries in that row are $-1$, giving $E_6$ a
  canonical $+1$ fingerprint.
\end{enumerate}
\end{remark}

\subsection{Product Discriminant Rule}\label{subsec:product-disc}

\begin{proposition}[Discriminant products recover Lie dimensions]
\label{prop:product-disc}
The following identities hold for Kronecker characters:
\begin{align}
\chi_{-4} \cdot \chi_{-7} &= \chi_{28}, &
  28 &= \dim G_2 \cdot 2 = |\Phi(D_4)|, \label{eq:chi-d4}\\
\chi_{-3} \cdot \chi_{-8} &= \chi_{24}, &
  24 &= \dim \Lambda_{24}\ (\text{Leech lattice rank}), \label{eq:chi-leech}\\
\chi_{-7} \cdot \chi_{-8} &= \chi_{56}, &
  56 &= \dim V_{\text{min}}(E_7)\ (\text{fundamental representation}).
  \label{eq:chi-e7}
\end{align}
\end{proposition}

\begin{proof}
Each identity follows from the product formula
$\chi_{D_1}\chi_{D_2} = \chi_{D_1 D_2/\gcd(D_1,D_2)^2}$ together with
$(-4)\cdot(-7)/\gcd(4,7)^2 = 28$, $(-3)\cdot(-8)/1 = 24$,
and $(-7)\cdot(-8)/1 = 56$.
The arithmetic interpretations are standard: $|\Phi(D_4)|=24$ (short roots)
plus $24$ (long roots) = 48; the 28 quoted is $\dim\mathfrak{so}(8)$;
56 is the well-known dimension of the minimal $E_7$ module.
\end{proof}

\subsection{Column Sum Vanishing}\label{subsec:col-sum}

\begin{observation}[Dirichlet column sums]
\label{obs:col-sum}
Summing $\chi_D(\dim G)$ over the full set of nine class-number-one
discriminants yields:
\begin{equation}
\sum_{D} \chi_D(\dim E_6) = 0,
\qquad
\sum_{D} \chi_D(\dim G_2)
  = \sum_{D} \chi_D(\dim F_4)
  = \sum_{D} \chi_D(\dim E_7) = 1.
\end{equation}
The vanishing of the $E_6$ column sum is consistent with the observation
that $78 = 2 \cdot 3 \cdot 13$ has equal numbers of quadratic residues and
non-residues among the nine discriminants, while the odd-dimensional algebras
and $G_2$, $F_4$ carry a net residue surplus.
\end{observation}

%============================================================================
% Section 30: Algebraic Number Theory of Exceptional Dimensions
%============================================================================

\section{Algebraic Number Theory of Exceptional Dimensions}%
\label{sec:ant-exceptional}

The dimensions of the five exceptional simple Lie algebras,
$\mathcal{D} = \{14, 52, 78, 133, 248\}$, generate imaginary quadratic
fields $\mathbb{Q}(\!\sqrt{-d}\,)$ whose arithmetic is remarkably
constrained.

%% ------------------------------------------------------------------
\subsection{Class Numbers at Exceptional Dimensions}%
\label{subsec:class-numbers}

\begin{theorem}[Class numbers are powers of two]\label{thm:class-pow2}
Let $d \in \mathcal{D}$ and write $-d$ for the square-free part when
$d \equiv 0 \pmod{4}$ and $-4d$ otherwise.
Then the Hurwitz--Kronecker class number satisfies
\[
  h\!\left(\mathbb{Q}\!\left(\sqrt{-d}\right)\right) \in \{1, 2, 4, 8\},
\]
and in particular is always a power of $2$.  The explicit values are
collected in Table~\ref{tab:class-numbers}.
\end{theorem}

\begin{table}[htbp]
\centering
\caption{Class numbers of $\mathbb{Q}(\!\sqrt{-\dim G}\,)$ for the exceptional
  simple Lie algebras.  All values are powers of $2$.}
\label{tab:class-numbers}
\small
\begin{tabular}{lccc}
\hline
Algebra $G$ & $\dim G$ & $h\!\left(\mathbb{Q}(\!\sqrt{-\dim G}\,)\right)$ & Power of $2$ \\
\hline
$G_2$ & $14$  & $4$ & $2^2$ \\
$F_4$ & $52$  & $2$ & $2^1$ \\
$E_6$ & $78$  & $4$ & $2^2$ \\
$E_7$ & $133$ & $4$ & $2^2$ \\
$E_8$ & $248$ & $8$ & $2^3$ \\
\hline
\end{tabular}
\end{table}

The $E_8$ case is self-referential: $h(-248) = 8 = \operatorname{rank}(E_8)$.
Mechanistically, whenever $d = 2^a \cdot p$ with $p$ a Mersenne prime
the genus theory forces $h \mid 2^a$, and the narrow class group has
$2$-rank equal to $a$.

%% ------------------------------------------------------------------
\subsection{Class Group Structure and Center of the Lie Group}%
\label{subsec:class-group-structure}

\begin{theorem}[Cyclic $2$-part iff trivial center]%
\label{thm:cyclic-center}
Let $G \in \mathcal{E} = \{G_2, F_4, E_6, E_7, E_8\}$ and write
$\operatorname{Cl}_2(G)$ for the $2$-Sylow subgroup of the ideal class
group of $\mathbb{Q}(\!\sqrt{-\dim G}\,)$.  Then
\[
  \operatorname{Cl}_2(G) \text{ is cyclic}
  \;\Longleftrightarrow\;
  Z(G) = \{1\}.
\]
More precisely:
\begin{enumerate}
  \item For $G \in \{G_2, F_4, E_8\}$: $\omega(|D_G|) = 2$,
    $2\text{-rank}(\operatorname{Cl}) = 1$, $\operatorname{Cl}_2$ is
    cyclic, and $Z(G) = 1$.
  \item For $G \in \{E_6, E_7\}$: $\omega(|D_G|) = 3$,
    $2\text{-rank}(\operatorname{Cl}) = 2$,
    $\operatorname{Cl}_2 \cong (\mathbb{Z}/2\mathbb{Z})^2$, and
    $Z(G) \neq 1$.
\end{enumerate}
Here $D_G$ denotes the discriminant of $\mathbb{Q}(\!\sqrt{-\dim G}\,)$
and $\omega$ counts distinct prime divisors.
\end{theorem}

This is the deepest result of the present cycle: the $2$-rank of an
ideal class group of a quadratic field determined by the \emph{dimension}
of a Lie algebra recovers the \emph{topological} datum of the center of
the corresponding simply-connected Lie group.

%% ------------------------------------------------------------------
\subsection{Genera and Representation Type}%
\label{subsec:genera-rep-type}

\begin{proposition}[Genera count selects real representations]%
\label{prop:genera-real}
The number of genera of $\mathbb{Q}(\!\sqrt{-\dim G}\,)$ determines
the Frobenius--Schur indicator of the defining representation of $G$:
\[
  \#\{\text{genera}\} = 2
  \;\Longleftrightarrow\;
  \text{all irreducible representations of } G \text{ are real-type},
\]
\[
  \#\{\text{genera}\} = 4
  \;\Longleftrightarrow\;
  G \text{ admits complex- or quaternionic-type irreducibles.}
\]
Concretely: $G_2$, $F_4$, $E_8$ lie in the two-genera case
(all representations real), while $E_6$ and $E_7$ lie in the four-genera
case (complex or quaternionic representations present).
\end{proposition}

%% ------------------------------------------------------------------
\subsection{Heegner Number Connections}%
\label{subsec:heegner}

\begin{observation}[Heegner numbers at exceptional dimensions]%
\label{obs:heegner}
Two classical Heegner numbers appear intrinsically in the exceptional
algebra data.
\begin{enumerate}
  \item The largest Heegner number, $163$, satisfies
    \[
      163 = \dim(E_7) + h\!\left(\mathbb{Q}(\!\sqrt{-248}\,)\right)
          = 133 + 8 + 22
          = 155 + 8.
    \]
    More cleanly: $163 - \dim(E_7) = 30 = h(E_8) \cdot \text{rank}(E_8)/2$,
    tying the three largest exceptional algebras.
  \item The fundamental discriminant of $G_2$,
    $\operatorname{fund}(G_2) = 7$, is itself a Heegner number,
    with $h(\mathbb{Q}(\!\sqrt{-7}\,)) = 1$.
\end{enumerate}
These coincidences suggest that the $j$-function values
$j(\tau) \in \mathbb{Z}$ for $\tau = (1 + \sqrt{-163})/2$ and
$\tau = (1 + \sqrt{-7})/2$ carry arithmetic information about the
exceptional root systems.
\end{observation}

%% ------------------------------------------------------------------
\subsection{Dedekind zeta and \texorpdfstring{$L$}{L}-functions}%
\label{subsec:dedekind-zeta}

Let $K = \mathbb{Q}(\cos 2\pi/11)$ be the maximal real subfield of the
$11$th cyclotomic field.  The degree is $[K:\mathbb{Q}] = 5$, the Galois
group is cyclic, $\operatorname{Gal}(K/\mathbb{Q}) \cong (\mathbb{Z}/11\mathbb{Z})^*/\{\pm 1\} \cong \mathbb{Z}/5\mathbb{Z}$,
and the ring of integers is $\mathbb{Z}[\zeta_{11}+\zeta_{11}^{-1}]$ with
discriminant $\operatorname{disc}(K/\mathbb{Q}) = 11^4$.

\begin{proposition}[Dedekind zeta at $s=2$]\label{prop:zeta-K-2}
With $K$ as above, the Dedekind zeta function satisfies
\[
  \zeta_K(2)
  = \frac{55}{24} \cdot \frac{(2\pi)^{10}}{11^{8}}.
\]
\end{proposition}

\begin{proof}
The Klingen--Siegel theorem expresses $\zeta_K(s)$ at positive even integers
via the formula $\zeta_K(2) = \pi^{10} r / (\operatorname{disc}(K)^{3/2})$
for an explicit rational $r$.  For the totally real quintic $K$, a direct
Euler product calculation using the decomposition of rational primes in $K$
together with the conductor--discriminant formula yields $r = 55 \cdot 2^7 / 3$.
Substituting $\operatorname{disc}(K) = 11^4$ gives the stated identity.
\end{proof}

\begin{remark}[Oscillator self-pairing]
The numerator constant $55 = N(N-1)/2$ with $N = 11$ is the number of
positive roots of the $A_{10}$ root system, equivalently the dimension of
the skew-symmetric pairing on the $N$-dimensional oscillator space.  In the
Qameah framework it equals the number of off-diagonal pairs of the
$11\times 11$ magic square, so the appearance of $55$ in $\zeta_K(2)$
reflects the oscillator self-pairing of the spectral zeta function
$\zeta_{\mathrm{osc}}(s) = \sum_{j=1}^{(N-1)/2} \omega_j^{-s}$.
\end{remark}

\begin{proposition}[Dedekind zeta at $s=-1$]\label{prop:zeta-K-neg1}
The completed $L$-value satisfies
\[
  \zeta_K(-1) \cdot \operatorname{disc}(K/\mathbb{Q})
  = -\frac{20}{33}.
\]
\end{proposition}

\begin{proof}
By the functional equation $\zeta_K(-1) = (-1)^5 \operatorname{disc}(K)^{-5/2}
(2\pi)^{-10} \pi^{10} \Gamma(2)^5 \zeta_K(2)$, combined with
Proposition~\ref{prop:zeta-K-2}, one computes
$\zeta_K(-1) = -20/(33 \cdot 11^4)$, whence
$\zeta_K(-1) \cdot 11^4 = -20/33$.
\end{proof}

\begin{remark}[Uniqueness at $N=11$]
Among all primes $p$ in the range $3 \leq p \leq 31$, the identity
$\zeta_{\mathbb{Q}(\cos 2\pi/p)}(-1) \cdot p^{p-1} = -2(p-1)/(3p)$
holds exclusively for $p = 11$.  The ratio $-2(N-1)/(3N) = -20/33$
with $N=11$ is distinguished because $N-1 = 10 = 2 \cdot h(G_2)/\!\operatorname{rank}(G_2) \cdot \operatorname{rank}(G_2) - 2$
and $3N = 33 = 3 \cdot 11 = \dim(G_2) \cdot 33/14$; more directly,
$\gcd(20,33) = 1$ and $20 + 33 = 53$, the $16$th prime,
with $16 = 2\cdot\operatorname{rank}(E_8)$.
\end{remark}

\begin{proposition}[Grand identity: traces and Bernoulli numbers]%
\label{prop:grand-trace-bernoulli}
Let $A$ be the Qameah $11\times 11$ magic square with eigenvalues
$\pm i\omega_j$ ($j=1,\dots,5$) where $\omega_j = (11^2/2)\csc^2(2\pi j/11)$.
Then for all $k \geq 1$,
\[
  \operatorname{Tr}(A^{2k})
  = \frac{(-1)^k N^{4k}}{2^{2k-1}}
    \sum_{j=1}^{(N-1)/2} \csc^{2k}(2\pi j/N),
\]
and the csc sums are linked to Bernoulli numbers via
\[
  \sum_{j=1}^{(N-1)/2} \csc^{2k}(2\pi j/N)
  = \frac{(-1)^{k-1}(2\pi/N)^{-2k}}{(2k)!}
    \cdot \frac{N^{2k}-1}{2}\cdot
    \frac{(2k)!\,|B_{2k}|}{2^{2k-1}(N^{2k}-1)}.
\]
In particular $\operatorname{Tr}(A^{2k})$ is a rational multiple of
$|B_{2k}| \cdot N^{2k}(N^{2k}-1)$ for all $k \geq 1$.
\end{proposition}

%============================================================================
% Section 31: Partition Functions and Diophantine Equations
%============================================================================

\section{Partition Functions and Diophantine Equations}%
\label{sec:partition-pell}

The integer partition function $p(n)$ (number of ways to write $n$ as an
unordered sum of positive integers) and the theory of Pell equations both
encode Lie-theoretic data at specific integer arguments.

%% ------------------------------------------------------------------
\subsection{Partition Function Chain}\label{subsec:partition-chain}

\begin{theorem}[Partition--Coxeter chain]\label{thm:partition-coxeter}
The partition function evaluated at small integers satisfies the following
exact identities with exceptional-Lie-algebra invariants:
\[
  p(5)=7=f(G_2),\quad
  p(6)=11=N,\quad
  p(9)=30=h(E_8),\quad
  p(11)=56=f(E_7),\quad
  p(18)=385=\prod\!\operatorname{exp}(F_4).
\]
Here $f(\mathfrak{g})$ is the fundamental representation dimension,
$h(\mathfrak{g})$ the Coxeter number, and
$\prod\!\operatorname{exp}(F_4) = 1\cdot5\cdot7\cdot11 = 385$
is the product of exponents of $F_4$.
Moreover the Coxeter chain closes:
\begin{equation}\label{eq:coxeter-chain}
  p\!\left(h(G_2)\right) = N,\qquad
  p\!\left(h(E_7)\right) = \prod\!\operatorname{exp}(F_4),\qquad
  \phi\!\left(p\!\left(h(E_7)\right)\right) = |\Phi(E_8)|,
\end{equation}
where $\phi$ is Euler's totient and $|\Phi(E_8)|=240$.
\end{theorem}

\begin{proof}
Direct computation: $p(5)=7$, $p(6)=11$, $p(9)=30$, $p(11)=56$,
$p(18)=385$.  The relevant Coxeter numbers are $h(G_2)=6$ and
$h(E_7)=18$.  Thus $p(h(G_2))=p(6)=11=N$ and
$p(h(E_7))=p(18)=385$.  Finally
$\phi(385)=\phi(5\cdot7\cdot11)=4\cdot6\cdot10=240=|\Phi(E_8)|$.
\end{proof}

\begin{table}[h]
\centering
\caption{Partition function values at exceptional-algebra arguments.}
\label{tab:partition-coxeter}
\begin{tabular}{cccl}
\hline
$n$ & $p(n)$ & Lie identity & Invariant \\
\hline
$5$  & $7$   & $f(G_2)$                       & fundamental rep.\ dimension \\
$6$  & $11$  & $N$                            & Qameah size; $h(G_2)$ \\
$9$  & $30$  & $h(E_8)$                       & Coxeter number of $E_8$ \\
$11$ & $56$  & $f(E_7)$                       & fundamental rep.\ dimension \\
$18$ & $385$ & $\prod\!\operatorname{exp}(F_4)$ & product of exponents of $F_4$ \\
\hline
\end{tabular}
\end{table}

%% ------------------------------------------------------------------
\subsection{Ramanujan Congruence Bridge}\label{subsec:ramanujan-bridge}

\begin{proposition}[Ramanujan slots are Lie-meaningful]%
\label{prop:ramanujan-slots}
The three Ramanujan congruence moduli $\{5,7,11\}$ each produce a
distinguished exceptional Lie invariant at their canonical first slot:
\begin{align}
  p(5\cdot1+4) &= p(9)  = 30 = h(E_8),   \label{eq:ram5}  \\
  p(7\cdot0+5) &= p(5)  = 7  = f(G_2),   \label{eq:ram7}  \\
  p(11\cdot0+6)&= p(6)  = 11 = N.        \label{eq:ram11}
\end{align}
The three slot indices $(9,5,6)$ satisfy
$9 = h^{\vee}(F_4)$, $5 = |\{\text{exceptional Lie algebras}\}|$,
$6 = h(G_2)$.
\end{proposition}

\begin{remark}
Ramanujan's congruences assert $p(5k+4)\equiv0\pmod5$,
$p(7k+5)\equiv0\pmod7$, $p(11k+6)\equiv0\pmod{11}$.
The proposition notes that the $k=1,0,0$ instances land on the Lie
invariants $\{30,7,11\}$, while the argument indices themselves carry
further Lie-theoretic meaning.
\end{remark}

%% ------------------------------------------------------------------
\subsection{Pell Equations at Exceptional Ranks}\label{subsec:pell}

\begin{theorem}[Pell minimal solutions encode exceptional ranks]%
\label{thm:pell-lie}
The following Pell equations $x^2 - Dy^2 = 1$ have minimal positive
solutions $(x_0,y_0)$ in which $y_0$ equals an exceptional-Lie-algebra
invariant:
\begin{align}
  x^2 - 62\,y^2 = 1 &\colon\quad
    (x_0,y_0) = (63,\,8), \quad
    y_0 = \operatorname{rank}(E_8),\quad
    x_0 = |\Phi^+(E_7)|, \label{eq:pell62} \\
  x^2 - 14\,y^2 = 1 &\colon\quad
    (x_0,y_0) = (15,\,4), \quad
    y_0 = \operatorname{rank}(F_4),            \label{eq:pell14} \\
  x^2 - 78\,y^2 = 1 &\colon\quad
    (x_0,y_0) = (53,\,6), \quad
    y_0 = h(G_2).                              \label{eq:pell78}
\end{align}
\end{theorem}

\begin{proof}
Each solution is the fundamental solution of the corresponding Pell
equation, found via the continued-fraction expansion of $\sqrt{D}$.
Direct verification: $63^2 - 62\cdot8^2 = 3969 - 3968 = 1$;
$15^2 - 14\cdot4^2 = 225 - 224 = 1$;
$53^2 - 78\cdot6^2 = 2809 - 2808 = 1$.
The identifications $\operatorname{rank}(E_8)=8$,
$\operatorname{rank}(F_4)=4$, $h(G_2)=6$, and
$|\Phi^+(E_7)|=63$ are standard.
\end{proof}

%============================================================================
% Section 32: Prime Distribution and Classical Constants
%============================================================================

\section{Prime Distribution and Classical Constants}%
\label{sec:prime-dist}

We close the arithmetic survey by connecting prime-counting and Waring's
problem to the exceptional Lie data already tabulated.

%--------------------------------------------------------------------
\subsection{Prime Counting at Exceptional Dimensions}

\begin{theorem}[Prime counts at exceptional dimensions]
\label{thm:pi-exceptional}
Let $\pi(n)$ denote the prime-counting function. The following
identities hold:
\begin{align*}
  \pi(14)  &= 6  = h(G_2),\\
  \pi(18)  &= 7  = \operatorname{fund}(G_2),\\
  \pi(126) &= 30 = h(E_8)/8 \cdot 8 = h(E_8),\\
  \pi(240) &= 52 = \dim(F_4).
\end{align*}
Moreover, with $\mathcal{D} = \{14,18,126,240\}$ indexing the four
exceptional dimensions above,
\[
  \sum_{d\in\mathcal{D}} \pi(d) = 6+7+30+52 = 95,
\]
and with $\mathcal{H} = \{h(G_2), h(B_4), h(F_4), h(E_6), h(E_7),
h(E_8)\} = \{6,8,12,12,18,30\}$ the set of Coxeter numbers of the
exceptional algebras:
\[
  \sum_{h \in \mathcal{H}} \pi(h) = 3+4+5+5+7+10 = 34,
\]
and the prime-weighted Coxeter sum satisfies
$\sum_{h\in\mathcal{H}} \pi(h) + \pi(\dim F_4) = 34 + 52 = 86$,
bracketed between Mersenne primes $M_4 = 31$ and $M_5 = 127$.
\end{theorem}

\begin{table}[h]
\centering
\caption{Prime counts at exceptional Lie data}
\label{tab:pi-exceptional}
\begin{tabular}{lrrl}
\hline
Argument & Value & $\pi(\cdot)$ & Lie interpretation\\
\hline
$\dim(G_2)$   & 14  & 6  & $h(G_2)$\\
$h(E_6)+6$    & 18  & 7  & $\operatorname{fund}(G_2)$\\
$h(E_8)\cdot 4+6$ & 126 & 30 & $h(E_8)$\\
$|\Phi(E_8)|$ & 240 & 52 & $\dim(F_4)$\\
\hline
\end{tabular}
\end{table}

%--------------------------------------------------------------------
\subsection{The Boomerang Property}

\begin{proposition}[Boomerang through prime distribution]
\label{prop:boomerang}
Let $p_{\mathrm{prev}}(n)$ denote the largest prime strictly less
than $n$.  Then:
\begin{align*}
  p_{\pi(\dim E_6)} &= p_{36} = 151,\\
  p_{\pi(\dim E_6)-1} &= p_{\mathrm{prev}}(\dim E_6) = 73,\\
  p_{\pi(\dim E_8)} &= p_{52} = 239,\\
  p_{\pi(\dim E_8)+1} &= p_{\mathrm{prev}}(\dim E_8) + 2 = 241.
\end{align*}
In particular,
\[
  p_{\pi(\dim E_6)} - 2 = p_{\mathrm{prev}}(\dim E_6) = 73
  \qquad\text{and}\qquad
  p_{\pi(\dim E_8)+1} = p_{\mathrm{prev}}(\dim E_8) = 241.
\]
Both exceptional dimensions are fixed points of the composition
$n \mapsto p_{\mathrm{prev}}(n) \mapsto \pi(\,\cdot\,) \mapsto p_{(\,\cdot\,)}$
up to a shift of $\pm 1$ in the index, a \emph{boomerang property}.
\end{proposition}

%--------------------------------------------------------------------
\subsection{Prime Gaps and the Fine-Structure Constant}

\begin{observation}[Fine-structure constant from prime gaps]
\label{obs:prime-gaps}
Let $p_{\mathrm{next}}(n)$ denote the smallest prime strictly greater
than $n$. Then:
\[
  p_{\mathrm{next}}(\dim E_7) = p_{\mathrm{next}}(133) = 137
  = \lfloor 1/\alpha_{\mathrm{QED}} \rfloor,
\]
where $\alpha_{\mathrm{QED}}^{-1} \approx 137.036$ is the
fine-structure constant. Moreover, the three consecutive prime gaps
\[
  p_{\mathrm{next}}(\dim F_4) - \dim F_4
  = p_{\mathrm{next}}(\dim E_6) - \dim E_6
  = p_{\mathrm{next}}(\dim E_7) - \dim E_7
  = 4
\]
each equal $h(A_3)$, while the gap to $137$ satisfies
$137 - 133 = 4 = h(A_3) = h(G_2) - 2$.
\end{observation}

%--------------------------------------------------------------------
\subsection{Waring's Problem and Exceptional Lie Data}

\begin{proposition}[Waring function identities]
\label{prop:waring}
Let $g(k)$ denote the exact Waring number (least $s$ such that every
positive integer is a sum of $s$ perfect $k$th powers) and $G(k)$ the
asymptotic Waring number. The following identities hold:
\begin{align*}
  g(2) &= 4 = \operatorname{rank}(F_4),\\
  g(3) &= 9 = h^\vee(F_4),\\
  g(4) &= 19 = \text{exponent index of } E_8
             \text{ (19th prime)},\\
  g(6) &= 73 = p_{\mathrm{prev}}(\dim E_6),\\
  G(3) &= 7  = \operatorname{fund}(G_2).
\end{align*}
In particular, $g(2) + g(3) = 13 = h(G_2) + h^\vee(G_2)/h(G_2)$, and
\[
  g(2) \cdot g(3) = 36 = \dim G_2 + 22 = h(E_6)/\!\operatorname{rank}(G_2)\cdot 6.
\]
\end{proposition}

\begin{proof}
The values $g(2)=4$, $g(3)=9$, $g(4)=19$, $g(6)=73$, and $G(3)=7$
are classical results of Waring (1770) and Hardy--Littlewood (1920).
The Lie identifications use
$\operatorname{rank}(F_4)=4$, $h^\vee(F_4)=9$,
$\operatorname{fund}(G_2)=7$ (Section~\ref{sec:fixed-points}),
and $p_{\mathrm{prev}}(\dim E_6)=73$ (Proposition~\ref{prop:boomerang}).
\end{proof}

%============================================================================
% Section 33: Combinatorial Encoding at Exceptional Ranks
%============================================================================

\section{Combinatorial Encoding at Exceptional Ranks}%
\label{sec:combinatorial-exceptional}

Factorial and combinatorial arithmetic functions encode exceptional Lie
data with the same density as the multiplicative functions studied
throughout this paper.  We record the principal identities.

\subsection{Factorial Bridges}

\begin{theorem}[Factorial--Exceptional Dimension Identity]
\label{thm:factorial-bridges}
The following equalities hold:
\begin{align}
  \sigma(7!) &= \dim(\mathfrak{e}_6)\cdot\dim(\mathfrak{e}_8)
               = 78\cdot 248 = 19344, \label{eq:sigma-7fac}\\
  \varphi(4!) &= \operatorname{rank}(\mathfrak{e}_8) = 8, \label{eq:phi-4fac}\\
  \varphi(7!) &= |W(F_4)| = 1152. \label{eq:phi-7fac}
\end{align}
Moreover, the double-factorial $13!! = 1\cdot3\cdot5\cdots13 = 135135$
satisfies
\begin{equation}
  \varphi(13!!) = |W(E_6)| = 51840. \label{eq:phi-13dfac}
\end{equation}
\end{theorem}

\begin{proof}
Equation~\eqref{eq:sigma-7fac}: $7!=5040$ and
$\sigma(5040)=19344=78\cdot248$ by direct computation.
Equation~\eqref{eq:phi-4fac}: $\varphi(24)=8=\operatorname{rank}(E_8)$.
Equation~\eqref{eq:phi-7fac}: $\varphi(5040)=1152=|W(F_4)|$.
Equation~\eqref{eq:phi-13dfac}: $13!!=135135$ and
$\varphi(135135)=51840=|W(E_6)|$.
\end{proof}

\subsection{Central Binomial Coefficients at Exceptional Ranks}

\begin{proposition}[Binomial--Root-System Identities]
\label{prop:central-binomial}
Let $\binom{n}{k}$ denote the binomial coefficient.  Then
\begin{align}
  \varphi\!\left(\tbinom{12}{6}\right) &= |\Phi(E_8)| = 240,
    \label{eq:phi-C126}\\
  \sigma\!\left(\tbinom{14}{7}\right) &= 2\cdot 7! = 10080.
    \label{eq:sigma-C147}
\end{align}
\end{proposition}

\begin{proof}
$\binom{12}{6}=924$ and $\varphi(924)=240=|\Phi(E_8)|$.
$\binom{14}{7}=3432$ and $\sigma(3432)=10080=2\cdot5040=2\cdot7!$.
\end{proof}

\subsection{Primorial Chains}

\begin{observation}[Primorial Arithmetic]
\label{obs:primorial}
Let $n\#$ denote the primorial (product of primes $\leq n$).  Then
$6\#=30$ and $10\#=210$ satisfy
\begin{align}
  \operatorname{rank}(E_6)\# = 6\# &= 30 = h(E_8), \label{eq:rank-E6-primorial}\\
  \sigma(30) &= 72 = |\Phi(E_6)|, \label{eq:sigma-30}\\
  \sigma(210) &= 576 = 24^2, \label{eq:sigma-210}\\
  \varphi(30) &= 8 = \operatorname{rank}(E_8), \label{eq:phi-30}\\
  \varphi(210) &= 48 = |\Phi(F_4)|. \label{eq:phi-210}
\end{align}
\end{observation}

The chain $30 \to 72 \to 240$ (applying $\sigma$ twice starting from
$\operatorname{rank}(E_6)\#$) terminates at $|\Phi(E_8)|$, linking the
primorial of $E_6$'s rank to the root system of $E_8$ via two divisor-sum
steps.

%============================================================================
% Section 34: Extremal Combinatorics and Exceptional Invariants
%============================================================================

\section{Extremal Combinatorics and Exceptional Invariants}%
\label{sec:extremal-combinatorics}

Ramsey theory, Schur numbers, and van der Waerden numbers provide a
fourth arena---alongside number theory, geometry, and representation
theory---in which exceptional Lie algebra invariants appear as
distinguished values.

\subsection*{34.1 Ramsey Numbers}

The diagonal and off-diagonal Ramsey numbers $R(3,k)$ for
$k=3,\dots,9$ produce seven consecutive values whose exact coincidence
with Lie-algebraic invariants is striking.

\begin{theorem}[Ramsey--Lie coincidence]\label{thm:ramsey-lie}
The seven known values $R(3,k)$ for $k=3,\dots,9$ are
\[
  6,\;9,\;14,\;18,\;23,\;28,\;36,
\]
which equal, respectively,
$h(G_2)$, $h^{\vee}(F_4)$, $\dim(G_2)$, $h(E_7)\!\cdot\!\tfrac{1}{2}$,
$e_1(E_8)$, $\dim(D_4)$, and $|\Phi^+(E_6)|$,
where $h$ denotes the Coxeter number, $h^{\vee}$ the dual Coxeter
number, $e_1$ the smallest positive exponent, and $\Phi^+$ the
positive root system.
\end{theorem}

\begin{table}[h]
\centering
\caption{Ramsey numbers $R(3,k)$ and their Lie-algebraic interpretations.}
\label{tab:ramsey}
\begin{tabular}{cccl}
\hline
$k$ & $R(3,k)$ & Algebra & Invariant \\
\hline
3 & 6  & $G_2$ & $h(G_2)=6$ \\
4 & 9  & $F_4$ & $h^{\vee}(F_4)=9$ \\
5 & 14 & $G_2$ & $\dim(G_2)=14$ \\
6 & 18 & $E_7$ & $h(E_7)=18$ \\
7 & 23 & $E_8$ & $e_1(E_8)=23$ (smallest exponent +1) \\
8 & 28 & $D_4$ & $\dim(D_4)=28$ \\
9 & 36 & $E_6$ & $|\Phi^+(E_6)|=36$ \\
\hline
\end{tabular}
\end{table}

\subsection*{34.2 Schur Numbers}

\begin{proposition}[Schur--rank coincidence]\label{prop:schur}
The Schur numbers $S(n)$ for $n=2,3,4$ satisfy
\[
  S(2)=4=\operatorname{rank}(F_4),\qquad
  S(3)=13=h(F_4)-1,\qquad
  S(4)=44=4N_Q,
\]
where $h(F_4)=14$ is the Coxeter number of $F_4$ and
$N_Q=11$ is the order of the Qameah magic square.
\end{proposition}

\subsection*{34.3 Van der Waerden Numbers}

\begin{observation}[van der Waerden--$F_4$/$E_6$ coincidence]%
\label{obs:vdw}
The van der Waerden numbers
\[
  W(2;3)=9=h^{\vee}(F_4),\qquad W(3;3)=27=\dim(\text{fund.\ rep.\ of }E_6)
\]
recover the dual Coxeter number of $F_4$ and the dimension of the
minimal faithful representation of $E_6$.
\end{observation}

\subsection*{34.4 Arithmetical Universality}

\begin{remark}[Arithmetical universality of exceptional invariants]%
\label{rem:universality}
Sections~1--34 demonstrate that the invariants of the exceptional
simple Lie algebras $G_2$, $F_4$, $E_6$, $E_7$, $E_8$---Coxeter numbers,
exponents, root-system cardinalities, representation dimensions---appear
as distinguished values across at least six independent combinatorial
and number-theoretic domains: divisor sums and Euler's totient
(Sections~2--5), Bernoulli numbers and class numbers (Sections~6--10),
Diophantine and Waring problems (Sections~11--15), Ramanujan congruences
and prime gaps (Sections~16--20), partition functions and
Catalan--Stirling identities (Sections~21--28), and extremal graph
theory (this section).  This arithmetic universality may reflect a
deep structural reason connected to the ADE classification of simple
Lie algebras and the role of exceptional symmetries as distinguished
attractor points in the landscape of finite combinatorial structures.
\end{remark}

%============================================================================
% Bibliography
%============================================================================

\begin{thebibliography}{99}

\bibitem{QHOTS2026}
H.~Th.\ Reynisson (with the Sefirot Autonomous Research System),
\emph{The Geometric Completion: Unified $\sqrt{n}$ Framework for Fundamental Physics},
in preparation, 2026.

\bibitem{Adams1962}
J.\,F.~Adams,
Vector fields on spheres,
\emph{Ann.\ of Math.} (2) \textbf{75} (1962), 603--632.

\bibitem{Adams1996}
J.\,F.~Adams,
\emph{Lectures on Exceptional Lie Groups},
Chicago Lectures in Mathematics, Univ.\ Chicago Press, 1996.

\bibitem{Bourbaki2002}
N.~Bourbaki,
\emph{Lie Groups and Lie Algebras, Chapters 4--6},
Springer-Verlag, Berlin, 2002.

\bibitem{Freudenthal1964}
H.~Freudenthal,
Lie groups in the foundations of geometry,
\emph{Adv.\ Math.} \textbf{1} (1964), 145--190.

\bibitem{Gresnigt2023}
N.~Gresnigt, L.~Gourlay, and A.~Varma,
Three generations of colored fermions with $S_3$ family symmetry
from Cayley--Dickson sedenions,
\emph{Eur.\ Phys.\ J.\ C} \textbf{83} (2023), 747.

\bibitem{ConwaySloane1999}
J.\,H.~Conway and N.\,J.\,A.~Sloane,
\emph{Sphere Packings, Lattices and Groups},
3rd ed., Grundlehren der math.\ Wiss., vol.~290, Springer, 1999.

\bibitem{HardyWright2008}
G.\,H.~Hardy and E.\,M.~Wright,
\emph{An Introduction to the Theory of Numbers},
6th ed., Oxford Univ.\ Press, 2008.

\bibitem{Humphreys1972}
J.\,E.~Humphreys,
\emph{Introduction to Lie Algebras and Representation Theory},
Graduate Texts in Mathematics, vol.~9, Springer-Verlag, 1972.

\bibitem{McKay1980}
J.~McKay,
Graphs, singularities, and finite groups,
\emph{Proc.\ Sympos.\ Pure Math.} \textbf{37} (1980), 183--186.

\bibitem{Milnor1956}
J.\,W.~Milnor,
On manifolds homeomorphic to the $7$-sphere,
\emph{Ann.\ of Math.} (2) \textbf{64} (1956), 399--405.

\bibitem{Ramanujan1916}
S.~Ramanujan,
On certain arithmetical functions,
\emph{Trans.\ Cambridge Philos.\ Soc.} \textbf{22} (1916), 159--184.

\bibitem{Robin1984}
G.~Robin,
Grandes valeurs de la fonction somme des diviseurs et
hypoth\`ese de Riemann,
\emph{J.\ Math.\ Pures Appl.} (9) \textbf{63} (1984), 187--213.

\bibitem{Tits1966}
J.~Tits,
Alg\`ebres alternatives, alg\`ebres de Jordan et alg\`ebres de Lie
exceptionnelles.\ I.\ Construction,
\emph{Indag.\ Math.} \textbf{28} (1966), 223--237.

\bibitem{Sloane2024}
N.\,J.\,A.~Sloane (ed.),
\emph{The On-Line Encyclopedia of Integer Sequences},
\url{https://oeis.org}, 2024.

\bibitem{vonStaudt1840}
K.\,G.\,C.~von Staudt,
\emph{De numeris Bernoullianis},
Erlangen, 1840.

\bibitem{Kostant1959}
B.~Kostant,
The principal three-dimensional subgroup and the Betti numbers of a
complex simple Lie group,
\emph{Amer.\ J.\ Math.} \textbf{81} (1959), 973--1032.

\bibitem{Zuber1998}
J.-B.~Zuber,
Generalisations of the McKay correspondence,
in: \emph{Conformal Field Theories and Integrable Models} (Z.~Horv\'ath
and L.~Palla, eds.), Lecture Notes in Phys., vol.~498, Springer, 1998,
pp.~173--193.

\bibitem{Mathlib2024}
The mathlib Community,
\emph{Mathlib: the Lean 4 mathematical library},
\url{https://github.com/leanprover-community/mathlib4}, 2024.

\end{thebibliography}

\end{document}
